\section{Lift estimates II: assembly}\label{sec:pf:lift2}

\providecommand{\lifttwomainT}{\mathrm{mainT}}
\providecommand{\lifttwolcv}{\operatorname{lcv}}
\providecommand{\lifttwork}{\operatorname{rk}}
\providecommand{\lifttwovdm}{\operatorname{vdm}}
\providecommand{\lifttwohdeg}{\widehat{\operatorname{deg}}}
\providecommand{\lifttwoLat}[1]{\Lambda^{\mathrm{#1}}}
\providecommand{\lifttwoP}[1]{\mathsf{#1}}
\providecommand{\lifttwoC}[1]{\mathsf{#1}}
\providecommand{\lifttwodF}{d_{\mathrm{Fl}}}
\providecommand{\lifttwobin}{B_T}
\providecommand{\lifttwozed}{\zeta}
\providecommand{\lifttwoCov}{\operatorname{Cov}}
\providecommand{\lifttwoComp}{\operatorname{Comp}}
\providecommand{\lifttwoEv}{\operatorname{Ev}}
\providecommand{\lifttwoL}{\mathcal{L}}
\providecommand{\lifttwoCv}{\mathcal{V}}

This section proves the lift estimates of Part~I (\cref{thm:covolume:lift}): from seven packages of
geometric inputs, among them the mean-slope statement $\sfname{DeepStR}$ of \cref{def:covolume:deepstr},
the development derives the proposition $\sfname{CovolumeBoundOfX}$ of \cref{def:covolX}, the archimedean
``$L_3$ condition'' on the base sides of the points of the admissible families. The derivation is an
assembly of twenty-six \emph{conditions}, each of the form ``eventually, at all points $p$ of all
admissible families and for all indices $i$, a certain inequality or structural statement holds at
$(p,i)$'' (\cref{tab:covolume:conditions}). We state each condition precisely and prove it from the
inputs and from the conditions proved before it. Three of the conditions (the height lower bound, the
Hilbert--Samuel lower bound for the component degrees, and the lower bound at the top index) are proved in
the companion section on lift estimates; we cite them.

\paragraph{The inputs.} We use the following abbreviations for the hypotheses of the theorem (all of them are
among, or derived from, the inputs of \cref{tab:covolume:inputs}).
\begin{description}[style=unboxed,leftmargin=0pt,font=\normalfont]
\item[(VE)] The two volume estimates: an upper bound for the arithmetic self-intersection of a Hermitian line bundle
on a model of the image of a lattice of cusp forms, in terms of the degree and the height of the image; and a
comparison of the number of small elements of a lattice of cusp forms with the number of small sections of the
corresponding line bundles on such a model.
\item[(X)] Smooth toroidal models of unitary moduli spaces with logarithmic Kodaira--Spencer isomorphism
\textup{(KS)}; seminegativity of subbundles of $\Lie G$ for semi-abelian $G$ over curves \textup{(Lie)}; and Green's
theorem on cuspidal characters \textup{(Gr)}.
\item[(Hg)] The Hodge package of Part~I: twenty-eight published statements on toroidal compactifications,
arithmetic volumes, de Rham cohomology of abelian schemes, ray class fields, vanishing theorems, Chow forms and
heights, GAGA and others (\cref{tab:covolume:inputs}).
\item[(H)] Toroidal compactifications of Siegel moduli spaces \textup{(Sg)}, toroidal charts, and the
compatibilities of arithmetic degrees.
\item[(GM)] Grothendieck--Messing deformation theory over nilpotent divided-power thickenings.
\item[(Q)] Kawamata--Viehweg vanishing, $L^2$-extension for log canonical pairs, Serre's theorem on holomorphic
sections, numerics of nef divisors, and existence, integrality and hyperplane-section compatibility of Chow forms.
\item[(DS)] The mean-slope statement $\sfname{DeepStR}$ (\cref{def:covolume:deepstr}).
\end{description}

\paragraph{The idea.} By the final assembly of Part~I (\cref{thm:covolume:assembly}), the covolume bound
follows from eight \emph{lift conditions} at the points: a Gram condition, a comparison of the normalised
covolume of $M_i$ with the slope of the lattice $\lifttwoLat{good}$ of good lifted forms, a truncation by an
idempotent of the group algebra of $G_p$, a bound for the rank of the bad part, a bound for the maximal
slope, size bounds, a comparison of the lattice $\lifttwoLat{type}$ with the corresponding lattice of
automorphic forms of the \emph{core} datum, and a lower bound $\mu_{\mathrm{core}}\ge-\lifttwomainT-CU$ for the
slope $\mu_{\mathrm{core}}$ of the latter. All but the last are comparisons between the lattices at a point
and the core lattices; they reduce to statements on Hermitian lattices under extension of the base field
(\cref{sec:pf:lift2:lattices}) and to properties of the points. The lower bound for $\mu_{\mathrm{core}}$ is
the heart of the matter. For $i+1<m$ it is obtained by \emph{descent} from a \emph{deep cover} $Y$, a unitary
Shimura variety of higher level over a larger coefficient field: the slope of the lattice of cusp forms of
the anchored weight $u_M$ on $Y$ is bounded below by an arithmetic Hilbert--Samuel argument on each
irreducible component (this is where $\sfname{DeepStR}$ enters, through the height lower bound), the passage
from $u_M$ to the shifted weight $G_s(-c_0)$ costs at most $\lifttwomainT$ (twisting by a flat line bundle),
multiplication by a short integral form of weight $c_0$ times the Hodge weight maps weight $G_s(-c_0)$ to
weight $G_k$ with controlled norms, and the slope inequality of \cref{lem:covolume:descent} together with a
rank comparison transports the lower bound to the lattice of the given type, which is finally compared with
the core lattice by pull-back along the level covering.

\paragraph{Organisation.} \Cref{sec:pf:lift2:conv} fixes the quantifier conventions and the combination
principle by which conditions are combined; \cref{sec:pf:lift2:lattices} collects the statements on
Hermitian lattices used below; \cref{sec:pf:lift2:objects} defines the objects at a point and on deep covers
and the conditions; \cref{sec:pf:lift2:numeric} establishes the numerical properties of the
families; \cref{sec:pf:lift2:point,sec:pf:lift2:liftpred} prove the conditions at a point;
\cref{sec:pf:lift2:deep} treats deep covers and the uniform lower bound; \cref{sec:pf:lift2:descent} the
descent; \cref{sec:pf:lift2:bad} the rank of the bad part; \cref{sec:pf:lift2:main} assembles the proof
of \cref{thm:covolume:lift}; and \cref{sec:pf:lift2:aux} proves auxiliary results on weight parameters, degrees of images,
toroidal flag data and cusp forms on deep covers, which are used above and in the companion section. Many results of this section have proofs in the development that rest on
further, deeper geometric statements (toroidal compactifications, Kodaira--Spencer maps, log canonical pairs,
flat line bundles, theta series). These are stated where they are used, with references to their formal
versions; their proofs are not reproduced here.

% ---------------------------------------------------------------------------
\subsection{Conventions}\label{sec:pf:lift2:conv}

We use the notions of \cref{def:families:closed} and \cref{sec:covolume:statement}: families
$F\in\mathrm{FamP}(\eps)$, the class $\mathrm{FamHyp}$, the quantifier $\mathrm{Closed}$, points
$p\in\mathrm{Pt}(F,x,\iota)$ ($x>x_1(F)$, $\iota\in I_x$), the structural conditions $\mathrm{PointConditions}'$
and $\mathrm{RecTied}$, and the predicate $\mathrm{EventuallyAtAllPoints}(F,Q)$, which we abbreviate to
$\lifttwoEv(F,Q)$: there is $x_0$ such that $Q(x,\iota,p)$ holds for all $x>x_1(F)$ with $x\ge x_0$, all
$\iota\in I_x$ and all $p\in\mathrm{Pt}(F,x,\iota)$ satisfying
$\mathrm{PointConditions}'(p)\wedge\mathrm{RecTied}(D(x,\iota),L(x,\iota))$. We call such points
\emph{admissible}.

At a point $p$ at $(x,\iota)$ we write $D=D(x,\iota)$ for its global data, with ranks $k,m,f,t$,
$r=k+m-1$, $n_V=r+t$, CM host field $E$ with auxiliary primes $d$, $g_0$, $g_1$, and totally complex field $K$
(the field over which the base side of the point is defined); its size data $L=L(x,\iota)$ have scale $x$, exponent
$\eps\in(0,\frac1{52})$, a parameter $J\ge1$, and we put
\begin{equation}\label{eq:pf:lift2:scale}
e_z=\eps/J,\qquad \ell=x^{\eps},\qquad b=e^{\ell},\qquad z=f\,r\,\ell .
\end{equation}
The family provides, for each $(x,\iota)$, a real parameter $P=P(x,\iota)$ (denoted $PP$ in Part~I). For
$i\in J_m=\{0,\dots,m-1\}$ we write $U=\mathrm{unitE}(p)=f\,m\,b\,\ell\log\log x$ (\cref{def:covolume:mainT})
and $\lifttwomainT=\lifttwomainT_i(p)$; these are the error unit and the main term of
\cref{def:covolume:mainT}.

\begin{definition}[Lift conditions and cover conditions]\label{pf:lift2:def:lifts}
Let $Q(F,x,\iota,p,i,P,C)$ be a statement about a point $p$ at $(x,\iota)$ of a family $F$, an index $i\in J_m$
and two reals $P,C$. The \emph{lift condition} $\lifttwoL(Q)$ of $Q$ is the statement
\[
\mathrm{Closed}\Bigl(\mathrm{FamHyp},\,(F,C_H)\mapsto\exists C\ge0:\
\lifttwoEv\bigl(F,\,p\mapsto\forall i\in J_m:\ Q(F,x,\iota,p,i,P(x,\iota),C)\bigr)\Bigr).
\]
If $Q(F,x,\iota,p,i,Y,P,C)$ involves in addition a deep cover $Y$ of $p$ at level $i+1$
(\cref{pf:lift2:def:deep}), its \emph{cover condition} $\lifttwoCv(Q)$ is the lift condition of the
statement ``if $i+1<m$, then $Q(F,x,\iota,p,i,Y,P,C)$ for every good deep cover $Y$''.
\formallabel{GlobalTheory.CovTop.LiftsCondition}
\formallabel{GlobalTheory.CovTop.CoversCondition}
\end{definition}

The constant $C$ is allowed to depend on $\eps$, on the family $F$ and on $C_H$, but not on $x$, $\iota$, $p$ or
$i$. The parameter $P$ is a datum of the family; it is never chosen in the proofs. Most conditions do not
involve $C$, or involve it only through the error unit $U$.

\begin{lemma}[Combination principle]\label{pf:lift2:combine}
\begin{enumerate}[label=(\alph*)]
\item If $\lifttwoEv(F,Q)$ and $\lifttwoEv(F,Q')$, then $\lifttwoEv(F,Q\wedge Q')$.
\item $\mathrm{Closed}(\mathcal C,P)$ and $\mathrm{Closed}(\mathcal C,P')$ imply $\mathrm{Closed}(\mathcal C,P\wedge P')$;
and if $\mathrm{Closed}(\mathcal C,P)$ and $P(F,C_H)$ implies $P'(F,C_H)$ for every $\eps>0$, every
$F\in\mathrm{FamP}(\eps)$ and every $C_H$, then $\mathrm{Closed}(\mathcal C,P')$.
\item Let $Q_1,\dots,Q_n,Q$ be statements as in \cref{pf:lift2:def:lifts} and assume $\lifttwoL(Q_1),\dots,
\lifttwoL(Q_n)$. Suppose that for every $\eps>0$, every $F\in\mathrm{FamP}(\eps)\cap\mathrm{FamHyp}$ and all
$C_1,\dots,C_n\ge0$ there are $C\ge0$ and $x_0'$ such that, for all admissible points $p$ at $(x,\iota)$ with
$x\ge x_0'$ and all $i\in J_m$, $Q_1(\dots,C_1)\wedge\dots\wedge Q_n(\dots,C_n)$ implies $Q(\dots,C)$. Then
$\lifttwoL(Q)$ holds.
\end{enumerate}
\formallabel{GlobalTheory.CovTop.EventuallyAtAllPoints.and}
\formallabel{Formalism.MainNotions.Closed.and}
\formallabel{Formalism.MainNotions.Closed.imp}
\end{lemma}

\begin{proof}
(a) Take the larger of the two thresholds $x_0$. (b) For the conjunction take the smaller of the two
$\eps_1$ and, for given $\eps$, the supremum of the two transported constants; a family whose constants are
large with respect to the supremum is large with respect to both. The second claim is immediate from the
definition. (c) Combine the $n$ closed statements by (b), fix $\eps$, $F$ and $C_H$ and
take the constants $C_1,\dots,C_n$ and thresholds they provide; combine the thresholds by (a), and take the
maximum with $x_0'$.
\end{proof}

All derivations of conditions from other conditions below are applications of
\cref{pf:lift2:combine}(c); we only state and prove the pointwise implication and the choice of $C$ and
$x_0'$. A condition $Q$ which holds at \emph{all} admissible points, for some constant $C$ depending only on
$F$, satisfies $\lifttwoL(Q)$ with $\eps_1=1$, an arbitrary transported constant and $x_0=0$.

% ---------------------------------------------------------------------------
\subsection{Hermitian lattices}\label{sec:pf:lift2:lattices}

We use the conventions of \cref{def:covolume:slopes}: $K$ is a number field, $V$ a finite-dimensional
$K$-vector space, and a family $\nu=(\nu_\sigma)_{\sigma\colon K\to\C}$ of functions $V\to\R$ \emph{has a Gram
frame} if there are a basis $(b_a)$ of $V$ and invertible complex matrices $X(v)$, one for each infinite place
$v$, such that $\nu_{\sigma_v}(x)=|X(v)^{-1}(\sigma_v(x_a))_a|$ for every $v$ and $x=\sum_ax_ab_a$, where
$\sigma_v$ is the embedding chosen for $v$. For a finitely generated $\OO_K$-submodule $L\subseteq V$ we have
the rank $\lifttwork L=\dim_KKL$, the slope $\mu_\nu(L)$ and the maximal slope $\mu_{\max,\nu}(L)$ of
\cref{def:covolume:slopes} (both $0$ if $\nu$ has no Gram frame), and we put
\begin{equation}\label{eq:pf:lift2:deg}
\lifttwohdeg_\nu(L)=\tfrac12\,\lifttwork(L)\,\mu_\nu(L).
\end{equation}
By \cref{rem:covolume:arakelov}, $\lifttwohdeg_\nu(L)$ is the Arakelov degree of $L$ divided by $[K:\Q]$
(this is half of the quantity denoted $\deg L$ in \cref{def:covolume:slopes}). Let $K$ be totally complex,
with $r_2=[K:\Q]/2$ complex places. If $L$ has rank $n$ and $e_1,\dots,e_n\in L$ is a $K$-basis of $KL$, then
\begin{equation}\label{eq:pf:lift2:gram}
r_2\,n\,\mu_\nu(L)=\log[L:\textstyle\sum_a\OO_Ke_a]-\sum_v\log\det G_v(e),
\end{equation}
where $G_v(e)=(\langle e_a,e_b\rangle_v)_{a,b}$ is the Gram matrix of $e$ for the Hermitian form at $v$ given by
the frame; this is a direct computation with the definition of $\mu_\nu$, using that the realisation of
$\OO_Ke$ has covolume $(2^{-r_2}|d_K|^{1/2})^n\prod_v\det G_v(e)$.

\begin{lemma}[Gram frames from inner products]\label{pf:lift2:hasgram}
Let $\nu=(\nu_\sigma)$ be a family of functions on a finite-dimensional $K$-vector space $V$. Suppose that for
every infinite place $v$ there are a complex inner product space $H_v$ and an additive map $\Phi_v\colon V\to
H_v$ which is $\sigma_v$-semilinear or $\bar\sigma_v$-semilinear, maps $K$-linearly independent families to
$\C$-linearly independent families, and satisfies $\nu_{\sigma_v}(x)=\|\Phi_v(x)\|$. Then $\nu$ has a Gram
frame.
\formallabel{GlobalTheory.HSApp.hasGram_of_inner}
\end{lemma}

\begin{proof}
Fix a basis $(b_a)_{a\le n}$ of $V$. For each $v$ the vectors $\Phi_v(b_a)$ are linearly independent, so
there is an invertible matrix $B_v$ with $\|\sum_ac_a\Phi_v(b_a)\|=|B_vc|$ for all $c\in\C^n$ (orthonormalise).
If $\Phi_v$ is $\sigma_v$-semilinear, $\|\Phi_v(x)\|=|B_v(\sigma_v(x_a))_a|$ and we take $X(v)=B_v^{-1}$; if it is
$\bar\sigma_v$-semilinear, $\|\Phi_v(x)\|=|\overline{B_v}(\sigma_v(x_a))_a|$ (complex conjugation preserves the
Euclidean norm) and we take $X(v)=\overline{B_v}^{-1}$.
\end{proof}

The next two lemmas compare slopes under an extension $K\subseteq K'$ of totally complex fields. Let $V'$ be a
$K'$-vector space and $\iota\colon V\to V'$ an additive map with $\iota(ax)=a\,\iota(x)$ for $a\in K$, which maps
$K$-linearly independent families to $K'$-linearly independent families; for $L\subseteq V$ put
$\iota_*L=\sum_{x\in L}\OO_{K'}\iota(x)$.

\begin{lemma}[Slopes under extension of scalars]\label{pf:lift2:slopeext}
Let $\nu,\nu'$ be families with Gram frames on $V$, $V'$, with $\nu_{\bar\sigma}=\nu_\sigma$ for all $\sigma$, let
$c\in\R$ with $\nu'_{\sigma'}(\iota x)\le e^c\nu_{\sigma'|_K}(x)$ for all $\sigma'\colon K'\to\C$ and $x\in V$,
and let $L\subseteq V$ be a nonzero finitely generated $\OO_K$-submodule. Then
\[
\mu_\nu(L)-2c\le\mu_{\nu'}(\iota_*L).
\]
\formallabel{GlobalTheory.SlopeExt.slope_sq_sub_le_slope_sq_span}
\end{lemma}

\begin{proof}
Let $n=\lifttwork L>0$, choose a $K$-basis $e_1,\dots,e_n\in L$ of $KL$ and put $L'=\iota_*L$, $d=[K':K]$. The
vectors $\iota(e_a)$ are $K'$-linearly independent and span $K'L'$, so $\lifttwork L'=n$; $L'$ is finitely
generated. Since $\OO_{K'}$ is flat over $\OO_K$, $L'/\sum_a\OO_{K'}\iota(e_a)\cong\OO_{K'}\otimes_{\OO_K}
(L/\sum_a\OO_Ke_a)$, and a finite $\OO_K$-module $M$ satisfies $|\OO_{K'}\otimes_{\OO_K}M|=|M|^d$; hence
$[L':\sum_a\OO_{K'}\iota(e_a)]=[L:\sum_a\OO_Ke_a]^d$. Let $w$ be a place of $K'$ above the place $v$ of $K$. The
hypothesis on the norms says that the Hermitian norm of $\sum_ac_a\iota(e_a)$ at $w$ is at most $e^c$ times
that of $\sum_ac_ae_a$ at $v$ for all $c\in K^n$ (using $\nu_{\bar\sigma}=\nu_\sigma$ if the embeddings of $w$
and $v$ are conjugate), hence, by density of $K$ in $\C$, for all $c\in\C^n$; thus $\det G'_w(\iota e)\le
e^{2nc}\det G_v(e)$. Every place $v$ of $K$ has exactly $d$ places above it, all complex, so $r_2'=d\,r_2$.
By \eqref{eq:pf:lift2:gram} for $L'$ and $L$,
\[
r_2'\,n\,\mu_{\nu'}(L')=d\log[L:\textstyle\sum\OO_Ke_a]-\sum_w\log\det G'_w(\iota e)
\ge d\bigl(r_2n\mu_\nu(L)\bigr)-2ncr_2',
\]
and dividing by $r_2'n$ gives the claim.
\end{proof}

\begin{lemma}[Comparison through the different]\label{pf:lift2:slopediff}
Let $\nu,\nu'$ be families with Gram frames on $V$, $V'$, with $\nu_{\bar\sigma}=\nu_\sigma$ and
$\nu'_{\sigma'}(\iota x)=\nu_{\sigma'|_K}(x)$ for all $\sigma'$, $x$. Let $L\subseteq V$ be an
$\OO_K$-submodule and $M\subseteq V'$ a finitely generated $\OO_{K'}$-submodule with $\iota(L)\subseteq M$ and
$\delta M\subseteq\iota_*L$ for every $\delta$ in the different ideal $\mathfrak D_{K'}$ of $K'/\Q$. Then
\[
\mu_{\nu'}(M)-\frac{2\log|d_{K'}|}{[K':\Q]}\le\mu_\nu(L).
\]
\formallabel{GlobalTheory.SlopeDifferent.slope_sq_sub_le_of_different_smul}
\end{lemma}

\begin{proof}
If $M=0$ then $L=0$ (as $\iota$ is injective), and by definition the slope of the zero module is
$-\log2+\log|d_K|/[K:\Q]$, resp.\ $-\log2+\log|d_{K'}|/[K':\Q]$; the inequality reduces to
$-\log|d_{K'}|/[K':\Q]\le\log|d_K|/[K:\Q]$, which holds as both discriminants are nonzero integers. Let
$M\ne0$. Since
$\mathfrak D_{K'}\ne0$, $K'M\subseteq K'\iota(KL)$; so $L$ is finitely generated (it embeds into $M$), and a
$K$-basis $e_1,\dots,e_n\in L$ of $KL$ gives a $K'$-basis $\iota(e_a)$ of $K'M$; thus $\lifttwork
M=\lifttwork L=n>0$. With $L'=\iota_*L\subseteq M$ we have $\mathfrak D_{K'}M\subseteq L'$, hence
$[M:L']\le[M:\mathfrak D_{K'}M]=N(\mathfrak D_{K'})^n=|d_{K'}|^n$. The Gram matrices satisfy $G'_w(\iota
e)=G_v(e)$ for $w$ above $v$ (equality of norms), and, as in the previous proof,
$[L':\sum\OO_{K'}\iota(e_a)]=[L:\sum\OO_Ke_a]^d$. By \eqref{eq:pf:lift2:gram},
\[
r_2'n\mu_{\nu'}(M)=\log[M:L']+d\log[L:\textstyle\sum\OO_Ke_a]-d\sum_v\log\det G_v(e)\le n\log|d_{K'}|+d\,r_2n\mu_\nu(L),
\]
and dividing by $r_2'n=d\,r_2n=n[K':\Q]/2$ gives the claim.
\end{proof}

For the maximal slope the development proves the analogous comparison:

\begin{lemma}[Maximal slopes under extension of scalars]\label{pf:lift2:slopemaxext}
In the situation of \cref{pf:lift2:slopeext}, assume moreover that for every $\sigma'$ there is $c_{\sigma'}>0$
with $e^{-c}\le c_{\sigma'}\le e^{c}$ and $\nu'_{\sigma'}(\iota x)=c_{\sigma'}\nu_{\sigma'|_K}(x)$ for all $x$. Then
$\mu_{\max,\nu'}(\iota_*L)\le\mu_{\max,\nu}(L)+2c$.
\end{lemma}

This is \lref{GlobalTheory.T.maxSlope_sq_fieldExtension_le_comparable_norms} (deferred). Its
proof reduces, by base change to a Galois closure, to the statement that the maximal slope of the extended
lattice $\OO_{K''}\otimes L$ of a Galois extension $K''/K$ equals that of $L$ (the maximal destabilising
submodule is Galois-stable), \lref{GlobalTheory.T.maxSlope_extendedLattice_galois_le}.

We shall also use the following results of Part~I, whose formal versions belong to this part of the
development.

\begin{lemma}\label{pf:lift2:partI}
\begin{enumerate}[label=(\alph*)]
\item (\cref{lem:covolume:lcv-slope}) If $\nu$ has a Gram frame, $b_1,\dots,b_d\in V$ are linearly independent,
$L\subseteq\sum_kKb_k$, $X$ is a frame on $K^d$ and $0<s\le1$ with $\|c\|_{X,v}=s\,\nu_{\sigma_v}(\sum c_kb_k)$
for all $v,c$, and $M=\{c\in K^d:\sum c_kb_k\in L\}$ is a lattice, then $\lifttwolcv(X,M)/d\le-\mu_\nu(L)$.
\item (\cref{lem:covolume:goodbad}) The good/bad cut, as stated there.
\item (\cref{lem:covolume:slope-good}) The slope of the good part, as stated there.
\item (\cref{lem:covolume:descent}) The descent inequality, as stated there.
\end{enumerate}
\formallabel{GlobalTheory.CovolCore.lcv_div_le_neg_slope_sq}
\formallabel{Formalism.ExtensionFiltration.goodBadCut}
\formallabel{GlobalTheory.CutNum.slope_good_ge}
\formallabel{GlobalTheory.DescentFamily.descent}
\end{lemma}

Part~I contains complete proofs of (a)--(d). The formal proof of (c) uses (b) and the fact that a Pythagorean
idempotent is self-adjoint for the Hermitian forms of a Gram frame; that of (d) uses \cref{lem:covolume:slopes}(a),(b)
and the equality of the slopes of $B$ and $B^n$.

% ---------------------------------------------------------------------------
\subsection{Points, deep covers and the conditions}\label{sec:pf:lift2:objects}

Fix an admissible point $p$ at $(x,\iota)$ and $i\in J_m$, and put $j=i+1$.

\begin{definition}[Lattices at a point]\label{pf:lift2:def:lat}
We use the objects of \cref{sec:covolume:assembly}. Let $W_i=K^{n_W(i)}$ be the $i$-th space of lifted forms of
the real datum of $p$, with its integral structure $M_K(W_i)$ and its archimedean realisations, and let
$\nu=(\nu_\sigma)_{\sigma\colon K\to\C}$, where $\nu_\sigma(f)$ is the norm of the realisation of $f$ at the place of
$\sigma$. Put
\[
\lifttwoLat{all}=\textstyle\sum_{y\in M_K(W_i)}\OO_Ky,\qquad
\lifttwoLat{good}=\lifttwoLat{type}\cap\operatorname{im}e_i,\qquad \lifttwoLat{bad}=\lifttwoLat{type}\cap\ker e_i,
\]
where $\lifttwoLat{type}\subseteq\lifttwoLat{all}$ is the lattice of integral forms of the type $(P_0,\psi_0^{-1})$
prescribed by the face of $p$ (a subgroup $P_0$ of the level group and a character $\psi_0$ of it), contained in the
$K$-subspace $W_i^{\mathrm{type}}$ of vectors of that type, and $e_i$ is the operator on $W_i$ given by the
idempotent of the group algebra of the finite group $G_p$ attached to the family. The \emph{core datum} $T$ of $p$
provides, at level $j$, a unitary moduli datum of signature $(p_j,q_j)$ over a coefficient field $F_T$, and for every
weight $G$ the $F_T$-space $V_T(G)$ of cusp forms of weight $G$, its lattice $\Lambda_T(G)$ of integral forms and the
Petersson-type norms $\nu^T_G$; we write $\mathcal G_T(G)$ if $\nu^T_G$ has a Gram frame and $\Lambda_T(G)$ is finitely
generated and spans $V_T(G)$. The weight $\kappa_i$ of $p$ determines a weight $G_k$ (\cref{pf:lift2:def:wt}),
and the point provides an embedding $e_i$ of the core cusp forms of weight $G_k$ into $W_i$ (the tie between the
core and the real datum). Finally
\[
\Lambda^{\mathrm{type}}_T=\Lambda_T(G_k)\cap F_T\text{-span}\{f:\ e_i(f)\in W_i^{\mathrm{type}}\},\qquad
\mu_{\mathrm{core}}=\mu_{\nu^T_{G_k}}\bigl(\Lambda^{\mathrm{type}}_T\bigr).
\]
\formallabel{GlobalTheory.CovTop.nrmR}
\formallabel{GlobalTheory.CovTop.latAll}
\formallabel{GlobalTheory.CovTop.latType}
\formallabel{GlobalTheory.CovTop.latGood}
\formallabel{GlobalTheory.CovTop.latBad}
\formal{GlobalTheory.CovTop.eGood}
\formal{GlobalTheory.CovTop.latTypeK}
\formal{GlobalTheory.CovTop.muCore}
\end{definition}

\begin{definition}[Weights]\label{pf:lift2:def:wt}
Let $\Psi$ be the CM type of the host datum ($|\Psi|=f$ at the points considered) and $N=p_j+q_j$. A \emph{weight}
of the datum at level $j$ is a function $u\colon\Psi\times\{0,\dots,N-1\}\to\Z$. The weights of automorphic line
bundles used by the development form a group $\mathfrak W(E)$, and $u\mapsto\underline u$ is the homomorphism from
weights of the datum to $\mathfrak W(E)$. Let $h$ be the Hodge weight: $h_{\tau,a}=0$ for $a<p_j$ and $h_{\tau,a}=-1$
for $a\ge p_j$. Put
\[
\lifttwovdm(u)=\prod_{\tau\in\Psi}\prod_{0\le a<l<N}(u_{\tau,a}-u_{\tau,l}),\qquad
\rho_B(u)=1-\frac{\lifttwovdm(u-3Nh)}{\lifttwovdm(u)},\qquad \lifttwodF=|\Psi|\binom N2
\]
(the Weyl numerator, the relative loss of a shift by $3N$ times the Hodge weight, and the dimension of the flag
variety). The weight $\kappa_i$ consists, for each $\tau\in\Psi$, of two finite sequences $\kappa^{\pm}_\tau$ of
non-negative integers, and of an even integer $n_i$; it determines a weight $G_k\in\mathfrak W(E)$ and the weight
$\lambda$ of the datum with $\lambda_{\tau,a}=\kappa^+_\tau(a)+n_i/2$ for $a<p_j$ and
$\lambda_{\tau,a}=-\kappa^-_\tau(q_j-1-(a-p_j))-n_i/2$ for $a\ge p_j$ (sequences extended by $0$). Let
\begin{gather*}
\kappa=\textstyle\sum_{\tau,a}|\lambda_{\tau,a}|,\qquad a_\tau=\lambda_{\tau,k-2},\qquad u_W=\lambda-a,\\
u_M=u_W-c_0h,\qquad G_s(s)=G_k+\underline{s\,h}\quad(s\in\Z),
\end{gather*}
where $a$ is regarded as a weight constant in the second variable (the \emph{anchor}, so that $u_W$ vanishes at the
index $k-2$) and $c_0=c_0(D,j)$ is the least multiple of $8$ which is $\ge r+\binom{fp_jq_j+1}2+2$ and $\ge8fr+4$;
one has $0\le c_0\le22f^2r^4$. Finally let
\[
T=\bigl\lceil P^2\lifttwodF^2k/m\bigr\rceil,\qquad \lifttwobin=\binom{T-1}{\lifttwodF},\qquad
\lifttwozed(u_1)=7\Bigl(\frac{\lifttwodF^2}{T}+\lifttwodF\,\rho_B(u_1)\Bigr),
\]
and call a weight $u_1$ a \emph{quotient weight}, $u_1\in\mathcal Q$, if for every $\tau$, $u_{1;\tau,k-2}=0$ and, for $0\le a<N-1$,
\[
u_{1;\tau,a}-u_{1;\tau,a+1}=\Bigl\lfloor\frac{u_{W;\tau,a}-u_{W;\tau,a+1}+\epsilon_a}{T+\lifttwodF+1}\Bigr\rfloor,\qquad
\epsilon_a=\begin{cases}-c_0-2N+2,&a+1=p_j,\\ 2,&\text{otherwise}.\end{cases}
\]
\formallabel{GlobalTheory.CovTop.kap}
\formallabel{GlobalTheory.CovTop.uW}
\formallabel{GlobalTheory.CovTop.lenT}
\formallabel{GlobalTheory.CovTop.binT}
\formallabel{GlobalTheory.CovTop.zed}
\formallabel{GlobalTheory.CovTop.IsQuot}
\formallabel{GlobalTheory.CovTop.Quo}
\formal{GlobalTheory.AllWt.wtOf}
\formal{GlobalTheory.CompHt.hodgeWt}
\formal{GlobalTheory.CovTop.laOf}
\formal{GlobalTheory.CovTop.uM}
\formal{GlobalTheory.CovTop.cZero}
\end{definition}

\begin{definition}[Deep covers]\label{pf:lift2:def:deep}
A \emph{deep cover} of $p$ at level $j$ consists of a coefficient datum whose field contains $F_T$, a
change-of-level datum for the unitary moduli space of $T$ at level $j$, giving a moduli space $X_Y$ over a
coefficient field $F_Y$ (a totally complex number field) with a level of modulus $M'_Y$, complex uniformisations at
the embeddings $\sigma\colon F_Y\to\C$, dictionaries between their components and frames of the local unitary group,
and a compact open subgroup $K_f$. The set of deep covers is $\lifttwoCov(p,i)$. For $e\in\N$ and $o\in\R$ the
development defines when a deep cover is \emph{$(e,o)$-good}, by fifteen conditions: the prime-to-level part of
the level is $15^e$; the level is admissible with trivial abelian part; the uniformisations are unitary and of test
type; the characters of the arithmetic groups are trivial; the dictionaries are unitary of level $K_f$; the base
points have absolute value $\xi_R(D)$; the function fields are algebraically closed in the appropriate sense; the
coefficient field is normal over $F_T$; $\mathrm{fs}(F_Y)\le o$, where
$\mathrm{fs}(F)=\log[F:\Q]+\log|d_F|/[F:\Q]$; and three normalisations of frames and models. A deep cover is
\emph{good} if it is $(\lfloor3\ell/\log3\rfloor,391z^2)$-good. (The development adds two further clauses,
normalising the frames and the models at the good embeddings; they are consequences of the first.)

For a deep cover $Y$ and a weight $G\in\mathfrak W(E)$ we have, as for the core datum, the $F_Y$-space $V_Y(G)$ of
cusp forms of weight $G$ on $X_Y$, the lattice $\Lambda_Y(G)$ of integral forms, the Petersson-type norms $\nu^Y_G$,
the slope $\mu_Y(G)=\mu_{\nu^Y_G}(\Lambda_Y(G))$ and the rank $\lifttwork_Y(G)=\lifttwork\Lambda_Y(G)$; $G$ is
\emph{good on $Y$}, written $\mathcal G_Y(G)$, if $\nu^Y_G$ has a Gram frame and $\Lambda_Y(G)$ is finitely generated
and spans $V_Y(G)$. For a weight $u$ of the datum we write $V_Y(u)=V_Y(\underline u)$, etc. The lattice of the given
type on $Y$ is the intersection $\Lambda^{\mathrm{deep}}_Y$ of $\Lambda_Y(G_k)$ with the $F_Y$-span of the pull-backs
to $X_Y$ of the core forms $f\in V_T(G_k)$ with $e_i(f)\in W_i^{\mathrm{type}}$.

Let $\lifttwoComp(Y)$ be the set of irreducible components of the generic fibre of $X_Y$ over $F_Y$. For a component
$Z$ and a weight $u$ of the datum, the cusp forms supported on $Z$ form a subspace $V_Z(u)$ with lattice
$\Lambda_Z(u)$ and norms $\nu^Z_u$, and $\mathcal G_Z(u)$ means that $\nu^Z_u$ has a Gram frame and $\Lambda_Z(u)$ is
finitely generated and spans $V_Z(u)$. Put $\mathrm{rk}_Z(u)=\lifttwork\Lambda_Z(u)$ and
$\lifttwohdeg_Z(u)=\lifttwohdeg_{\nu^Z_u}(\Lambda_Z(u))$; let $\mathrm{dg}_Z(u)$ be the degree of the image of $Z$
under the cusp forms in $V_Z(u)$, i.e.\ the degree of an integral Chow form (in dimension $\lifttwodF$) of the ideal
of algebraic relations among these forms, and $\mathrm{ht}_Z(u)$ the normalised height of that Chow form with respect
to the Hermitian lattice $(\Lambda_Z(u),\nu^Z_u)$; both are $0$ if no such Chow form exists.
\formal{GlobalTheory.CovTop.Deep}
\formal{GlobalTheory.CovTop.Deep.Good}
\formal{GlobalTheory.CovTop.GoodCov}
\formallabel{GlobalTheory.CovTop.goodCov_iff}
\formal{GlobalTheory.AllWt.GoodW}
\formal{GlobalTheory.CovTop.latDeep}
\formal{GlobalTheory.CompHt.compDeg}
\formal{GlobalTheory.CompHt.compHeight}
\end{definition}

\begin{definition}[Predicates at a point]\label{pf:lift2:def:pred}
Let $\delta_x=\exp(-x^{e_z/7})$ and $A,C\in\R$. The following statements about $(p,i)$ are defined:
\begin{description}[style=unboxed,leftmargin=0pt]
\item[$\lifttwoP{GramAt}$] $\nu$ has a Gram frame, $\lifttwoLat{all}$ is finitely generated,
$\lifttwoLat{good}\ne0$, and $\mathcal G_T(G_k)$ holds for the core datum at level $j$.
\formallabel{GlobalTheory.CovTop.GramAt}
\item[$\lifttwoP{RecAt}$] $\lifttwolcv(X_{W,i},M_i)/D\le-\mu_\nu(\lifttwoLat{good})$, where $D$, $X_{W,i}$, $M_i$ are
the data of the base side $S=p.S(\mathrm{false})$ of $p$ (\cref{def:covolX}).
\formallabel{GlobalTheory.CovTop.RecAt}
\item[$\lifttwoP{CutAt}$] $e_i^2=e_i$; $\nu_\sigma(e_iy+w-e_iw)^2=\nu_\sigma(e_iy)^2+\nu_\sigma(w-e_iw)^2$ for all
$\sigma,y,w$; $|G_p|\,e_i(\lifttwoLat{type})\subseteq\lifttwoLat{type}$; and $2\log|G_p|\le b$.
\formallabel{GlobalTheory.CovTop.CutAt}
\item[$\lifttwoP{BadAt}$] $\lifttwork\lifttwoLat{bad}\le3\,\delta_x\lifttwork\lifttwoLat{type}$.
\formallabel{GlobalTheory.CovTop.BadAt}
\item[$\lifttwoP{LamAt}(A)$] $\mu_{\max,\nu}(\lifttwoLat{all})\le x^Ab$.
\formallabel{GlobalTheory.CovTop.LamAt}
\item[$\lifttwoP{SizeAt}(P,A)$] $1\le b\le U\le x^Ab$, $|\lifttwomainT|\le x^Ab$, $|\lifttwomainT|\le P\kappa$,
$b\le P\kappa$, and, if $j<m$: a quotient weight exists, and every quotient weight $u_1$ satisfies
$3\lifttwodF+5\le T$, $0\le\rho_B(u_1)<1$ and $\lifttwozed(u_1)<1$.
\formallabel{GlobalTheory.CovTop.SizeAt}
\item[$\lifttwoP{MeasAt}(C)$] for every $\sigma\colon K\to\C$ there is $c>0$ with $|\log c|\le CU$ and
$\nu_\sigma(e_i(f))=c\,\nu^T_{G_k,\sigma|F_T}(f)$ for all $f\in V_T(G_k)$.
\formallabel{GlobalTheory.CovTop.MeasAt}
\item[$\lifttwoP{RatAt}$] $\lifttwoLat{type}$ is the $\OO_K$-span of $M_K(W_i)\cap W_i^{\mathrm{type}}$.
\formallabel{GlobalTheory.CovTop.RatAt}
\item[$\lifttwoP{CoreAt}(C)$] $\mu_{\mathrm{core}}-CU\le\mu_\nu(\lifttwoLat{type})$.
\formallabel{GlobalTheory.CovTop.CoreAt}
\item[$\lifttwoP{LowAt}(C)$] $-\lifttwomainT-CU\le\mu_{\mathrm{core}}$.
\formallabel{GlobalTheory.CovTop.LowAt}
\item[$\lifttwoP{AnchorAt}(C)$] $\bigl|\lifttwomainT-\sum_{\tau\in\Psi}(p_j-q_j)\,a_\tau\log\Omega(p)\bigr|\le CU$,
$\bigl(\sum_\tau|a_\tau|\bigr)q_j\ell\le CU$, and $f\,r\,n_V\,\ell\le b$.
\formallabel{GlobalTheory.CovTop.AnchorAt}
\end{description}
\end{definition}

\begin{definition}[Predicates at a deep cover]\label{pf:lift2:def:predY}
Let $Y\in\lifttwoCov(p,i)$, $C\in\R$ and $u_1$ a weight. The following statements are defined
($\Sigma_Z$ denotes a sum over $\lifttwoComp(Y)$):
\begin{description}[style=unboxed,leftmargin=0pt]
\item[$\lifttwoP{DeepGoodAt}$] $\mathcal G_Y(G)$ for $G\in\{G_k,G_s(-c_0),G_s(c_0),\underline{u_M}\}$;
$\lifttwork_Y(G_s(-c_0))>0$; $\lifttwoComp(Y)$ is finite and nonempty; $\mathcal G_Z(u_M)$ for all $Z$.
\formallabel{GlobalTheory.CovTop.DeepGoodAt}
\item[$\lifttwoP{FlatDegAt}$] for every component $Z$ and every $\tau\in\Psi$ the automorphic line bundle on $Z$ of
the weight $-\mathbf 1_\tau$ (equal to $-1$ at $\tau$ and $0$ elsewhere) has a constant (flat) section,
and every such section has degree $\frac{p_j-q_j}2\log\Omega(p)+q_j(\log d-\log\mathrm{rd}_E)$.
\formallabel{GlobalTheory.CovTop.FlatDegAt}
\formallabel{GlobalTheory.FlatLine.FlatSec}
\item[$\lifttwoP{DegAt}(u_1)$] there is $\delta>0$ with $\delta\le\mathrm{dg}_Z(u_1)\le\delta/(1-\rho_B(u_1))$ for
all $Z$.
\formallabel{GlobalTheory.CovTop.DegAt}
\item[$\lifttwoP{ShareAt}(u_1)$] $\lifttwobin\,\mathrm{dg}_Z(u_1)\le\mathrm{rk}_Z(u_M)$ and
$(1-\lifttwozed(u_1))\,\mathrm{rk}_Z(u_M)\le\lifttwobin\,\mathrm{dg}_Z(u_1)$ for all $Z$.
\formallabel{GlobalTheory.CovTop.ShareAt}
\item[$\lifttwoP{HeightAt}(C,u_1)$] there are $c_s,c_c,X\ge0$ and $\mathrm{sl}\colon\lifttwoComp(Y)\to\R$ with, for
all $Z$, $0\le\mathrm{ht}_Z(u_1)+c_s$ and $(1-\rho_B(u_1))\max(0,\mathrm{sl}(Z)+c_s)\le\mathrm{ht}_Z(u_1)+c_s+c_c$;
$-X\le\Sigma_Z\mathrm{sl}(Z)/|\lifttwoComp(Y)|$; and $TX\le CU$, $Tc_c\le CU$,
$c_s\bigl(6\lifttwodF+10+T(\lifttwozed(u_1)+2\rho_B(u_1))\bigr)\le CU$.
\formallabel{GlobalTheory.CovTop.HeightAt}
\item[$\lifttwoP{HSAt}(C,u_1)$] there are $A_0>0$, $R\ge0$ with $A_0\log\max(A_0,1)+R\le CU$ and, for all $Z$,
\[
\lifttwobin\,\mathrm{dg}_Z(u_1)\bigl(T\,\mathrm{ht}_Z(u_1)-3(\lifttwodF+1)|\mathrm{ht}_Z(u_1)|-A_0\log(2+|\mathrm{ht}_Z(u_1)|)-R\bigr)
\le\lifttwohdeg_Z(u_M).
\]
\formallabel{GlobalTheory.CovTop.HSAt}
\item[$\lifttwoP{SumAt}$] $2\,\Sigma_Z\lifttwohdeg_Z(u_M)/\Sigma_Z\mathrm{rk}_Z(u_M)\le\mu_Y(u_M)$.
\formallabel{GlobalTheory.CovTop.SumAt}
\item[$\lifttwoP{UAt}(C)$] $-CU\le\mu_Y(u_M)$.
\formallabel{GlobalTheory.CovTop.UAt}
\item[$\lifttwoP{RingAt}$] $\mu_{\nu^Y_{G_k}}(\Lambda^{\mathrm{deep}}_Y)-2b\le\mu_{\mathrm{core}}$.
\formal{GlobalTheory.CovTop.RingAt}
\item[$\lifttwoP{MapsAt}$] there are $n\in\N$ and $F_Y$-linear maps $\varphi_1,\dots,\varphi_n\colon V_Y(G_s(-c_0))\to
V_Y(G_k)$ with $\varphi_h(\Lambda_Y(G_s(-c_0)))\subseteq\Lambda^{\mathrm{deep}}_Y$, $\bigcap_h\ker\varphi_h=0$,
$\sum_h\nu^Y_{G_k,\sigma}(\varphi_hf)^2\le e^{2b}\nu^Y_{G_s(-c_0),\sigma}(f)^2$ for all $\sigma,f$, and
$n\,\lifttwork\Lambda^{\mathrm{deep}}_Y\le\lifttwork_Y(G_s(c_0))$.
\formallabel{GlobalTheory.CovTop.MapsAt}
\item[$\lifttwoP{RkAt}$] $\lifttwork_Y(G_s(c_0))\le\bigl(1+b/(P\kappa)\bigr)\lifttwork_Y(G_s(-c_0))$.
\formallabel{GlobalTheory.CovTop.RkAt}
\item[$\lifttwoP{GminAt}$] there is a subset $S\subseteq\Lambda^{\mathrm{deep}}_Y$ with
$\nu^Y_{G_k,\sigma}(y)\le e^{P\kappa}$ for all $y\in S$ and all $\sigma$, whose $F_Y$-span contains
$\Lambda^{\mathrm{deep}}_Y$.
\formallabel{GlobalTheory.CovTop.GminAt}
\item[$\lifttwoP{FlatAt}(C)$] $\mu_Y(u_M)-\lifttwomainT-CU\le\mu_Y(G_s(-c_0))$.
\formal{GlobalTheory.CovTop.FlatAt}
\end{description}
\end{definition}

\begin{definition}[The conditions]\label{pf:lift2:def:conds}
The conditions are the following lift conditions ($P=P(x,\iota)$ and $C$ the existential constant of
\cref{pf:lift2:def:lifts}):
\begin{center}\small
\begin{tabular}{@{}lll@{\qquad}lll@{}}
\toprule
condition & definition & & condition & definition &\\
\midrule
$\lifttwoC{Gram}$ & $\lifttwoL(\lifttwoP{GramAt})$ & & $\lifttwoC{Size}$ & $\lifttwoL(\lifttwoP{SizeAt}(P,C))$ &\\
$\lifttwoC{Meas}$ & $\lifttwoL(\lifttwoP{MeasAt}(C))$ & & $\lifttwoC{Anchor}$ & $\lifttwoL(\lifttwoP{AnchorAt}(C))$ &\\
$\lifttwoC{Rat}$ & $\lifttwoL(\lifttwoP{RatAt})$ & & $\lifttwoC{Core}$ & $\lifttwoL(\lifttwoP{CoreAt}(C))$ &\\
$\lifttwoC{Rec}$ & $\lifttwoL(\lifttwoP{RecAt})$ & & $\lifttwoC{Cut}$ & $\lifttwoL(\lifttwoP{CutAt})$ &\\
$\lifttwoC{Bad}$ & $\lifttwoL(\lifttwoP{BadAt})$ & & $\lifttwoC{Lam}$ & $\lifttwoL(\lifttwoP{LamAt}(C))$ &\\
$\lifttwoC{Low}$ & $\lifttwoL(\lifttwoP{LowAt}(C))$ & & $\lifttwoC{LowDef}$ & $\lifttwoL(j=m\Rightarrow\lifttwoP{LowAt}(C))$ &\\
$\lifttwoC{DeepEx}$ & $\lifttwoL(j<m\Rightarrow\exists Y\text{ good})$ & & $\lifttwoC{DeepGood}$ & $\lifttwoCv(\lifttwoP{DeepGoodAt})$ &\\
$\lifttwoC{FlatDeg}$ & $\lifttwoCv(\lifttwoP{FlatDegAt})$ & & $\lifttwoC{Deg}$ & $\lifttwoCv(\forall u_1\in\mathcal Q\colon\lifttwoP{DegAt}(u_1))$ &\\
$\lifttwoC{Share}$ & $\lifttwoCv(\forall u_1\in\mathcal Q\colon\lifttwoP{ShareAt}(u_1))$ & & $\lifttwoC{Height}$ & $\lifttwoCv(\forall u_1\in\mathcal Q\colon\lifttwoP{HeightAt}(C,u_1))$ &\\
$\lifttwoC{HS}$ & $\lifttwoCv(\forall u_1\in\mathcal Q\colon\lifttwoP{HSAt}(C,u_1))$ & & $\lifttwoC{Sum}$ & $\lifttwoCv(\lifttwoP{SumAt})$ &\\
$\lifttwoC{ULow}$ & $\lifttwoCv(\lifttwoP{UAt}(C))$ & & $\lifttwoC{Ring}$ & $\lifttwoCv(\lifttwoP{RingAt})$ &\\
$\lifttwoC{Maps}$ & $\lifttwoCv(\lifttwoP{MapsAt})$ & & $\lifttwoC{Rk}$ & $\lifttwoCv(\lifttwoP{RkAt})$ &\\
$\lifttwoC{Gmin}$ & $\lifttwoCv(\lifttwoP{GminAt})$ & & $\lifttwoC{Flat}$ & $\lifttwoCv(\lifttwoP{FlatAt}(C))$ &\\
\bottomrule
\end{tabular}
\end{center}
The formal counterparts are listed in the order of the table, row by row.
\formallabel{GlobalTheory.CovTop.GramCondition}
\formallabel{GlobalTheory.CovTop.SizeCondition}
\formallabel{GlobalTheory.CovTop.MeasureCondition}
\formallabel{GlobalTheory.CovTop.AnchorCondition}
\formallabel{GlobalTheory.CovTop.LiftRational}
\formallabel{GlobalTheory.CovTop.CoreBound'}
\formallabel{GlobalTheory.CovTop.LiftRecursive}
\formallabel{GlobalTheory.CovTop.LiftTruncation}
\formallabel{GlobalTheory.CovTop.LiftBadSet}
\formallabel{GlobalTheory.CovTop.LiftLambdaBound}
\formallabel{GlobalTheory.CovTop.LiftLowBound}
\formallabel{GlobalTheory.CovTop.LowDefectCondition}
\formallabel{GlobalTheory.CovTop.DeepExistence}
\formallabel{GlobalTheory.CovTop.DeepGoodness}
\formallabel{GlobalTheory.CovTop.FlatDegreeCondition}
\formallabel{GlobalTheory.CovTop.CoverDegree}
\formallabel{GlobalTheory.CovTop.ShareCondition}
\formallabel{GlobalTheory.CovTop.HeightCondition}
\formallabel{GlobalTheory.CovTop.CoverHeightSum}
\formallabel{GlobalTheory.CovTop.SumBoundCondition}
\formallabel{GlobalTheory.CovTop.UniformLowerBound}
\formallabel{GlobalTheory.CovTop.RingCondition}
\formallabel{GlobalTheory.CovTop.MapsCondition}
\formallabel{GlobalTheory.CovTop.CoverRank}
\formallabel{GlobalTheory.CovTop.MinimalGenerator}
\formallabel{GlobalTheory.CovTop.FlatnessCondition}
\end{definition}

The eight lift conditions of Part~I are $\lifttwoC{Gram}$, $\lifttwoC{Rec}$, $\lifttwoC{Cut}$, $\lifttwoC{Bad}$,
$\lifttwoC{Lam}$, $\lifttwoC{Size}$, $\lifttwoC{Core}$ and $\lifttwoC{Low}$.

% ---------------------------------------------------------------------------
\subsection{Numerical properties of the families}\label{sec:pf:lift2:numeric}

Recall $b'=xb/\exp(x^{\eps/3})$ and $w''(n)=b'/(xn)$ (scale data of the size data), where $n$ is the order
parameter of the ranks of $D$.

\begin{definition}[Host and weight numerics]\label{pf:lift2:def:num}
Let $P,A,M\in\R$. The \emph{host numerics} $\mathcal H(P,A,M)$ of $(D,L)$ is the conjunction of: $M\le\ell$,
$10\le\ell\le x$, $1\le\log\log x$; $8\le f$; $[E:\Q]=2f$; $\log(fr)\le\ell$ and $\log n_V\le\ell$; $\log g_k\le9\ell+1$
for $k=0,1$; $\log|d_E|\le4f\ell$; $\log d\le\ell$; the splitting bound of the host is $\ge5$; $\gcd(15,g_0g_1)=1$;
$z^{12}\le P$; $1\le A$, $P\le x^A$ and $f\,r\,n_V\,\ell\le x^A$; $(x^A)^M\le b$; and $\sqrt b\le w''(n)$.

The \emph{weight numerics} $\mathcal N(P)$ of $(p,i)$ is the conjunction of: $(p_j,q_j)=(k+j-1,m-j)$ and $p_j+q_j=r$;
$|\Psi|=f$; $\lifttwodF=f\binom r2$; $n_i=n_V$; the weight of $\kappa_i$ is dominant; $U=f\,m\,b\,\ell\log\log x$;
$|\log\Omega(p)|\le3\ell$; $|\lambda_{\tau,a}|\le4b$ for all $\tau,a$; $w''(n)\le2a_\tau$ for all $\tau$;
$P^2rmk\lifttwodF^3\le b$ and $PT\lifttwodF^2\ell\le b$; and, if $j<m$, there are $h_0\in\Z$ and $S,\eta\ge0$ with
$P(S+\eta)\le100w''(n)$, $4T^2\le h_0$ and $PTrf^2m^2k\lifttwodF\le h_0$, such that for every weight $u_1$ the record of
divided weights built from $(k-1,p_j,T,\lifttwodF,c_0,h_0,S,\eta,u_W,u_1)$ satisfies the hypotheses of the
divided-weight construction of the development (among them $4\lifttwodF+4\le T$, $24NT\le h_0$, lower bounds for the
consecutive gaps of $u_W$ and the bound $|u_{W;\tau,a}|+c_0+4N\le2b$).
\formallabel{GlobalTheory.CovTop.HostNum}
\formallabel{GlobalTheory.CovTop.WtNum}
\formal{GlobalTheory.MeanTop.Wts.Hyp}
\end{definition}

\begin{proposition}[Numerics hold eventually]\label{pf:lift2:numholds}
\begin{enumerate}[label=(\alph*)]
\item For every $M\in\R$ the lift condition of $\mathcal H(P,C,M)$ holds; here the existential constant $C$ plays
the role of $A$. In particular, eventually $M\le\ell$, $10\le\ell$, $1\le\log\log x$, $5\le x^{e_z/5}$ and
$x\ge10^{60}$.
\formallabel{GlobalTheory.CovTop.HostNAlt}
\formallabel{GlobalTheory.CovTop.Sizes.hostN_alt}
\formallabel{GlobalTheory.CovTop.Sizes.eventually_scale}
\item The lift condition of $\mathcal N(P)$ holds.
\formallabel{GlobalTheory.CovTop.LiftWeightNumerator}
\formallabel{GlobalTheory.CovTop.WtShape.wtN_z}
\formallabel{GlobalTheory.T.lift_weight_numerator}
\end{enumerate}
\end{proposition}

\begin{proof}
(a) Take $\eps_1=1$ and any transported constant. Given $F$, put $A=22+|C_t(F)|$, with $C_t(F)$ a constant of the
family. Since $x^{\eps}\to\infty$ and $\log(x^{AM})=o(x^\eps)$, for large $x$ we have $M\le\ell$, $10\le\ell$,
$1\le\log\log x$, $(x^A)^M\le e^{x^\eps}=b$, $5\le x^{e_z/5}$ and $x\ge10^{60}$; the bound $\log g_k\le9\ell+1$ comes
from the family hypothesis $g_k\le\exp(9\ell+1)$. The remaining clauses of $\mathcal H$ are consequences of these and
of the parameter data of the family (its parameter choice, the parameters $P$ and the size bounds), by
\lref{GlobalTheory.CovTop.Sizes.hostNum_of_sizes} (deferred), which derives the host numerics from the
eventual scale bounds and the parameter data.
(b) At every admissible point with $x\ge\max(8,3^{4/\eps})$ (so that $x^{\eps/4}\ge3$), the clauses of $\mathcal N$
follow from the parameter data of the family by \lref{GlobalTheory.CovTop.WtShape.wtNum_at} (deferred), which
derives the weight numerics at a single point from its size data; take $C=0$.
\end{proof}

\begin{lemma}[Size bounds at a point]\label{pf:lift2:sizeat}
If $\mathcal H(P,A,1)$ and $\mathcal N(P)$ hold, then $\lifttwoP{SizeAt}(P,A+2)$ holds at $(p,i)$.
\formallabel{GlobalTheory.CovTop.SizeRows.sizeAt_of}
\formallabel{GlobalTheory.CovTop.SizeRows.mainT_le}
\formallabel{GlobalTheory.CovTop.SizeRows.rho_bounds}
\end{lemma}

\begin{proof}
We have $x\ge10$, $1\le x^A\le b$ (as $A\ge1$ and $(x^A)^1\le b$), and $r\le n_V$, $m\le r$, $f\ge8$, $m\ge2$,
$\ell\ge10$. (i) $U=fmb\ell\log\log x\ge b\ge1$. (ii) $fm\ell\le frn_V\ell\le x^A$ and $\log\log x\le x\le x^2$, so
$U\le x^{A+2}b$. (iii) By definition, $\lifttwomainT=\sum_\rho\bigl((k+j-1)-(m-j)\bigr)w_j(\rho)\log\Omega(p)$, the sum
over the $f$ real places $\rho$ of the host datum, where $w_j(\rho)>0$ is the entry $e_i$ of the Harish-Chandra
parameter at $\rho$ (positivity follows from the windows of the size data). Since $|(k+j-1)-(m-j)|\le r$ and
$|\log\Omega|\le3\ell$, $|\lifttwomainT|\le3r\ell\sum_\rho w_j(\rho)$. The size data bound each $w_j(\rho)$ by
$\sum_a|\lambda_{\tau,a}|+r$ for the element $\tau\in\Psi$ corresponding to $\rho$ (the first entry of $\lambda_\tau$ is
the first entry of $\kappa^+_\tau$ plus $n_V/2$, and the window of the size data relates it to $w_j(\rho)$), and
some $w_j(\rho)$ is at least $b$ (a thick place); hence $\sum_\rho w_j(\rho)\le\kappa+fr$ and $b\le\kappa+r$. From
$z^{12}\le P$ and $frr\ell\le x^A\le b$ one gets $6400r^2\le P$, $9r\ell\le P$, $100\le P$ and $20fr\le b$; hence
$b\le2\kappa$, $\sum_\rho w_j(\rho)\le3\kappa$ and $|\lifttwomainT|\le9r\ell\kappa$. (iv) $\kappa\le4bfr$ (there are
$fr$ entries, each at most $4b$), so $|\lifttwomainT|\le36fr^2\ell b\le36x^Ab\le x^{A+2}b$. (v) $|\lifttwomainT|\le
9r\ell\kappa\le P\kappa$ and $b\le2\kappa\le P\kappa$. (vi) Let $j<m$ and take $h_0,S,\eta$ from $\mathcal N(P)$. A
quotient weight exists: define $u_1$ by summing the prescribed floor quotients of the gaps of $u_W$ from the index
$k-2$. Let $u_1$ be a quotient weight. The hypotheses of the divided-weight construction give $4\lifttwodF+4\le T$,
whence $3\lifttwodF+5\le T$ ($\lifttwodF\ge1$). They also imply that $u_1$ and $u_1-3Nh$ are strictly decreasing in
$a$, that for $a<l$ the ratio
$\bigl((u_1-3Nh)_{\tau,a}-(u_1-3Nh)_{\tau,l}\bigr)/(u_{1;\tau,a}-u_{1;\tau,l})$ lies in $[1-6NT/h_0,1]$ (it equals $1$
unless $a<p_j\le l$, and in that case the gap of $u_1$ across $p_j$ is at least $h_0/(2T)$), and hence, multiplying
the at most $|\Psi|N^2$ ratios, $0\le\rho_B(u_1)<1$ and $\rho_B(u_1)h_0\le6|\Psi|N^3T=6fr^3T$. With
$h_0\ge PTrf^2m^2k\lifttwodF\ge6400\,r^3fT\lifttwodF$ this gives $6400\lifttwodF\rho_B(u_1)\le6$; and
$T\ge P^2\lifttwodF^2\ge100\lifttwodF^2$. Hence $\lifttwozed(u_1)\le7(\frac1{100}+\frac6{6400})<1$.
\end{proof}

\begin{proposition}[The size condition]\label{pf:lift2:size}
$\lifttwoC{Size}$ holds.
\formallabel{GlobalTheory.T.cov_size_bound}
\formallabel{GlobalTheory.CovTop.SizeRows.size_alt_bis}
\formallabel{GlobalTheory.CovTop.SizeRows.size_of_alt}
\end{proposition}

\begin{proof}
By \cref{pf:lift2:numholds} with $M=1$ and \cref{pf:lift2:sizeat}, with $C=A+2$ and \cref{pf:lift2:combine}(c).
\end{proof}

% ---------------------------------------------------------------------------
\subsection{Conditions at a point}\label{sec:pf:lift2:point}

\begin{proposition}[Gram condition]\label{pf:lift2:gram}
At every admissible point and every $i$, $\lifttwoP{GramAt}$ holds; in particular $\lifttwoC{Gram}$ holds.
\formallabel{GlobalTheory.T.covering_gram_of_prints}
\formallabel{GlobalTheory.CovTop.covering_gram_matrix_of_local_traces_alt}
\formallabel{GlobalTheory.CovTop.GramPt.hasGram_of_unif}
\formallabel{GlobalTheory.T.row_norm_has_gram_of_conditions}
\formallabel{GlobalTheory.T.row_lattice_fg_of_conditions}
\formallabel{GlobalTheory.T.row_good_lattice_ne_bot_of_conditions}
\end{proposition}

\begin{proof}
(1) $\nu$ has a Gram frame: the primitive-rows condition contained in $\mathrm{PointConditions}'$ provides, for each
infinite place $v$ of $K$, the realisation $W_i\to H_v$ into a complex inner product space, semilinear for the
embedding of $v$ or its conjugate and mapping independent families to independent families, with
$\nu_{\sigma_v}(f)$ the norm of the realisation; apply \cref{pf:lift2:hasgram}.
(2) $\lifttwoLat{all}$ is finitely generated: there is an additive map $\iota\colon V_T(G_k)\to W_i$, semilinear for
$F_T\to K$, with $\iota(f)=e_i(f)$, mapping independent families to independent families, and with
$\iota(\Lambda_T(G_k))=M_K(W_i)$ (\lref{GlobalTheory.T.exists_semilinear_emb_supLattice_onto_realisation},
deferred: the tie of the point identifies the integral structure of $W_i$ with the image of the core lattice);
$\Lambda_T(G_k)$ is finitely generated by the laws of the core datum; and the $\OO_K$-span of the image of a finitely
generated $\OO_{F_T}$-module under a semilinear map is finitely generated (it is spanned by the images of
generators).
(3) $\lifttwoLat{good}\ne0$: the primitive-rows condition gives $d_W(i)=D>0$ for the base side, and that
$M_i=\{c\in K^{D}:\sum_kc_kb_k\in\lifttwoLat{good}\}$ is a lattice in $K^D$, where $(b_k)$ is a basis of the space
of good lifted forms of the given type; a lattice in $K^D$ with $D>0$ is nonzero, and $c\mapsto\sum c_kb_k$ is
injective.
(4) $\mathcal G_T(G_k)$ holds for the core datum by
\lref{GlobalTheory.T.row_good_weight_dictionary_of_conditions} (deferred), which states that the Petersson norms
of the core datum at the point have a Gram frame and that its lattice of integral forms is finitely generated and
spanning; its proof constructs the Gram frame from the $L^2$ inner products on adelic quotients and the linear
independence of the corresponding adelic functions. The four statements hold at all admissible points, so
$\lifttwoC{Gram}$ holds with $\eps_1=1$, $C=0$ and $x_0=0$.
\end{proof}

\begin{remark}\label{pf:lift2:gramrem}
The formal statement of \cref{pf:lift2:gram} assumes the Hodge package \textup{(Hg)}, but its proof does not use it: the
Gram condition is a consequence of the point conditions alone.
\end{remark}

\begin{definition}[Measure ratio]\label{pf:lift2:def:rat}
For an admissible point $p$ and $i\in J_m$ let $\mathrm{rat}(p,i)=\mu_W(\mathcal F_W)\prod_\tau d_\tau$, where $\mu_W$ is
the Haar measure of the local unitary group used by the realisation of $W_i$, $\mathcal F_W$ a fundamental domain of
its level subgroup, and $d_\tau$ the dimensions attached to the lowest $K$-types.
\formallabel{GlobalTheory.CovTop.MeasRd.ratAt}
\end{definition}

\begin{proposition}[Measure condition]\label{pf:lift2:meas}
Assume \textup{(Hg)}. Then $\lifttwoC{Meas}$ holds.
\formallabel{GlobalTheory.T.covering_meas_of_prints}
\formallabel{GlobalTheory.CovTop.MeasRd.covering_measure_reduction_alt}
\formallabel{GlobalTheory.CovTop.MeasRd.measAt_of}
\end{proposition}

\begin{proof}
The development proves: (a) for every $D$ the \emph{measure normalisation} (for every level datum the index of the
level subgroup is bounded, the lowest $K$-types have dimension $\ge1$, and for every Haar measure and fundamental
domain the realisation inner products are $\mathrm{rat}$ times the probability-normalised Petersson inner products,
with explicit upper bounds for $\log\mathrm{rat}$ under the size bounds of the data)
(\lref{GlobalTheory.T.normalized_covolume_eq_one}, deferred; it uses the compatibility of Tamagawa measures,
covolume ratios equal to indices and bounds for Weyl dimensions); (b) consequently, at every admissible point,
$\nu_\sigma(e_i(f))=\sqrt{\mathrm{rat}(p,i)}\,\nu^T_{G_k,\sigma|F_T}(f)$ for all $\sigma$ and $f$
(\lref{GlobalTheory.T.exactAt_of_oneMeasure}, deferred: the measure normalisation identifies the two norms up to
the factor $\sqrt{\mathrm{rat}}$); (c) the closed bounds $0<\mathrm{rat}(p,i)$, $\log\mathrm{rat}(p,i)\le K_1U$
(\lref{GlobalTheory.T.exists_log_ratAt_le_mul_unitE}, deferred: an upper bound for the measure ratio by a
multiple of the error unit) and $-K_2U\le\log\mathrm{rat}(p,i)$
(\lref{GlobalTheory.T.exists_neg_unitE_le_log_ratAt}, deferred: the corresponding lower bound; it compares the
Haar measure with a product of local masses through the lower bound for nef shares). Given (b) and (c),
$c=\sqrt{\mathrm{rat}}$ satisfies $|\log c|=\frac12|\log\mathrm{rat}|\le\frac{K_1+K_2}2U$ (if $U<0$, the two bounds
force $K_1=K_2=0$ and $\log\mathrm{rat}=0$), so $\lifttwoP{MeasAt}((K_1+K_2)/2)$ holds. Conclude by
\cref{pf:lift2:combine}(c).
\end{proof}

\begin{proposition}[Rationality]\label{pf:lift2:rat}
At every admissible point, $\lifttwoP{RatAt}$ holds; in particular $\lifttwoC{Rat}$ holds.
\formallabel{GlobalTheory.T.liftRational_of_prints}
\formallabel{GlobalTheory.CovTop.covering_rational_of_local_traces_alt}
\formallabel{GlobalTheory.CovTop.ratAt_of_rows_alt}
\end{proposition}

\begin{proof}
The development proves the following statement on forms
(\lref{GlobalTheory.T.integralForms_eq_span_integralRationalForms}, deferred): if the $K$-space of forms has a
rational structure over the subfield $F_T$ which is preserved by the level group $P_0$, and the products
$\chi(u)\psi_0(u)^{-1}$ of the character of the action and the inverse of the type character take values in $F_T$, then
the lattice of integral forms of type $(P_0,\psi_0^{-1})$ is spanned by the integral forms of that type. At an
admissible point the rational structure is given by the core datum, it is preserved by the group action, and the
values $\chi(u)\psi_0(u)^{-1}$ lie in $F_T$ by the point conditions and $\mathrm{RecTied}$. This holds at all
admissible points; take $\eps_1=1$, $C=0$, $x_0=0$.
\end{proof}

\begin{proposition}[Anchor condition]\label{pf:lift2:anchor}
$\lifttwoC{Anchor}$ holds; more precisely, if $\mathcal N(P)$ and $\mathcal H(P',A,1)$ hold, then
$\lifttwoP{AnchorAt}(100)$ holds at $(p,i)$.
\formallabel{GlobalTheory.T.anchor_bound}
\formallabel{GlobalTheory.CovTop.anchor_holds_alt}
\formallabel{GlobalTheory.CovTop.AnchorRd.anchorAt_of}
\end{proposition}

\begin{proof}
The last clause $frn_V\ell\le b$ is $frn_V\ell\le x^A\le b$. Since $(k+j-1)-(m-j)=p_j-q_j$, and the real places of the
host datum are in bijection with $\Psi$ compatibly with the Harish-Chandra entries,
$\lifttwomainT=\sum_{\tau}(p_j-q_j)w_j(\tau)\log\Omega$. The size data place each anchor $a_\tau=\lambda_{\tau,k-2}$
within $S_c+n_V+m+t+2$ of $w_j(\tau)$, where $S_c$ is the cloud width of the size data; the parameter choice gives
$rS_c\le16b$, and $rn_V\le b$. Hence each term of $\lifttwomainT-\sum_\tau(p_j-q_j)a_\tau\log\Omega$ is at most
$r(S_c+n_V+m+t+2)\cdot3\ell\le20b\cdot3\ell$ in absolute value, and the sum is at most $60fb\ell\le100U$. For the
second clause, $|a_\tau|\le4b$, $q_j\le m$ and $|\Psi|=f$ give $(\sum_\tau|a_\tau|)q_j\ell\le4fbm\ell\le100U$. Combine
\cref{pf:lift2:numholds}(a) (with $M=1$) and (b) by \cref{pf:lift2:combine}(c), with $C=100$.
\end{proof}

% ---------------------------------------------------------------------------
\subsection{The lift predicates at a point}\label{sec:pf:lift2:liftpred}

\begin{proposition}[Recursion]\label{pf:lift2:rec}
If $p$ satisfies $\mathrm{PointConditions}'$ and $\lifttwoP{GramAt}$ holds at $(p,i)$, then $\lifttwoP{RecAt}$ holds
at $(p,i)$ (\cref{lem:covolume:rec}). Consequently $\lifttwoC{Gram}$ implies $\lifttwoC{Rec}$.
\formallabel{GlobalTheory.CovTop.recAt_of_gramAt_alt}
\formallabel{GlobalTheory.CovTop.liftRecursive_of_gram}
\formallabel{GlobalTheory.T.lift_recursion_of_gram_bound}
\end{proposition}

\begin{proof}
This is \cref{lem:covolume:rec}; we give the details. The primitive-rows and type-separation conditions give
$d_W(i)=D>0$ (\lref{Formalism.IndefR.Links.complex_data_conditions_base_level_of_primitive_links}: the primitive
links of a point give positive base level and the complex data conditions); the frames of the lifted forms identify
the frame norm of $X_{W,i}$ on $K^D$ with $s\,\nu$ for $s=n_P^{-1/2}\in(0,1]$, where $n_P\ge1$ is the cardinality of a
finite set of the face (\lref{GlobalTheory.T.liftForms_frames_majorantNorm_eq_petersson}, deferred: the majorant
norm of the lifted frames is the scaled Petersson norm); $\lifttwoLat{good}$ lies in the span of the basis $(b_k)$ of
good lifted forms of the given type; and $M_i$ is a lattice by the primitive-rows condition.
\Cref{pf:lift2:partI}(a) gives the claim. The second statement follows by \cref{pf:lift2:combine}(c), since
$\mathrm{PointConditions}'$ holds at admissible points.
\end{proof}

\begin{proposition}[Core comparison]\label{pf:lift2:core}
If $\lifttwoP{GramAt}$, $\lifttwoP{MeasAt}(C)$ and $\lifttwoP{RatAt}$ hold at $(p,i)$, then $\lifttwoP{CoreAt}(2C)$
holds. Consequently $\lifttwoC{Gram}$, $\lifttwoC{Meas}$ and $\lifttwoC{Rat}$ imply $\lifttwoC{Core}$.
\formallabel{GlobalTheory.CoreSlope.coreAt_of}
\formallabel{GlobalTheory.CoreSlope.coreLe_of_gram_meas_rat_alt}
\formallabel{GlobalTheory.T.coreLe_of_gram_meas_liftRational}
\end{proposition}

\begin{proof}
Let $\iota\colon V_T(G_k)\to W_i$ be the semilinear map of the proof of \cref{pf:lift2:gram}: $\iota(f)=e_i(f)$, $\iota$
preserves linear independence and $\iota(\Lambda_T(G_k))=M_K(W_i)$. The coefficient field $F_T$ is totally complex (no
embedding of it is invariant under complex conjugation), and $\nu^T_{G_k,\bar\sigma}=\nu^T_{G_k,\sigma}$. By
$\lifttwoP{MeasAt}(C)$, $\nu_{\sigma'}(\iota x)=c_{\sigma'}\nu^T_{G_k,\sigma'|F_T}(x)\le e^{CU}\nu^T_{G_k,\sigma'|F_T}(x)$. We
claim $\lifttwoLat{type}=\iota_*\Lambda^{\mathrm{type}}_T$. By $\lifttwoP{RatAt}$, $\lifttwoLat{type}$ is spanned by
$M_K(W_i)\cap W_i^{\mathrm{type}}=\iota(\Lambda_T(G_k))\cap W_i^{\mathrm{type}}$. If $x\in\Lambda_T(G_k)$ and
$\iota(x)\in W_i^{\mathrm{type}}$, then $x$ is itself one of the core forms whose image is of the given type, so
$x\in\Lambda^{\mathrm{type}}_T$. Conversely, $\iota$ maps the $F_T$-span of the core forms of the given type into
$W_i^{\mathrm{type}}$ (a $K$-subspace), so $\iota(\Lambda^{\mathrm{type}}_T)\subseteq M_K(W_i)\cap
W_i^{\mathrm{type}}$. This proves the claim. In particular $\Lambda^{\mathrm{type}}_T\ne0$, since
$0\ne\lifttwoLat{good}\subseteq\lifttwoLat{type}$, and $\Lambda^{\mathrm{type}}_T$ is finitely generated, being a
submodule of the finitely generated module $\Lambda_T(G_k)$ over the noetherian ring $\OO_{F_T}$. By
\cref{pf:lift2:slopeext} with $c=CU$,
$\mu_{\mathrm{core}}-2CU\le\mu_\nu(\iota_*\Lambda^{\mathrm{type}}_T)=\mu_\nu(\lifttwoLat{type})$. For the conditions take
$2C$ and \cref{pf:lift2:combine}(c).
\end{proof}

\begin{proposition}[Maximal slope]\label{pf:lift2:lam}
Assume \textup{(Hg)}. Then $\lifttwoC{Gram}$ and $\lifttwoC{Meas}$ imply $\lifttwoC{Lam}$.
\formallabel{GlobalTheory.T.liftLambdaBound_of_gram_meas}
\formallabel{GlobalTheory.CovTop.LamAll.covering_lambda_all_places_alt}
\formallabel{GlobalTheory.CovTop.LamAll.slopeMax_all_le}
\end{proposition}

\begin{proof}
The development proves the closed bound $\mu_{\max,\nu^T_{G_k}}(\Lambda_T(G_k))\le x^{A_1}b$ for the core lattice
(\lref{UpperTheory.Lit.Tam.integral_cusp_slope_power_bound_alt_alt}, deferred; its proof compares the
Petersson norms with Tamagawa measures and bounds the mean logarithmic loss of sup norms linearly in the weight).
Fix $\eps$, $F$, $C_H$ and let $A_1$, $C_m$ (from $\lifttwoC{Meas}$), $A$ (from \cref{pf:lift2:numholds}(a) with
$M=1$) be the constants of the closed statements. At an admissible point with $x$ large: (1) by $\lifttwoP{GramAt}$,
$\lifttwoP{MeasAt}(C_m)$, $\Lambda^{\mathrm{all}}=\iota_*\Lambda_T(G_k)$ (as in the previous proof) and
\cref{pf:lift2:slopemaxext} with $c=C_mU$, $\mu_{\max,\nu}(\lifttwoLat{all})\le\mu_{\max,\nu^T_{G_k}}(\Lambda_T(G_k))+2C_mU
\le x^{A_1}b+2C_mU$; (2) by the host and weight numerics, $U=fmb\ell\log\log x\le frn_V\ell\cdot b\,x\le x^{A+1}b$;
(3) since $x\ge10$, $1+2C_m\le x^{\lceil2C_m\rceil}$ (Bernoulli's inequality), and therefore
\[
\mu_{\max,\nu}(\lifttwoLat{all})\le x^{A_1}b+2C_mx^{A+1}b\le x^{A_1+A+1}b\,(1+2C_m)\le x^{A_1+A+1+\lceil2C_m\rceil}b.
\]
Hence $\lifttwoP{LamAt}(A_1+A+1+\lceil2C_m\rceil)$ holds.
\end{proof}

\begin{proposition}[Truncation]\label{pf:lift2:cut}
If $p$ is admissible and $\mathcal H(P,A,3)$ holds, then $\lifttwoP{CutAt}$ holds at $(p,i)$. Consequently
$\lifttwoC{Cut}$ holds.
\formallabel{GlobalTheory.CovTop.cutAt_of_rows_alt}
\formallabel{GlobalTheory.T.liftTruncation_of_prints}
\formallabel{GlobalTheory.CovTop.covering_cut_of_local_traces_alt}
\end{proposition}

\begin{proof}
Let $c\colon G_p\to K$ be the coefficient function of the idempotent of the group algebra which defines $e_i$, so that
$e_i(y)=\sum_gc(g)\rho(g)y$ for the action $\rho$ of $G_p$ on $W_i$. At admissible points: $c$ is an idempotent of the
group algebra, $\sum_gc(g)c(g^{-1}k)=c(k)$; it is a class function; $|G_p|c(g)$ is integral; and
$\sigma(c(g^{-1}))=\overline{\sigma(c(g))}$ for every $\sigma$
(\lref{GlobalTheory.T.group_idempotent_coefficients_alt}, deferred: the coefficients of the idempotent defining
the truncation have these four properties). The idempotence of $c$ gives $e_i^2=e_i$. The realisations at all places
are $G_p$-invariant for the inner products (\lref{GlobalTheory.T.row_inner_product_invariant_alt}, deferred: the
inner products of the realisations are invariant under the action of $G_p$); with the last property of $c$ this makes
$e_i$ self-adjoint, hence an orthogonal projection, which is the Pythagorean identity. The action preserves
$M_K(W_i)$, $c$ is a class function and $|G_p|c$ is integral, so $|G_p|e_i$ maps the lattice of integral forms of the
given type into itself. Finally $2\log|G_p|\le b$ under $\mathcal H(P,A,3)$
(\lref{GlobalTheory.T.two_mul_log_card_Gp_le_b}, deferred; $|G_p|$ is bounded in terms of the ranks and the
levels). The condition follows from \cref{pf:lift2:numholds}(a) with $M=3$ and \cref{pf:lift2:combine}(c).
\end{proof}

% ---------------------------------------------------------------------------
\subsection{Deep covers and the uniform lower bound}\label{sec:pf:lift2:deep}

\begin{proposition}[Existence of good deep covers]\label{pf:lift2:deepex}
Assume \textup{(Hg)} (only the ray class field statement of the package is used). Then $\lifttwoC{DeepEx}$ holds.
\formallabel{GlobalTheory.T.covering_deepEx_of_prints}
\formallabel{GlobalTheory.CovTop.DeepExist.exists_deep_covering_level_alt}
\formallabel{GlobalTheory.CovTop.DeepExist.size_le}
\formallabel{GlobalTheory.CovTop.RecTied.core_size_le}
\end{proposition}

\begin{proof}
The development proves (\lref{GlobalTheory.T.exists_good_cover_over_sized_field}, deferred): if the ray class
field statement holds, the core datum satisfies its laws, $1\le j\le m$, the splitting bound of the host is $\ge5$
and $\gcd(15,g_0g_1)=1$, then for every $e\in\N$ there is a deep cover $Y$ of $p$ at level $j$ which is
$(e,o)$-good for $o=\mathrm{fs}(F_T)+20(\log(g_0g_115^e)+\log d+1)$. Let $p$ be admissible with $\mathcal H(P,A,0)$
(which holds eventually by \cref{pf:lift2:numholds}(a)) and $j<m$; take $e=\lfloor3\ell/\log3\rfloor$. The condition
$\mathrm{RecTied}$ bounds $\log[F_T:\Q]$ and $\log|d_{F_T}|/[F_T:\Q]$ by $19fr^2\ell+198f^2\ell$, and twice this is
at most $217z^2$ as $\ell\ge2$; so $\mathrm{fs}(F_T)\le217z^2$. With $\log g_k\le9\ell+1$, $\log d\le\ell$,
$e\log15\le8\ell\le z$ and $\ell\ge10$ one gets $217z^2+20(\log g_0+\log g_1+e\log15+\log d+1)\le391z^2$. Since being
$(e,o)$-good is monotone in $o$, $Y$ is good. Conclude by \cref{pf:lift2:combine}(c).
\end{proof}

\begin{proposition}[Goodness of deep covers]\label{pf:lift2:deepgood}
Assume \textup{(Hg)}, \textup{(GM)}, \textup{(Q)}, \textup{(Sg)}, \textup{(KS)} and \textup{(Lie)}. Then
$\lifttwoC{DeepGood}$ holds.
\formallabel{GlobalTheory.T.deepGood_of_toroidal_and_semiabelian}
\formallabel{GlobalTheory.CovTop.DeepGoodAll.deepGood_alt_alt}
\formallabel{GlobalTheory.CovTop.good_comp_nonempty}
\end{proposition}

\begin{proof}
For $\tau\in\Psi$ and a weight $u$ write $\mathrm{gap}(u_\tau,a)=u_{\tau,a}-u_{\tau,a+1}$. The development proves:
\begin{enumerate}[label=(\roman*)]
\item (\lref{GlobalTheory.T.weight_conditions_at_index'}, deferred) eventually at all admissible points, if
$j<m$: the core datum satisfies its laws, $1\le p_j$, $1\le q_j$, $[E:\Q]>4$, $1\le j\le m$, $n_i$ is even,
$u_{M,\tau}$ is non-increasing with $\mathrm{gap}(u_{M,\tau},p_j-1)\ge1$ for all $\tau$, and for
$s\in\{-c_0,0,c_0\}$ the weight $\lambda+sh$ has the same two properties;
\item (\lref{GlobalTheory.T.rank_pos_negative_weight_bundle_on_covers}, deferred) the cover condition of
$\lifttwork_Y(G_s(-c_0))>0$;
\item (\lref{GlobalTheory.T.covering_shift_preserves_goodAt}, deferred) for a good deep cover, every even $n_i$
and every $s$ such that $\lambda+sh$ is non-increasing with gap $\ge1$ at $p_j-1$, $\mathcal G_Y(G_s(s))$;
\item (\cref{pf:lift2:regdom}) for a good deep cover and every non-increasing weight $u$ with gap $\ge1$ at $p_j-1$:
$\mathcal G_Y(\underline u)$ and $\mathcal G_Z(u)$ for all components $Z$;
\item the set of components is finite (\lref{Formalism.L2.Pj.CompFin.finite_comp}: the generic fibre of a
moduli space has finitely many irreducible components) and, for a good deep cover, nonempty (the moduli space of a
deep cover has a point, since the rank of the core datum at level $j$ is positive).
\end{enumerate}
At a good deep cover of an admissible point with $j<m$, (iii) with $s=0$ (note $G_s(0)=G_k$), $s=-c_0$ and $s=c_0$
gives the first three goodness statements, (iv) with $u=u_M$ the fourth and the statements on components, (ii) the
positivity of the rank and (v) the rest. Combine (i) and (ii) by \cref{pf:lift2:combine}(c).
\end{proof}

\begin{proposition}[Flat degrees]\label{pf:lift2:flatdeg}
Assume \textup{(Hg)}. Then $\lifttwoP{FlatDegAt}$ holds for every good deep cover of every admissible point; in
particular $\lifttwoC{FlatDeg}$ holds.
\formallabel{GlobalTheory.T.flatDeg_of_scholzeWeinstein_of_prints}
\formallabel{GlobalTheory.CovTop.FlatDegProof.flatDeg_alt_alt}
\formallabel{GlobalTheory.CovTop.FlatDegProof.flatDegAt_of_st}
\formallabel{GlobalTheory.T.logPeriod_eq_kroneckerCharacter_logPeriod_of_host}
\formallabel{GlobalTheory.LawsDeep.laws_deep}
\end{proposition}

\begin{proof}
The development proves the \emph{flat degree statement} from the package
(\lref{UpperTheory.Lit.Flat.flat_tensor_swap_isIso}, deferred): for every CM host field $E$ with Kronecker
character $\chi$, every datum, level, moduli space, coefficient datum and uniformisations satisfying the laws, the
constant (flat) sections of the line bundles of weight $-\mathbf 1_\tau$ exist and have degree
$\frac{p-q}2\operatorname{Re}\log\Omega_E(\chi)+q(\log d-\log\mathrm{rd}_E)$, where $\log\Omega_E(\chi)$ is the
logarithmic period attached to $\chi$; its proof goes through residues of the logarithmic derivative of the $L$-function and Serre's GAGA
for sections (\lref{GlobalTheory.T.exists_const_section_kronecker_char}, deferred: the existence of constant
sections with the stated degree for a Kronecker character). At an admissible point the host is a host field and
$\log\Omega(p)=\operatorname{Re}\log\Omega_E(\chi)$ for a Kronecker character $\chi$ modulo $d$ (a consequence of the
laws of the core datum); and the uniformisations of a good deep cover satisfy the laws (the level of a good deep cover is a
multiple of the level of the core datum, good at every prime, and the laws of the moduli space follow from those of
the core datum). Hence $\lifttwoP{FlatDegAt}$ holds at all admissible points, and $\lifttwoC{FlatDeg}$ holds with
$\eps_1=1$, $C=0$, $x_0=0$.
\end{proof}

The next proposition says that the images of the components under the cusp forms of the quotient weight have, up
to the factor $1-\rho_B(u_1)$, the same degree. By \cref{pf:lift2:wtpt}, for $j<m$ every quotient weight is the weight
$u_1=u_1(P(p,i))$ of the canonical weight parameter (\cref{pf:lift2:def:wtpar}); $u_1$ and $u_1'=u_1-3Nh$ are strictly
decreasing in the second variable; and $u_M$ differs from the weight $w_T=\mathrm{wt}_T(P(p,i))=\kappa_{\mathrm{Fl}}+u_2+
(\lifttwodF+T)u_1$ by a weight constant in the second variable, where $\kappa_{\mathrm{Fl}}$ is the canonical weight of
the flag variety and $u_2$ is $u_1$ plus the remainders of the division.

\begin{proposition}[Degrees of components]\label{pf:lift2:deg}
Assume \textup{(X)}, \textup{(Hg)}, \textup{(Sg)}, \textup{(KS)}, \textup{(GM)}, \textup{(Q)} and \textup{(Lie)}.
Then $\lifttwoC{DeepGood}$, $\lifttwoC{Size}$ and $\lifttwoC{FlatDeg}$ imply $\lifttwoC{Deg}$; the constant is $C=0$.
\formallabel{GlobalTheory.T.coverDegree_of_grothendieck_deformation}
\formallabel{GlobalTheory.CovTop.DegProof.covering_degree_eq_alt}
\end{proposition}

\begin{proof}
Let $v_0(Z)$ (the volume constant of $Z$) and $\mathrm{top}_Z(u)=v_0(Z)\lifttwovdm(u)$ be as in
\cref{pf:lift2:def:proport}. At a good deep cover of an admissible point with $j<m$ the development provides:
(a) \emph{proportionality}: for strictly decreasing $u$ the growth rate of $\dim V_Z(\tau u)$ in $\tau$ equals
$\mathrm{top}_Z(u)$, and $v_0(Z)>0$ (\cref{pf:lift2:proport}; this is the proportionality principle of Hirzebruch and
Mumford \cite{Mumford1977}, established through the toroidal compactification);
(b) $v_0(Z)$ does not depend on $Z$ (\lref{GlobalTheory.T.compactification_top_chern_eq_volume}, deferred: the components of a good deep cover
all have the same volume constant, computed as a top Chern number on a toroidal compactification);
(c) \emph{birationality}: if all consecutive gaps of $u$ are at least $(\lifttwodF+3)^2+2N$, the cusp forms of weight $u$
on $Z$ generate the function field of $Z$ (\cref{pf:lift2:biron}); the gap hypothesis for $u_1$ holds because the
canonical weight parameter is large (\cref{pf:lift2:wtbig}) and $G(\lifttwodF,N)\ge(\lifttwodF+3)^2+2N$;
(d) a \emph{valuation multiplicity} estimate comparing the weights $u_1$ and $u_1-3Nh$ on every component
(\lref{GlobalTheory.T.valuation_multiplicity_run_estimate_alt}, a cover condition, deferred).
Under (a)--(d) and the strict monotonicity of $u_1$ and $u_1-3Nh$, the development shows for all components
$Z,Z'$ (\lref{GlobalTheory.T.componentImageDegree_between_of_proportionality}, deferred)
\[
\mathrm{top}_{Z'}(u_1-3Nh)\le\mathrm{dg}_Z(u_1)\le\mathrm{top}_{Z'}(u_1-3Nh)/(1-\rho_B(u_1)).
\]
Fix a component $Z'$ (they exist) and take $\delta=\mathrm{top}_{Z'}(u_1-3Nh)$; it is positive since $v_0(Z')>0$ and
the Weyl numerator of a strictly decreasing weight is positive. This is $\lifttwoP{DegAt}(u_1)$. The closed statements
used are (d), the large gaps of $u_1$, the properties of $u_1$ from $\lifttwoC{Size}$, the size bounds for deep
covers (\cref{pf:lift2:sizecov}) and $\lifttwoC{DeepGood}$; combine them by \cref{pf:lift2:combine}(c).
\end{proof}

Note that $\mathrm{top}_{Z'}(u_1)=\mathrm{top}_{Z'}(u_1-3Nh)/(1-\rho_B(u_1))$ by definition of $\rho_B$: the degree of
the image lies between the volumes of the line bundles of weights $u_1-3Nh$ and $u_1$.

\begin{lemma}[Binomial comparisons]\label{pf:lift2:binom}
\begin{enumerate}[label=(\alph*)]
\item Let $D\le M\le N'$ be integers and $x\ge0$ with $N'-M\le x(N'-D+1)$. Then
$\binom{N'}{k}(1-kx)\le\binom Mk$ for $0\le k\le D$.
\formallabel{GlobalTheory.choose_mul_one_sub_le_choose}
\item Let $D,T,a$ be integers and $\rho,\delta,d$ reals with $2\le D$, $D+1\le T$, $0\le\rho\le\frac12$,
$T+D+3\le a\le(T+D+3)/(1-\rho)$, $\delta,d\ge0$ and $d\le\binom{a+D}D\delta$. Then
\[
\Bigl(1-7\Bigl(\frac{D^2}T+D\rho\Bigr)\Bigr)d\le\binom{T-1}D\delta.
\]
\formallabel{GlobalTheory.share_bound_of_jet_count}
\end{enumerate}
\end{lemma}

\begin{proof}
(a) Induction on $k$; $k=0$ is trivial. Using $\binom{N'}{k+1}(k+1)=\binom{N'}k(N'-k)$ and the same for $M$, it suffices
to show $\binom{N'}k(N'-k)(1-(k+1)x)\le\binom Mk(M-k)$. By induction $\binom{N'}k(1-kx)(M-k)\le\binom Mk(M-k)$, and the
difference $(N'-k)(1-(k+1)x)-(1-kx)(M-k)=(N'-M)-x(N'-k)-xk(N'-M)$ is $\le0$ because
$x(N'-k)\ge x(N'-D+1)\ge N'-M$ (as $k+1\le D$).
(b) Put $x=4D/T+\rho\ge0$. Apply (a) with $N'=a+D$, $M=T-1$ and $k=D$: the hypothesis $(a+D)-(T-1)\le
x(a+1)$ follows from $a(1-\rho)\le T+D+3$ and $4D\le\frac{4D}T(a+1)$ (as $a+1\ge T$). Hence
$\binom{a+D}D(1-Dx)\le\binom{T-1}D$. Since $1-7(D^2/T+D\rho)\le1-D(4D/T+\rho)=1-Dx$, if the left side of the claim is
$\le0$ there is nothing to prove; otherwise $\bigl(1-7(\frac{D^2}T+D\rho)\bigr)d\le(1-Dx)\binom{a+D}D\delta\le
\binom{T-1}D\delta$.
\end{proof}

\begin{proposition}[Shares]\label{pf:lift2:share}
Under the hypotheses of \cref{pf:lift2:deg}, $\lifttwoC{DeepGood}$, $\lifttwoC{Size}$, $\lifttwoC{FlatDeg}$ and
$\lifttwoC{Deg}$ imply $\lifttwoC{Share}$.
\formallabel{GlobalTheory.T.share_of_deep_good_size_flat_deg_cover_degree}
\formallabel{GlobalTheory.CovTop.ShareRowsNef.share_alt_alt}
\formallabel{GlobalTheory.CovTop.ShareRowsNef.nef_share_lower_bound}
\formallabel{GlobalTheory.CovTop.ShareRowsNef.nef_share_upper_bound}
\end{proposition}

\begin{proof}
Fix a good deep cover $Y$ of an admissible point with $j<m$, a quotient weight $u_1$ (the canonical one) and a
component $Z$. \emph{Lower bound.} By the flat degree condition, twisting by a weight constant in the second variable
does not change the ranks on components (\cref{pf:lift2:twistrank}(a)), so
$\mathrm{rk}_Z(u_M)=\mathrm{rk}_Z(w_T)$. The cusp forms of weight $u_1$ are birational on $Z$
(\cref{pf:lift2:biron12}), and the weights satisfy
the large-gap condition; under these two hypotheses $\lifttwobin\,\mathrm{dg}_Z(u_1)\le\mathrm{rk}_Z(w_T)$. The formal
proof of this inequality bounds the rank from below by counting sections of the nef and big Hodge line bundle on a
resolution of the image (Hilbert--Samuel for the birational map,
\lref{GlobalTheory.T.choose_mul_componentDegree_le_dimension}, deferred: if the relevant Hodge line bundle is nef
and big and the map given by the forms of weight $u_1$ is birational, then the dimension of the space of cusp forms
of weight $w_T$ on $Z$ is at least $\lifttwobin\,\mathrm{dg}_Z(u_1)$).
\emph{Upper bound.} Put $\rho=\rho_B(u_1)\in[0,\frac12]$ (by $\lifttwozed(u_1)<1$ we even have $7\lifttwodF\rho<1$),
$n=T+\lifttwodF+3$ and $a=\lfloor n/(1-\rho)\rfloor\ge n$. By \cref{pf:lift2:deg} applied with $Z'=Z$,
$n\,\mathrm{top}_Z(u_1)\le n\,\mathrm{dg}_Z(u_1)/(1-\rho)<(a+1)\mathrm{dg}_Z(u_1)$ (as $\mathrm{dg}_Z(u_1)>0$ by
$\lifttwoP{DegAt}$). The \emph{jet bound} (\lref{GlobalTheory.T.jet_rank_complex_le}, deferred) states: if all
consecutive gaps of $u_1$ are large, $u_M$ is non-increasing, $n u_1-u_M$ is non-increasing and $n\,\mathrm{top}_Z(u_1)<
(a+1)\mathrm{dg}_Z(u_1)$, then $\mathrm{rk}_Z(u_M)\le\binom{a+\lifttwodF}{\lifttwodF}\mathrm{dg}_Z(u_1)$. Since
$\lifttwodF\ge2$ and $\lifttwodF+1\le T$, \cref{pf:lift2:binom}(b) gives
$(1-\lifttwozed(u_1))\mathrm{rk}_Z(u_M)\le\lifttwobin\,\mathrm{dg}_Z(u_1)$. The monotonicity hypotheses come from
\cref{pf:lift2:wtpt}(c),(e). Combine the closed statements by
\cref{pf:lift2:combine}(c).
\end{proof}

\begin{proposition}[Sum of component degrees]\label{pf:lift2:sum}
Let $Y$ be a good deep cover of an admissible point and $u$ a weight with $\mathcal G_Y(\underline u)$, finitely
many components and $\mathcal G_Z(u)$ for all $Z$. Then
$2\,\Sigma_Z\lifttwohdeg_Z(u)/\Sigma_Z\mathrm{rk}_Z(u)\le\mu_Y(u)$. Consequently $\lifttwoC{DeepGood}$ implies
$\lifttwoC{Sum}$.
\formallabel{GlobalTheory.CovSum.sum_deg_div_le_mu}
\formallabel{GlobalTheory.T.sum_component_degrees_div_le_latticeSlope}
\formallabel{GlobalTheory.CovTop.sumC_of_deepGood_alt}
\formallabel{GlobalTheory.T.sum_c_of_deep_good}
\formallabel{GlobalTheory.CovSum.log_two_le}
\formallabel{GlobalTheory.CovSum.slope_sq_nonneg_of_rk_zero}
\end{proposition}

\begin{proof}
The development proves that $\Lambda_Y(\underline u)$ has rank $\Sigma_Z\mathrm{rk}_Z(u)$ and degree at least the sum of
the degrees of the component lattices: $\Sigma_Z\lifttwohdeg_Z(u)\le\lifttwohdeg_{\nu^Y_u}(\Lambda_Y(\underline u))$
(\lref{GlobalTheory.T.componentDegrees_sum_le_total_degree}, deferred; the space of cusp forms is the direct sum of
the spaces supported on the components). If the rank $\rho_u$ of $\Lambda_Y(\underline u)$ is positive, then
$\lifttwohdeg=\frac12\rho_u\mu_Y(u)$ by \eqref{eq:pf:lift2:deg}, and dividing gives the claim. If $\rho_u=0$, the
left side is $0$ (division by $0$ is $0$ in the development), and $\mu_Y(u)$ is the slope of a module of rank $0$,
namely $-\log2+\log|d_{F_Y}|/[F_Y:\Q]$, which is $\ge0$ because $F_Y$ contains a conjugate of the host field $E$ and
hence $\log|d_{F_Y}|/[F_Y:\Q]\ge\log2$. For the condition apply this with $u=u_M$; the hypotheses are parts of
$\lifttwoP{DeepGoodAt}$ and of the goodness of $Y$.
\end{proof}

\begin{lemma}[Averaging over components]\label{pf:lift2:avg}
Let $\mathcal Z$ be a nonempty finite set and let $T,d,A_0,c_s,X,c_c,\zeta,\rho,R,B,\delta',U,C_1,C_2,\mu$ be reals and
$h,\mathrm{sl},\delta,\mathrm r,g\colon\mathcal Z\to\R$ functions with $d\ge0$, $T\ge3d+5$, $A_0>0$, $c_s,X,c_c,R\ge0$,
$0\le\zeta<1$, $0\le\rho<1$, $B>0$, $\delta'>0$, $U\ge1$, and for all $Z\in\mathcal Z$:
\begin{enumerate}[label=(\arabic*)]
\item $\delta'\le\delta(Z)\le\delta'/(1-\rho)$;
\item $B\delta(Z)\le\mathrm r(Z)$ and $(1-\zeta)\mathrm r(Z)\le B\delta(Z)$;
\item $0\le h(Z)+c_s$ and $(1-\rho)\max(0,\mathrm{sl}(Z)+c_s)\le h(Z)+c_s+c_c$;
\item $B\delta(Z)\bigl(Th(Z)-3(d+1)|h(Z)|-A_0\log(2+|h(Z)|)-R\bigr)\le g(Z)$;
\end{enumerate}
and moreover $-X\le\frac1{|\mathcal Z|}\sum_Z\mathrm{sl}(Z)$, $TX\le C_1U$, $Tc_c\le C_1U$,
$c_s(6d+10+T(\zeta+2\rho))\le C_1U$, $A_0\log\max(A_0,1)+R\le C_2U$ and $2\sum_Zg(Z)/\sum_Z\mathrm r(Z)\le\mu$. Then
$-2(3C_1+C_2+4)U\le\mu$.
\formallabel{GlobalTheory.CovTop.FromComp.core}
\formallabel{Formalism.ExtensionFiltration.meanHeightLowerBound}
\formallabel{Formalism.ExtensionFiltration.sumComponents'}
\end{lemma}

\begin{proof}
\emph{Step 1 (mean height).} Let $n=|\mathcal Z|$ and $m_0=\max(0,(1-\rho)c_s-X-c_c)$. Then
$m_0\le\frac1n\sum_Z(h(Z)+c_s)$: the right side is
$\ge0$ by (3), and from (3), $(1-\rho)(\mathrm{sl}(Z)+c_s)-c_c\le h(Z)+c_s$; summing and using
$\sum\mathrm{sl}\ge-nX$ gives $\sum(h+c_s)\ge n\bigl((1-\rho)(c_s-X)-c_c\bigr)\ge n\bigl((1-\rho)c_s-X-c_c\bigr)$.

\emph{Step 2 (pointwise).} For $y\ge0$, $A_0\log(2+y)\le y+A_0\log\max(A_0,1)+2$: if $A_0\ge1$ this is
$A_0\log\frac{2+y}{A_0}\le2+y-A_0$, and if $A_0<1$ it follows from $\log(2+y)\le1+y$. Since $h\ge-c_s$ we have
$|h|\le h+2c_s$. Hence, with $N_0=(T-3d-4)c_s+2c_s(3d+4)+A_0\log\max(A_0,1)+2+R$,
\begin{align*}
Th-3(d+1)|h|-A_0\log(2+|h|)-R&\ge Th-(3d+4)(h+2c_s)-A_0\log\max(A_0,1)-2-R\\
&=(T-3d-4)(h+c_s)-N_0,
\end{align*}
and by (4), $B\delta(Z)\bigl((T-3d-4)(h(Z)+c_s)-N_0\bigr)\le g(Z)$.

\emph{Step 3 (summation).} Let $W=\sum_ZB\delta(Z)>0$ and $S=\sum_Z\mathrm r(Z)$; by (2), $W\le S$. By (1) and
$(T-3d-4)(h+c_s)\ge0$,
\[
\sum_ZB\delta(Z)(T-3d-4)(h(Z)+c_s)\ge B\delta'(T-3d-4)\sum_Z(h(Z)+c_s)\ge B\delta'n(T-3d-4)m_0,
\]
and (1) also gives $(1-\rho)W\le nB\delta'$. Hence $\sum_Zg(Z)\ge W\bigl((T-3d-4)m_0(1-\rho)-N_0\bigr)$.

\emph{Step 4.} From $m_0\ge(1-\rho)c_s-X-c_c$ we get $(1-\rho)m_0\ge(1-2\rho)c_s-X-c_c$, so
\begin{align*}
(T-3d-4)m_0(1-\rho)-N_0&\ge-T(X+c_c)-c_s\bigl(6d+10+T(\zeta+2\rho)\bigr)\\
&\quad-\bigl(A_0\log\max(A_0,1)+4\bigr)-R=:-Q,
\end{align*}
using $(T-3d-4)(2\rho c_s+X+c_c)\le T(2\rho c_s+X+c_c)$; note $Q\ge0$.

\emph{Step 5.} If the factor $(T-3d-4)m_0(1-\rho)-N_0$ is $\ge0$, then $\sum g\ge0\ge-QS$; otherwise, since $W\le S$,
$\sum g\ge W\cdot(\text{factor})\ge S\cdot(\text{factor})\ge-QS$. In both cases $\sum g/S\ge-Q$.
Finally $Q\le3C_1U+C_2U+4\le(3C_1+C_2+4)U$ as $U\ge1$, and $\mu\ge2\sum g/S\ge-2(3C_1+C_2+4)U$.
\end{proof}

\begin{proposition}[Uniform lower bound]\label{pf:lift2:ulow}
$\lifttwoC{DeepGood}$, $\lifttwoC{Size}$, $\lifttwoC{Deg}$, $\lifttwoC{Share}$, $\lifttwoC{Height}$, $\lifttwoC{HS}$ and
$\lifttwoC{Sum}$ imply $\lifttwoC{ULow}$. More precisely, if at $(p,i)$ with $j<m$ and a good deep cover $Y$ the
predicates $\lifttwoP{DeepGoodAt}$, $\lifttwoP{SizeAt}(P,A)$, $\lifttwoP{DegAt}(u_1)$, $\lifttwoP{ShareAt}(u_1)$,
$\lifttwoP{HeightAt}(C_1,u_1)$, $\lifttwoP{HSAt}(C_2,u_1)$ (for all quotient weights $u_1$) and $\lifttwoP{SumAt}$ hold,
then $\lifttwoP{UAt}(2(3C_1+C_2+4))$ holds.
\formallabel{GlobalTheory.CovTop.FromComp.atPoint}\formallabel{GlobalTheory.CovTop.FromComp.u_low_of_cover_hyps}
\formallabel{GlobalTheory.T.uniform_lower_bound_of_sums'}
\end{proposition}

\begin{proof}
By $\lifttwoP{SizeAt}$ there is a quotient weight $u_1$, and $3\lifttwodF+5\le T$, $0\le\rho_B(u_1)<1$,
$\lifttwozed(u_1)<1$, $U\ge b\ge1$. Apply \cref{pf:lift2:avg} with $\mathcal Z=\lifttwoComp(Y)$ (finite and nonempty by
$\lifttwoP{DeepGoodAt}$), $d=\lifttwodF$, $\rho=\rho_B(u_1)$, $\zeta=\lifttwozed(u_1)\ge0$, $B=\lifttwobin>0$ (as
$T-1\ge\lifttwodF$), $h=\mathrm{ht}_\bullet(u_1)$, $\delta=\mathrm{dg}_\bullet(u_1)$ and $\delta'$ from
$\lifttwoP{DegAt}$, $\mathrm r=\mathrm{rk}_\bullet(u_M)$, $g=\lifttwohdeg_\bullet(u_M)$, $\mu=\mu_Y(u_M)$, and
$c_s,c_c,X,\mathrm{sl}$ from $\lifttwoP{HeightAt}$, $A_0,R$ from $\lifttwoP{HSAt}$. The hypotheses of the lemma are
exactly the predicates. The condition follows by \cref{pf:lift2:combine}(c) with $C=2(3C_1+C_2+4)$.
\end{proof}

The conditions $\lifttwoC{Height}$ and $\lifttwoC{HS}$ are proved in the companion section on lift estimates:
\lref{SelfIntersectionVolumeEstimates.height_lower_bound_general_alt} states that \textup{(VE)}, \textup{(X)},
\textup{(GM)}, \textup{(Q)}, \textup{(Hg)}, \textup{(H)}, \textup{(DS)}, $\lifttwoC{DeepGood}$, $\lifttwoC{Size}$,
$\lifttwoC{Deg}$ and $\lifttwoC{FlatDeg}$ imply $\lifttwoC{Height}$; and
\lref{GlobalTheory.T.coverHeightSum_of_flatDeg_coverDegree_share} states that \textup{(X)}, \textup{(Hg)},
\textup{(GM)}, \textup{(Q)}, \textup{(Sg)}, \textup{(KS)}, \textup{(Lie)}, $\lifttwoC{DeepGood}$, $\lifttwoC{Size}$,
$\lifttwoC{FlatDeg}$, $\lifttwoC{Deg}$ and $\lifttwoC{Share}$ imply $\lifttwoC{HS}$.

% ---------------------------------------------------------------------------
\subsection{Descent}\label{sec:pf:lift2:descent}

In this subsection $Y$ is a good deep cover of an admissible point $p$ and $j=i+1<m$. Pull-back of cusp forms along
the change of level defines, for every weight $G$, an injective map $\pi_G\colon V_T(G)\to V_Y(G)$, semilinear for the
inclusion $F_T\subseteq F_Y$, which maps $\Lambda_T(G)$ into $\Lambda_Y(G)$ and preserves the norms,
$\nu^Y_{G,\sigma}(\pi_Gx)=\nu^T_{G,\sigma|F_T}(x)$. The finite group $\mathrm{Aut}_Y$ of automorphisms of the level
structure of $Y$ acts on $X_Y$; we write $g^*$ for the pull-back by $g\in\mathrm{Aut}_Y$. For a subgroup $S$ and a
character $\chi\colon S\to F_Y^\times$ we write $\Lambda_Y(G)^{\chi}$ for the lattice of $f\in\Lambda_Y(G)$ with
$g^*f=\chi(g)f$ for $g\in S$.

\begin{proposition}[Comparison with the core]\label{pf:lift2:ring}
Let $Y$ be a deep cover of $p$ at level $j$ which is $(e,o)$-good, and assume $\mathcal G_T(G_k)$ for the core and
$\mathcal G_Y(G_k)$. Then
\[
\mu_{\nu^Y_{G_k}}(\Lambda^{\mathrm{deep}}_Y)-\frac{2\log|d_{F_Y}|}{[F_Y:\Q]}\le\mu_{\mathrm{core}}.
\]
Consequently, under the Hodge package \textup{(Hg)}, $\lifttwoC{Gram}$ and $\lifttwoC{DeepGood}$ imply $\lifttwoC{Ring}$, with
$C=0$.
\formallabel{GlobalTheory.CoreRing.slope_sq_latOver_sub_le}
\formallabel{GlobalTheory.T.cuspFormLattice_slope_sq_goodCover_le_core}
\formallabel{GlobalTheory.CovTop.RingStep.ring_of_covering_descent_alt}
\formallabel{GlobalTheory.T.covering_ring_of_prints_gram_deepGood}
\end{proposition}

\begin{proof}
Let $\pi=\pi_{G_k}$ and let $C\subseteq V_T(G_k)$ be the $F_T$-span of the core forms of the given type, so that
$\Lambda^{\mathrm{type}}_T=\Lambda_T(G_k)\cap C$. The development proves for the level covering
(\lref{GlobalTheory.T.levelCover_cuspForm_pullback_injective_descent}, deferred): $\pi$ has the properties
listed above, and if $g\in\Lambda_Y(G_k)$ lies in the $F_Y$-span of $\pi(C)$, then $\delta g$ lies in the
$\OO_{F_Y}$-span of $\pi(\Lambda_T(G_k)\cap C)$ for every $\delta$ in the different ideal of $F_Y$ (a Galois descent
by traces). Moreover, a cusp form on $Y$ is determined by the condition of being the pull-back of a given core form
(\lref{GlobalTheory.T.existsUnique_cuspForm_pullback_goodCover}, deferred). Hence every element of the
$F_Y$-span of the pull-backs of the core forms of the given type lies in the $F_Y$-span of $\pi(C)$ (by induction on
linear combinations: if $h$ is the pull-back of $f\in C$, then $h-\pi f$ is the pull-back of $0$, hence $0$), and
so $\Lambda^{\mathrm{deep}}_Y=\Lambda_Y(G_k)\cap F_Y\pi(C)$ satisfies the hypotheses of \cref{pf:lift2:slopediff} with
$L=\Lambda^{\mathrm{type}}_T$, $M=\Lambda^{\mathrm{deep}}_Y$ (finitely generated, being a submodule of
$\Lambda_Y(G_k)$) and $\iota=\pi$; $F_T$ is totally complex and $\nu^T$ is invariant under complex conjugation of the
embedding. \Cref{pf:lift2:slopediff} gives the displayed inequality.

For the condition, let $p$ be admissible with $\mathcal H(P,A,4)$ (\cref{pf:lift2:numholds}(a) with $M=4$),
and $Y$ good. Then $\log|d_{F_Y}|/[F_Y:\Q]\le\mathrm{fs}(F_Y)\le391z^2$. From $z\ge1$, $z^{12}\le P\le x^A$ and
$x^A\ge x\ge10$ we get $391z^2\le391x^A\le(x^A)^4\le b$. Hence $\mu(\Lambda^{\mathrm{deep}}_Y)-2b\le\mu_{\mathrm{core}}$.
\end{proof}

\begin{proposition}[Descent maps]\label{pf:lift2:maps}
Assume \textup{(Hg)}, \textup{(GM)}, \textup{(Q)}, \textup{(Sg)}, \textup{(KS)} and \textup{(Lie)}. Then $\lifttwoC{Gram}$, $\lifttwoC{DeepGood}$, $\lifttwoC{Size}$ and
$\lifttwoC{FlatDeg}$ imply $\lifttwoC{Maps}$, with $C=0$.
\formallabel{GlobalTheory.CovTop.DescentMaps.covering_descent_maps_compatible_alt}
\formallabel{GlobalTheory.T.covering_maps_of_gram_deep_size_flatDeg}
\end{proposition}

\begin{proof}
Let $W=\underline{c_0h}$, so that $G_s(-c_0)+W=G_k$ and $G_k+W=G_s(c_0)$. The development provides, eventually at
good deep covers of admissible points:
\begin{enumerate}[label=(\roman*)]
\item a subgroup $S\subseteq\mathrm{Aut}_Y$ and a character $\chi$ of $S$ with $\Lambda^{\mathrm{deep}}_Y=
\Lambda_Y(G_k)^{\chi}$ and $\log|S|\le45z^2$
(\lref{GlobalTheory.T.exists_psi_isTypeLevel_latDeep_eq_latN}, deferred: the type of the core forms is a level
type, read off from a frame of the level);
\item a constant $C_u\ge0$ and an integral cusp form $u\in\Lambda_Y(W)$ such that the determinant of the matrix
$\bigl((hg)^*u\bigr)_{h,g\in S}$ is a non-zero-divisor in the ring of sections, and such that
$\nu^Y_{G+W,\sigma}(f\cdot g^*u)\le\exp(C_uz^9)\,\nu^Y_{G,\sigma}(f)$ for every $g\in\mathrm{Aut}_Y$, every weight $G$ of
even parity and every $f\in V_Y(G)$ (\lref{GlobalTheory.T.exists_short_integral_vector_of_determinant_bound},
deferred; the existence of a form with non-zero-divisor level determinant is
\lref{GlobalTheory.T.orbits_separated_by_regular_cusp}, and the norm bound is a comparison of the norms of short
generators);
\item the level group preserves cusp forms, integral forms and Petersson norms in every weight of even parity
(\lref{GlobalTheory.LevelIsometry.keeps_cusp_lattice_norm}, deferred).
\end{enumerate}
For $f\in V_Y(G_s(-c_0))$ and $h\in S$ put
\[
\varphi_h(f)=\sum_{g\in S}\chi(g)^{-1}g^*\bigl(f\cdot h^*u\bigr)\in V_Y(G_k).
\]
(a) $\varphi_h(f)$ is a $\chi$-eigenform, and it is integral if $f$ is (\lref{GlobalTheory.LevDesc.cuspPh_mem_chiLat},
deferred: $u$ is integral, pull-backs preserve integrality, and the values of $\chi$ are roots of unity since $S$ is
finite); so $\varphi_h(\Lambda_Y(G_s(-c_0)))\subseteq\Lambda^{\mathrm{deep}}_Y$.
(b) Since the action is a right action, $\varphi_h(f)=\sum_g\chi(g)^{-1}(g^*f)\,(hg)^*u$; if $\varphi_h(f)=0$ for all $h$,
then the vector $(\chi(g)^{-1}g^*f)_g$ is annihilated by the matrix $((hg)^*u)_{h,g}$, whose determinant is a
non-zero-divisor, so $f=0$. (c) By (iii) and (ii), $\nu(\varphi_h(f))\le\sum_g\nu(f\cdot h^*u)\le|S|\exp(C_uz^9)\nu(f)$,
hence $\sum_h\nu(\varphi_hf)^2\le|S|^3\exp(C_uz^9)^2\nu(f)^2$. (d) The map $(F_h)_h\mapsto\sum_hF_h\cdot h^*u$ from
$(V_Y(G_k)^\chi)^S$ to $V_Y(G_s(c_0))$ is injective by the same determinant argument, so
$|S|\lifttwork\Lambda^{\mathrm{deep}}_Y\le\lifttwork_Y(G_s(c_0))$ (\lref{GlobalTheory.CovolRank.card_mul_rkK_chiLat_le_rankW},
deferred). It remains to see $|S|^3\exp(2C_uz^9)\le e^{2b}$. Under $\mathcal H(P,A,2C_u+200)$ we have
$z\ge\ell\ge2C_u+200$, and $z^{12}\le P\le x^A\le(x^A)^{2C_u+200}\le b$; with $\log|S|\le45z^2$,
$3\log|S|+2C_uz^9\le135z^2+2C_uz^9\le z^{10}\le z^{12}\le b\le2b$. Thus $\lifttwoP{MapsAt}$ holds with $n=|S|$.
\end{proof}

\begin{proposition}[Rank comparison]\label{pf:lift2:rank}
Assume \textup{(X)}, \textup{(Hg)}, \textup{(Sg)}, \textup{(KS)}, \textup{(GM)}, \textup{(Q)} and \textup{(Lie)}. Then $\lifttwoC{DeepGood}$, $\lifttwoC{Size}$, $\lifttwoC{FlatDeg}$,
$\lifttwoC{Deg}$ and $\lifttwoC{Share}$ imply: eventually, at good deep covers, for $s\in\{1,2\}$,
\[
\lifttwork_Y(G_s(sc_0))\le\Bigl(1+\min\Bigl(\frac12,\frac b{P\kappa}\Bigr)\Bigr)\lifttwork_Y(G_s(-c_0)).
\]
In particular $\lifttwoC{Rk}$ holds, with $C=0$.
\formallabel{GlobalTheory.CovTop.HodgeSteps.rk_le_alt_alt}\formallabel{GlobalTheory.T.rank_growth_bound_small_scale_alt_alt}
\formallabel{GlobalTheory.CovTop.covering_rank_of_descent_steps_alt}\formallabel{GlobalTheory.T.coverRank_of_coverDegree_and_share}
\end{proposition}

\begin{proof}
Let $u_1$ be the quotient weight, $\rho=\rho_B(u_1)$, $\zeta=\lifttwozed(u_1)$ and $d=\lifttwodF$. Two weights that differ
by a weight constant in the second variable have the same component ranks, and the rank of $\Lambda_Y$ in a shifted
weight $G_s(s')$ is the sum over components of the ranks of the weight $u_W+s'h$
(\cref{pf:lift2:twistrank}, which uses $\lifttwoC{FlatDeg}$); in particular
$\lifttwork_Y(G_s(-c_0))=\sum_Z\mathrm{rk}_Z(u_M)$ and $\lifttwork_Y(G_s(sc_0))=\sum_Z\mathrm{rk}_Z(u_M+(s+1)c_0h)$, the
latter weight being good on $Y$ by the regularity statements used in \cref{pf:lift2:deepgood}. The size bounds give
$d\ge3$, $P\ge100$, $T\ge1000d^2$ and $\zeta\le\frac12$, $\zeta P\kappa\le b$
(\cref{pf:lift2:sizecov}). An elementary estimate for binomial coefficients
(choose $a=\lfloor(T+d+4)/(1-\rho)\rfloor$ and compare $\binom{a+d}d$ with $\binom{T-1}d$ through the ratio
$(1+y)^d\le e^{dy}$, $y=(2d+5+\rho T)/((1-\rho)(T-d))$) gives an integer $a$ with
\[
(T+d+4)<(a+1)(1-\rho),\qquad\binom{a+d}d\le(1+\zeta)\binom{T-1}d.
\]
For each component $Z$, $\mathrm{dg}_Z(u_1)>0$ and $(1-\rho)\mathrm{top}_Z(u_1)\le\mathrm{dg}_Z(u_1)$ by
\cref{pf:lift2:deg}, so $(T+d+4)\mathrm{top}_Z(u_1)<(a+1)\mathrm{dg}_Z(u_1)$, and the jet bound
(\lref{GlobalTheory.T.jet_rank_complex_le}) applies to $u_1$ and the weight $u_M+(s+1)c_0h$ (whose monotonicity
properties are part of \cref{pf:lift2:wtpt}(e)):
\[
\mathrm{rk}_Z(u_M+(s+1)c_0h)\le\binom{a+d}d\mathrm{dg}_Z(u_1)\le(1+\zeta)\lifttwobin\,\mathrm{dg}_Z(u_1)
\le(1+\zeta)\mathrm{rk}_Z(u_M),
\]
the last step by $\lifttwoP{ShareAt}$. Summing over $Z$ and using $\zeta\le\min(\frac12,b/(P\kappa))$ gives the claim.
$\lifttwoP{RkAt}$ is the case $s=1$, as $\min(\frac12,b/(P\kappa))\le b/(P\kappa)$.
\end{proof}

\begin{proposition}[Short generators]\label{pf:lift2:gmin}
Assume \textup{(Hg)}. Then $\lifttwoC{Gram}$, $\lifttwoC{DeepGood}$ and $\lifttwoC{Size}$ imply $\lifttwoC{Gmin}$, with $C=0$.
\formallabel{GlobalTheory.CovGminTop.covering_gmin_of_gram_deep_size}
\formallabel{GlobalTheory.T.covering_gmin_of_prints_gram_deep_size}
\end{proposition}

\begin{proof}
\emph{Twisted traces.} Let $H$ be a finite group acting $F_Y$-linearly on $V=V_Y(G_k)$ (on the right), preserving a
lattice $\Lambda$ and a family of norms $\nu$ with a Gram frame, $\nu_{\bar\sigma}=\nu_\sigma$,
$\nu_\sigma(ax)=|\sigma(a)|\nu_\sigma(x)$, and let $\chi\colon H\to F_Y^\times$ be a character with values integral units of
absolute value $1$. If $S_0\subseteq\Lambda$ spans $V$ and $\nu_\sigma(y)\le e^B$ on $S_0$, then the twisted traces
$\mathrm{tr}_\chi(y)=\sum_h\chi(h)^{-1}h^*y$, $y\in S_0$, lie in $\Lambda$, are $\chi$-eigenvectors, satisfy
$\nu_\sigma(\mathrm{tr}_\chi y)\le|H|e^B$ (triangle inequality for norms with a Gram frame), and span the
$\chi$-eigenspace (the trace is $|H|$ times the projection onto it). \emph{Application.} By
\lref{GlobalTheory.T.exists_good_cover_at_minimal_level} (deferred) there is a constant $C_1\ge0$ such that at
good deep covers satisfying the size bounds, $\Lambda_Y(G_k)$ contains a spanning set of forms of norms
$\le\exp(C_1z^2\kappa)$ at all places. With $S$, $\chi$ from \cref{pf:lift2:maps}(i), $\Lambda^{\mathrm{deep}}_Y=\Lambda_Y(G_k)^\chi$,
and the twisted traces give a spanning set of $\Lambda^{\mathrm{deep}}_Y$ inside it with norms $\le|S|\exp(C_1z^2\kappa)$.
Finally $\log|S|+C_1z^2\kappa\le45z^2+C_1z^2\kappa\le P\kappa$, since $\kappa\ge1$, $z^{12}\le P$ and, for $x$ large,
$z\ge\ell\ge C_1+45$. This is $\lifttwoP{GminAt}$.
\end{proof}

\begin{proposition}[Flatness]\label{pf:lift2:flat}
Assume \textup{(Hg)}. Then $\lifttwoC{DeepGood}$, $\lifttwoC{FlatDeg}$ and $\lifttwoC{Anchor}$ imply $\lifttwoC{Flat}$; if
$\lifttwoC{Anchor}$ holds with the constant $C_A$, $\lifttwoC{Flat}$ holds with $11C_A+83$.
\formallabel{GlobalTheory.FlatSplit.flat_at_of_split_covering}
\formallabel{GlobalTheory.FlatSplit.flat_alt_alt}\formallabel{GlobalTheory.T.flat_of_prints_deepGood_flatDeg_anchor}
\end{proposition}

\begin{proof}
Let $c_1(\tau)=-\frac1{|\lifttwoComp(Y)|}\sum_Z\deg_{\mathrm{flat}}(Z,\tau)$ be the average degree of the flat sections
of weight $-\mathbf 1_\tau$. The development proves the twisting statement
(\lref{GlobalTheory.T.flatLine_thetaRoot_twist_twistingSection}, deferred): under the de~Rham statement of \textup{(Hg)}
and the laws of $Y$, for every $w\colon\Psi\to\Z$ and $m_0\in\N$, twisting by the weight constant $w$ (with parity
shift $m_0$) preserves goodness and ranks, and raises degrees by at least the rank times
\[
\theta(w,m_0)=\sum_\tau w_\tau c_1(\tau)-q_j\log\mathrm{rd}_E\sum_\tau|w_\tau|-m_0fr\Bigl(\frac{\log2\pi}2+\log\mathrm{rd}_E\Bigr)
-\sum_{\ell'\in\mathcal B}\log\ell',
\]
where $\mathcal B$ is the set of primes dividing $2d\,d_EM'_Y$. The weight $G_s(-c_0)$ is the twist of $\underline{u_M}$
by the anchor $a$ with parity shift $n_i/2$ (both define the same line bundles), so
$\lifttwork_Y(G_s(-c_0))=\lifttwork_Y(u_M)>0$ and, by \eqref{eq:pf:lift2:deg}, $\mu_Y(G_s(-c_0))\ge\mu_Y(u_M)+2\theta(a,n_i/2)$.
It remains to show $2\theta(a,n_i/2)\ge-\lifttwomainT-(11C_A+83)U$. The development proves the closed statement
$|c_1(\tau)+\frac{p_j-q_j}2\log\Omega(p)|\le3q_j\ell$ and $\sum_{\ell'\in\mathcal B}\log\ell'\le40\ell$
(\lref{GlobalTheory.T.abs_flat_c1_add_logOm_le_and_logSum_badPrimes_le}, deferred; the first part is the flat
degree condition). Then $\sum_\tau a_\tau c_1(\tau)\ge-\frac12\sum_\tau(p_j-q_j)a_\tau\log\Omega-3q_j\ell\sum|a_\tau|$, and
$\lifttwoP{AnchorAt}(C_A)$ bounds $\sum(p_j-q_j)a_\tau\log\Omega$ by $\lifttwomainT+C_AU$ and $q_j\ell\sum|a_\tau|$ by
$C_AU$; $\log\mathrm{rd}_E\le2\ell$ (from $\log|d_E|\le4f\ell$ and $[E:\Q]=2f$), $2m_0\le n_V$, $\log2\pi\le2\pi-1$ give
$m_0fr(\frac{\log2\pi}2+\log\mathrm{rd}_E)\le\frac{3}2frn_V\ell\le\frac32b\le\frac32U$; and $40\ell\le40U$. Collecting,
$2\theta\ge-\lifttwomainT-C_AU-6C_AU-4C_AU-3U-80U$.
\end{proof}

\begin{proposition}[Lower bound for the core slope]\label{pf:lift2:low}
\begin{enumerate}[label=(\alph*)]
\item (\cref{lem:covolume:low}) Let $C_F,C_U\ge0$. If $\lifttwoP{SizeAt}(P,A)$ holds at $(p,i)$ and $Y$ is a deep cover
with $\lifttwoP{DeepGoodAt}$, $\lifttwoP{RingAt}$, $\lifttwoP{MapsAt}$, $\lifttwoP{RkAt}$, $\lifttwoP{GminAt}$,
$\lifttwoP{FlatAt}(C_F)$ and $\lifttwoP{UAt}(C_U)$, then $\lifttwoP{LowAt}(2(C_F+C_U)+9)$ holds at $(p,i)$.
\item $\lifttwoC{Size}$, $\lifttwoC{LowDef}$, $\lifttwoC{DeepEx}$, $\lifttwoC{DeepGood}$, $\lifttwoC{Ring}$,
$\lifttwoC{Maps}$, $\lifttwoC{Rk}$, $\lifttwoC{Gmin}$, $\lifttwoC{Flat}$ and $\lifttwoC{ULow}$ imply $\lifttwoC{Low}$.
\end{enumerate}
\formallabel{GlobalTheory.CovTop.LowCover.lowAt_of_cover}
\formallabel{GlobalTheory.CovTop.LowCover.low_of_cover_hypotheses}\formallabel{GlobalTheory.T.liftLowBound_of_cover_estimates}
\end{proposition}

\begin{proof}
(a) is \cref{lem:covolume:low}, whose proof applies the descent inequality \cref{pf:lift2:partI}(d) to
$A=\Lambda_Y(G_s(-c_0))$ (nonzero by $\lifttwoP{DeepGoodAt}$), $B=\Lambda^{\mathrm{deep}}_Y$, the maps of $\lifttwoP{MapsAt}$
(with $P=b$), the generators of $\lifttwoP{GminAt}$ (with $\Lambda=P\kappa$), $\eps=b/(P\kappa)$ and the rank bound
$n\lifttwork B\le\lifttwork_Y(G_s(c_0))\le(1+\eps)\lifttwork A$ from $\lifttwoP{MapsAt}$ and $\lifttwoP{RkAt}$. (b) Let
$C_D$, $C_F$, $C_U$ be the constants of $\lifttwoC{LowDef}$, $\lifttwoC{Flat}$, $\lifttwoC{ULow}$ and take
$C=\max(C_D,2(C_F+C_U)+9)$. At an admissible point, if $j=m$ then $\lifttwoP{LowAt}(C_D)$ holds by $\lifttwoC{LowDef}$; if
$j<m$, $\lifttwoC{DeepEx}$ provides a good deep cover $Y$, and (a) gives $\lifttwoP{LowAt}(2(C_F+C_U)+9)$. Since $U\ge0$
(from $\lifttwoP{SizeAt}$), $\lifttwoP{LowAt}(C')$ implies $\lifttwoP{LowAt}(C)$ for $C\ge C'$.
\end{proof}

The condition $\lifttwoC{LowDef}$ (the bound at the top index $j=m$, where no deep cover is used) is proved in the
companion section: \lref{GlobalTheory.T.lowDef_of_gram_of_anchor} states that \textup{(Hg)}, $\lifttwoC{Gram}$ and
$\lifttwoC{Anchor}$ imply $\lifttwoC{LowDef}$.

% ---------------------------------------------------------------------------
\subsection{The bad part}\label{sec:pf:lift2:bad}

\begin{proposition}[Rank of the bad part]\label{pf:lift2:bad}
Assume \textup{(X)}, \textup{(Hg)}, \textup{(GM)}, \textup{(Q)}, \textup{(Sg)}, \textup{(KS)} and \textup{(Lie)}. Then $\lifttwoC{Gram}$, $\lifttwoC{DeepEx}$, $\lifttwoC{DeepGood}$,
$\lifttwoC{Size}$, $\lifttwoC{FlatDeg}$, $\lifttwoC{Deg}$ and $\lifttwoC{Share}$ imply $\lifttwoC{Bad}$, with $C=0$.
\formallabel{GlobalTheory.CovTop.covering_bad_locus_of_descent_steps_alt}
\formallabel{GlobalTheory.T.liftBadSet_of_deep_size_flatDeg_share}
\end{proposition}

\begin{proof}
The development proves two closed statements:
\begin{enumerate}[label=(\roman*)]
\item at the top index: under \textup{(Hg)} and Green's theorem on cuspidal characters (part of \textup{(X)}), $\lifttwoC{Gram}$
implies the lift condition of ``$j=m\Rightarrow\lifttwoP{BadAt}$''
(\lref{GlobalTheory.T.lifts_badAt_last_of_green_cuspidal}, deferred);
\item at the lower indices: under the hypotheses of the proposition, the cover condition of $\lifttwoP{BadAt}$
(\lref{GlobalTheory.T.covers_badAt_of_deep_good_size_flat_degree}, deferred).
\end{enumerate}
Both proofs follow the same pattern. A finite group $\Gamma$ acting on a finite-dimensional complex representation
$V$ is attached to the point (at the top index; for lower indices the level group of $Y$ acting on the complexified
cusp forms of weight $G_k$), together with a subgroup $P_b$ with a character $\lambda_b$ (a Whittaker datum) and two
subgroups $H_0,H_1$ (one for each auxiliary prime $g_0,g_1$), such that $\lifttwork\lifttwoLat{type}$ is the dimension
of the $\lambda_b$-eigenspace of $P_b$ and $\lifttwork\lifttwoLat{bad}$ is the dimension of the part of that eigenspace
killed by an idempotent attached to the types. At the top index, $H_k$ acts freely on the index set of an internal direct
sum decomposition of $V$ permuted by $\Gamma$ (\lref{GlobalTheory.T.deformation_class_has_representative_alt}); at
lower indices, the corresponding regularity of embeddings, with multiplicities $n_k$ satisfying
$n_k|H_k|\le\frac32|P_b|\dim V^{\lambda_b}$, is proved using the existence of a cusp form whose level determinant is a
non-zero-divisor (\lref{GlobalTheory.T.orbits_separated_by_regular_cusp}, which states that for the level group of
a type frame, and for each of its two halves, some cusp form $u$ of weight $\underline{c_0h}$ has
$\det((hg)^*u)_{h,g}$ a non-zero-divisor in the ring of sections). In both cases this bounds the bad part by the sum,
over the two auxiliary primes, of the local shares of the non-good tuples, times $c\dim V^{\lambda_b}$ with $c=1$,
resp.\ $c=\frac32$. Finally there are nontrivial additive characters of $\Z/g_k$ for which this sum of local shares is at
most $\exp(-x^{e_z/7})$, for every choice of non-good tuples
(\lref{GlobalTheory.CovTop.LocShareNum.local_share_count_le_exp_alt}, deferred; this is where Green's theorem on
cuspidal characters is used). Hence $\lifttwork\lifttwoLat{bad}\le c\,\delta_x\lifttwork\lifttwoLat{type}\le3\delta_x
\lifttwork\lifttwoLat{type}$.

Given (i) and (ii), let $p$ be admissible. If $j=m$, (i) applies. If $j<m$, $\lifttwoC{DeepEx}$ provides a good deep
cover, and (ii) applies to it. Combine by \cref{pf:lift2:combine}(c).
\end{proof}

% ---------------------------------------------------------------------------
\subsection{Proof of the lift estimates}\label{sec:pf:lift2:main}

\begin{lemma}\label{pf:lift2:loss}
For $c>0$ and $A\in\R$, $e^{-x^c}x^A\to0$ as $x\to\infty$. Consequently, for $c>0$ and $A,B,C\in\R$ there is $x_0$
with $12e^{-x^c}(x^A+x^B+C)\le1$ for all $x\ge x_0$.
\formallabel{GlobalTheory.CovTop.Assembly.tendsto_exp_mul_rpow}
\formallabel{GlobalTheory.CovTop.Assembly.exists_loss_le}
\end{lemma}

\begin{proof}
With $y=x^c\to\infty$, $e^{-x^c}x^A=e^{-y}y^{A/c}\to0$. Hence $12e^{-x^c}(x^A+x^B+Cx^0)\to0<1$.
\end{proof}

\begin{theorem}[Final assembly; \cref{thm:covolume:assembly}]\label{pf:lift2:assembly}
The conditions $\lifttwoC{Gram}$, $\lifttwoC{Rec}$, $\lifttwoC{Cut}$, $\lifttwoC{Bad}$, $\lifttwoC{Lam}$,
$\lifttwoC{Size}$, $\lifttwoC{Core}$ and $\lifttwoC{Low}$ imply $\sfname{CovolumeBoundOfX}$. More precisely, if at
$(p,i)$ the predicates $\lifttwoP{GramAt}$, $\lifttwoP{RecAt}$, $\lifttwoP{CutAt}$, $\lifttwoP{BadAt}$,
$\lifttwoP{LamAt}(A)$, $\lifttwoP{SizeAt}(P,B)$, $\lifttwoP{CoreAt}(C_1)$ and $\lifttwoP{LowAt}(C_2)$ hold, with
$C_1,C_2\ge0$ and $12\exp(-x^{e_z/7})(x^A+x^B+C_1+C_2+1)\le1$, then
$\lifttwolcv(X_{W,i},M_i)/D\le\lifttwomainT+(C_1+C_2+1)U$.
\formallabel{GlobalTheory.CovTop.Assembly.lcv_le}
\formallabel{GlobalTheory.CovTop.Assembly.covolX_of_lift_bounds_and_coreLe}
\formallabel{GlobalTheory.T.covolX_of_truncation_badSet_lambda_lowBound}
\end{theorem}

\begin{proof}
This is \cref{thm:covolume:assembly}; its proof in Part~I applies \cref{pf:lift2:partI}(c) to
$L=\lifttwoLat{type}\subseteq M=\lifttwoLat{all}$, $e=e_i$, $g=|G_p|$, $\delta=\delta_x$ and
$Q=(x^A+x^B+C_1+C_2+1)U$, and then $\lifttwoP{RecAt}$, $\lifttwoP{CoreAt}$ and $\lifttwoP{LowAt}$. The threshold $x_0'$ is
given by \cref{pf:lift2:loss} with $c=e_z/7$ ($e_z>0$ is the common exponent of the family), and the
covolume bound holds for all $C\ge C_0=C_1+C_2+1$ since $U\ge1$.
\end{proof}

\begin{theorem}[Lift estimates; \cref{thm:covolume:lift}]\label{pf:lift2:main}
Assume the inputs \textup{(VE)}, \textup{(X)}, \textup{(Hg)}, \textup{(H)}, \textup{(GM)}, \textup{(Q)} and \textup{(DS)} of \cref{sec:pf:lift2}. Then all conditions of
\cref{pf:lift2:def:conds} hold, and $\sfname{CovolumeBoundOfX}$ holds.
\formallabel{SelfIntersectionVolumeEstimates.topCovolG_alt_alt}
\end{theorem}

\begin{proof}
Let \textup{(Sg)} be the Siegel toroidal input contained in \textup{(H)}, and \textup{(KS)}, \textup{(Lie)} the unitary toroidal input and the Lie
seminegativity contained in \textup{(X)}. We establish the conditions in the following order.
\begin{enumerate}[label=(\arabic*)]
\item $\lifttwoC{Gram}$ (\cref{pf:lift2:gram}), $\lifttwoC{Size}$ (\cref{pf:lift2:size}), $\lifttwoC{Meas}$
(\cref{pf:lift2:meas}, from \textup{(Hg)}), $\lifttwoC{Anchor}$ (\cref{pf:lift2:anchor}).
\item $\lifttwoC{DeepEx}$ (\cref{pf:lift2:deepex}, from \textup{(Hg)}); $\lifttwoC{DeepGood}$ (\cref{pf:lift2:deepgood}, from
\textup{(Hg)}, \textup{(GM)}, \textup{(Q)}, \textup{(Sg)}, \textup{(KS)}, \textup{(Lie)}); $\lifttwoC{FlatDeg}$ (\cref{pf:lift2:flatdeg}, from \textup{(Hg)}).
\item $\lifttwoC{Deg}$ (\cref{pf:lift2:deg}) and $\lifttwoC{Share}$ (\cref{pf:lift2:share}), from \textup{(X)}, \textup{(Hg)}, \textup{(Sg)}, \textup{(KS)},
\textup{(GM)}, \textup{(Q)}, \textup{(Lie)} and the conditions of (1), (2).
\item $\lifttwoC{Height}$ by \lref{SelfIntersectionVolumeEstimates.height_lower_bound_general_alt}, from \textup{(VE)}, \textup{(X)},
\textup{(GM)}, \textup{(Q)}, \textup{(Hg)}, \textup{(H)}, \textup{(DS)}, $\lifttwoC{DeepGood}$, $\lifttwoC{Size}$, $\lifttwoC{Deg}$, $\lifttwoC{FlatDeg}$; and
$\lifttwoC{HS}$ by \lref{GlobalTheory.T.coverHeightSum_of_flatDeg_coverDegree_share}, from \textup{(X)}, \textup{(Hg)}, \textup{(GM)}, \textup{(Q)},
\textup{(Sg)}, \textup{(KS)}, \textup{(Lie)}, $\lifttwoC{DeepGood}$, $\lifttwoC{Size}$, $\lifttwoC{FlatDeg}$, $\lifttwoC{Deg}$, $\lifttwoC{Share}$.
\item $\lifttwoC{Sum}$ (\cref{pf:lift2:sum}) and $\lifttwoC{ULow}$ (\cref{pf:lift2:ulow}).
\item $\lifttwoC{LowDef}$ by \lref{GlobalTheory.T.lowDef_of_gram_of_anchor}, from \textup{(Hg)}, $\lifttwoC{Gram}$,
$\lifttwoC{Anchor}$; $\lifttwoC{Ring}$ (\cref{pf:lift2:ring}), $\lifttwoC{Maps}$ (\cref{pf:lift2:maps}),
$\lifttwoC{Rk}$ (\cref{pf:lift2:rank}), $\lifttwoC{Gmin}$ (\cref{pf:lift2:gmin}), $\lifttwoC{Flat}$ (\cref{pf:lift2:flat});
then $\lifttwoC{Low}$ (\cref{pf:lift2:low}).
\item $\lifttwoC{Rec}$ (\cref{pf:lift2:rec}), $\lifttwoC{Cut}$ (\cref{pf:lift2:cut}), $\lifttwoC{Bad}$
(\cref{pf:lift2:bad}, which uses Green's theorem from \textup{(X)}), $\lifttwoC{Lam}$ (\cref{pf:lift2:lam}), $\lifttwoC{Rat}$
(\cref{pf:lift2:rat}) and $\lifttwoC{Core}$ (\cref{pf:lift2:core}).
\end{enumerate}
Finally \cref{pf:lift2:assembly} gives $\sfname{CovolumeBoundOfX}$.
\end{proof}

\begin{remark}[On the hypotheses]\label{pf:lift2:hyprem}
(1) The inputs \textup{(VE)} and \textup{(DS)} and the remaining parts of \textup{(H)} are used only through the height lower bound of the companion
section, i.e.\ in $\lifttwoC{Height}$. (2) Green's theorem on cuspidal characters is used only for $\lifttwoC{Bad}$.
(3) The formal statements of \cref{pf:lift2:gram} and of the second part of \cref{pf:lift2:ring} assume the Hodge
package \textup{(Hg)}, which their proofs do not use. The Gram condition, the rationality and the flat degree condition
hold at all admissible points, with no threshold. (4) The rank comparison is
proved in a stronger form (with the factor $1+\min(\frac12,b/(P\kappa))$ and also for $2c_0$), of which only the case
$s=1$ is used. (5) In the sum bound \cref{pf:lift2:sum}, when the lattice has rank $0$ the inequality is
$0\le\mu_Y(u)$, which holds because the slope of the zero module is $-\log2$ plus the logarithmic root discriminant of
$F_Y$.
\end{remark}

% ---------------------------------------------------------------------------
\subsection{Proofs of auxiliary results}\label{sec:pf:lift2:aux}

This subsection proves results which are used above and in the companion section on lift estimates. We keep the
notation of \cref{sec:pf:lift2:objects}; deeper ingredients are cited with a statement of what is used.

\subsubsection*{Weight parameters and numerical bounds}

\begin{definition}[Weight parameters]\label{pf:lift2:def:wtpar}
Let $j$ be a level of the core datum, $N=p_j+q_j$, $d=\lifttwodF$, and for a weight $v$ write
$\mathrm{gap}(v_\tau)_l=v_{\tau,l}-v_{\tau,l+1}$ ($l+1<N$). Let $\kappa_{\mathrm{Fl}}$ be the weight with
$\kappa_{\mathrm{Fl};\tau,l}=q_j-p_j+2l+1$ for $l<p_j$ and $\kappa_{\mathrm{Fl};\tau,l}=2l-3p_j-q_j+1$ for $l\ge p_j$ (the
canonical weight of the flag variety). A \emph{weight parameter} is a quadruple $P=(u,c,a,T)$ of a weight $u$, an
integer $c$, an index $a<N$ and $T\in\N$. Put $T_1=T+d+1$ and $u^-=u-c\,h$, and let $u_1=u_1(P)$ and $\varrho=\varrho(P)$
be the weights which vanish at the index $a$ and whose gaps are the quotient and the remainder of the Euclidean
division by $T_1$ of the gaps of $u^--\kappa_{\mathrm{Fl}}$:
\[
\mathrm{gap}((u^--\kappa_{\mathrm{Fl}})_\tau)_l=T_1\,\mathrm{gap}(u_{1,\tau})_l+\mathrm{gap}(\varrho_\tau)_l,\qquad
0\le\mathrm{gap}(\varrho_\tau)_l<T_1 .
\]
Put $u_2=u_2(P)=u_1+\varrho$ and $\mathrm{wt}_t(P)=\kappa_{\mathrm{Fl}}+u_2+(d+t)\,u_1$ for $t\in\N$. The parameter $P$ is
\emph{large} if every consecutive gap of $u_1(P)$ is at least
$G(d,N)=(d+3)\bigl(\binom{d+1}2+4N+d+3\bigr)$, $d+1\le T$ and $d\ge1$. The \emph{canonical weight parameter} of $(p,i)$
is $P(p,i)=(u_W,c_0,k-2,T)$, with $T=\lceil P^2d^2k/m\rceil$ as in \cref{pf:lift2:def:wt}. For a weight $u$ we write
$u'=u-3Nh$, so that $\rho_B(u)=1-\lifttwovdm(u')/\lifttwovdm(u)$.
\formal{GlobalTheory.CompHS.WtPar}
\formallabel{GlobalTheory.FlagTan.kap}
\formallabel{GlobalTheory.HSApp.WtBig}
\formallabel{GlobalTheory.CompHt.uPrime}
\formallabel{GlobalTheory.CompHt.rhoB}
\end{definition}

Since $\mathrm{gap}(h_\tau)_l=1$ if $l+1=p_j$ and $0$ otherwise, and $\mathrm{gap}(\kappa_{\mathrm{Fl};\tau})_l=2N-2$ if
$l+1=p_j$ and $-2$ otherwise, we have
\begin{equation}\label{eq:pf:lift2:gapminus}
\mathrm{gap}((u^--\kappa_{\mathrm{Fl}})_\tau)_l=\mathrm{gap}(u_\tau)_l+\epsilon_l(c),\qquad
\epsilon_l(c)=\begin{cases}-c-2N+2,&l+1=p_j,\\2,&\text{otherwise}.\end{cases}
\end{equation}
In particular the quotient weights of \cref{pf:lift2:def:wt} are exactly the weights which vanish at $k-2$ and have the
gaps of $u_1(P(p,i))$; as a weight is determined by its value at one index and its gaps, $u_1(P(p,i))$ is the only
quotient weight. Moreover, by construction,
\begin{equation}\label{eq:pf:lift2:uminus}
u^-_{\tau,l}=\mathrm{wt}_T(P)_{\tau,l}+\bigl(u^-_{\tau,a}-\kappa_{\mathrm{Fl};\tau,a}\bigr)
\end{equation}
for all $\tau,l$, since $u^--\kappa_{\mathrm{Fl}}-(u^-_{\tau,a}-\kappa_{\mathrm{Fl};\tau,a})=T_1u_1+\varrho$ and
$\mathrm{wt}_T(P)=\kappa_{\mathrm{Fl}}+\varrho+T_1u_1$.

\begin{lemma}\label{pf:lift2:expl}
For $L\in\R$ let $e(L)=\lfloor3L/\log3\rfloor\in\N$ (so that the level of a good deep cover has prime-to-level part
$15^{e(\ell)}$). If $e\in\N$ and $e\log3\le3L$, then $e\le e(L)$. If $L\ge10$, then $e(L)\ge4$.
\formallabel{GlobalTheory.CovTop.expL}
\formallabel{GlobalTheory.CovTop.le_expL}
\formallabel{GlobalTheory.ThetaSteep.four_le_expL}
\end{lemma}

\begin{proof}
The first claim is the definition of the integer part, as $\log3>0$. For the second, $\log3\le3-1=2$, so
$4\log3\le8\le3L$.
\end{proof}

\begin{proposition}[Large gaps]\label{pf:lift2:wtbig}
The lift condition of the statement ``if $j<m$, the canonical weight parameter $P(p,i)$ is large'' holds, with the
constant $0$.
\formallabel{GlobalTheory.HSApp.wtBig_run_alt}
\end{proposition}

\begin{proof}
The parameter data of the family provide two gap sizes $g_S$ and $h_0$ at every $(x,\iota)$ (an expression in the
ranks and a constant of the family, resp.\ a size parameter of $L$ divided by $m$). We use:
\begin{enumerate}[label=(\roman*)]
\item (\lref{GlobalTheory.T.runSizes_quotientGaps_large}) the closed statement: eventually, for all $\iota$,
$d=f\binom r2\ge1$, $d+1\le T$, and with $T_1=T+d+1$ and $G=G(d,r)$,
\[
G\,T_1+4\le\min(g_S,h_0),\qquad G\,T_1+22f^2r^4+2r+1\le h_0 ;
\]
\item (\lref{GlobalTheory.T.anchoredWeight_indefiniteLift_gaps}) at every point and for $j<m$: $N=r$,
$\lifttwodF=f\binom r2$, and for every $\tau$ and $l$ with $l+1<r$, $\mathrm{gap}(u_{W,\tau})_l\ge\min(g_S,h_0)-2$ if
$l+1\ne p_j$, and $\mathrm{gap}(u_{W,\tau})_l\ge h_0+r-3$ if $l+1=p_j$.
\end{enumerate}
Let $p$ be admissible with $x$ beyond the threshold of (i), and $j<m$. By (ii), $d=\lifttwodF$ and $T$ are the numbers
of (i). By \eqref{eq:pf:lift2:gapminus} with $c=c_0$, the gap of $u_1$ at $l$ is the integer part of
$(\mathrm{gap}(u_{W,\tau})_l+\epsilon_l(c_0))/T_1$, and it is $\ge G$ as soon as
$\mathrm{gap}(u_{W,\tau})_l+\epsilon_l(c_0)\ge GT_1$. If $l+1\ne p_j$, then
$\mathrm{gap}(u_{W,\tau})_l+2\ge\min(g_S,h_0)\ge GT_1+4$. If $l+1=p_j$, then, as $c_0\le22f^2r^4$,
\[
\mathrm{gap}(u_{W,\tau})_l-c_0-2r+2\ge h_0-r-1-c_0\ge GT_1+22f^2r^4+r-c_0\ge GT_1 .
\]
Together with $d+1\le T$ and $d\ge1$ from (i), $P(p,i)$ is large.
\end{proof}

\begin{proposition}[Weights at a point]\label{pf:lift2:wtpt}
The lift condition of the following statement holds, with the constant $0$: if $j<m$, then
\begin{enumerate}[label=(\alph*)]
\item $p_j\ge1$, $q_j\ge1$ and $[E:\Q]/2>2$;
\item every quotient weight equals $u_1=u_1(P(p,i))$;
\item for every $\tau$, $u_{1,\tau}$ and $u'_{1,\tau}$ are strictly decreasing;
\item there is $w\colon\Psi\to\Z$ with $|w_\tau|\le3r$ for all $\tau$ and $u_M=\mathrm{wt}_T(P(p,i))+w$ (the function $w$
regarded as a weight constant in the second variable);
\item for $s\in\{0,1,2,3\}$ and every $\tau$, the sequences $u_{M,\tau}+sc_0h_\tau$ and
$(T+\lifttwodF+3+[s\ne0])\,u_{1,\tau}-(u_{M,\tau}+sc_0h_\tau)$ are non-increasing.
\end{enumerate}
In the development this lift condition is stated under the hypothesis $\lifttwoC{Size}$, which is not used.
\formallabel{GlobalTheory.T.weight_parameter_strict_antitone_alt}
\formallabel{GlobalTheory.QuoWt.weights_at_point_alt}
\formallabel{GlobalTheory.QuoWt.atPoint}
\formallabel{GlobalTheory.QuoWt.dominance}
\formallabel{GlobalTheory.QuoWt.arith}
\formallabel{GlobalTheory.QuoWt.quo_unique}
\formallabel{GlobalTheory.QuoWt.isQuot_u₁}
\formallabel{GlobalTheory.MeanTop.divided_strictAnti}
\end{proposition}

\begin{proof}
By \cref{pf:lift2:numholds} it suffices to prove the statement at a point where $\mathcal N(P)$ and $\mathcal H(P,A,1)$
hold; let $j<m$ and let $h_0$ be the integer provided by $\mathcal N(P)$. Write $N=r$, $d=\lifttwodF$, $T_1=T+d+1$ and
$g_l=\mathrm{gap}(u_{W,\tau})_l$. The hypotheses of the divided-weight construction contained in $\mathcal N(P)$
include: $4d+4\le T$; $4(c_0+2N+T+d+1)\le h_0$ and $24NT\le h_0$; $(T+d+2)(T+d+2+2N)\le g_l+2$ for $l+1\ne p_j$; and
$h_0\le g_l$ for $l+1=p_j$.

(a) $p_j=k+j-1\ge1$, and $q_j=m-j\ge1$ as $j<m$; $[E:\Q]/2=f\ge8$.

(b) By \eqref{eq:pf:lift2:gapminus} and the remark following it.

(c) If $l+1\ne p_j$, the gap of $u_1$ at $l$ is $\lfloor(g_l+2)/T_1\rfloor\ge\lfloor(T_1+1)(T_1+1+2N)/T_1\rfloor\ge1$,
and $u'_1$ has the same gap there. If $l+1=p_j$, then $g_l-c_0-2N+2\ge\frac34h_0\ge18NT\ge(3N+1)\frac54T\ge(3N+1)T_1$
(using $c_0+2N\le h_0/4$ and $T_1\le\frac54T$), so the gap of $u_1$ at $l$ is at least $3N+1$ and that of $u'_1$, which
is smaller by $3N$, is at least $1$.

(d) Here $P(p,i)$ has $u=u_W$, $c=c_0$, so $u^-=u_W-c_0h=u_M$. At the anchor $a=k-2<p_j$ we have $u_{W,\tau,a}=0$ and
$h_{\tau,a}=0$, so \eqref{eq:pf:lift2:uminus} gives $u_M=\mathrm{wt}_T(P(p,i))+w$ with
$w_\tau=-\kappa_{\mathrm{Fl};\tau,k-2}=-(m+k-2j-2)$, and $|w_\tau|\le m+k\le3r$.

(e) First, an arithmetic consequence of $\mathcal N(P)$ and $\mathcal H(P,A,1)$: since $z\ge f r\ell\ge480$ and
$z^{12}\le P$ we have $P\ge40$; moreover $4d+4\le T$, $4T^2\le h_0$, $PTNf^2m^2kd\le h_0$, $4(c_0+2N+T+d+1)\le h_0$,
$0\le c_0\le22f^2N^4$, $d=fN(N-1)/2$, $N+4\le2k$, $m\ge2$, $N\ge6$ and $f\ge8$. These give
\begin{equation}\label{eq:pf:lift2:arith}
(T+c_0)T_1+c_0+2N-2\le h_0,\qquad 2N+d\le T .
\end{equation}
Indeed $2N\le3d+4$ as $d\ge4N(N-1)\ge2N$; and $T^2\le h_0/4$, $T(d+1)\le T^2$, $c_0+2N-2\le h_0/4$, while
$PTNf^2m^2kd\ge40\cdot4\cdot T\,N f^2\,\frac{N+4}2\,\frac{fN(N-1)}2\ge176f^2N^4T$ gives $c_0T_1\le44f^2N^4T\le h_0/4$.
Now let $s\le3$. The gaps of $u_{M,\tau}+sc_0h_\tau$ are $g_l\ge0$ for $l+1\ne p_j$ and $g_l+(s-1)c_0\ge h_0-c_0\ge0$ for
$l+1=p_j$. Let $M_s=T_1+2$ if $s=0$ and $M_s=T_1+3$ otherwise, and let $q_l=\mathrm{gap}(u_{1,\tau})_l$. The gap of
$M_su_{1,\tau}-(u_{M,\tau}+sc_0h_\tau)$ at $l$ is $M_sq_l-g_l$ if $l+1\ne p_j$ and $M_sq_l-g_l+(1-s)c_0$ if $l+1=p_j$. In
the first case $g_l\le(q_l+1)T_1-3$ by definition of $q_l$, so the gap is $\ge2q_l-T_1+3\ge0$, because
$q_l\ge T_1+1$. In the second case $g_l\le(q_l+1)T_1+c_0+2N-3$ and, by \eqref{eq:pf:lift2:arith}, $q_l\ge T+c_0$; so for
$s=0$ the gap is $\ge2q_l-T_1-2N+3\ge T+2c_0-d-2N+2\ge0$, and for $1\le s\le3$ it is
$\ge3q_l-T_1-2N+3-sc_0\ge2T-d-2N+2\ge0$, by \eqref{eq:pf:lift2:arith}.
\end{proof}

\begin{definition}[Theta exponents]\label{pf:lift2:def:theta}
For the datum at level $j$ let $\Phi$ be the set of pairs $(\tau,l)$ with $\tau\colon E\to\C$ and $0\le l<\mathrm{mult}(\tau)$
(the multiplicity of $\tau$ is $p_j$ for $\tau\in\Psi$ and $q_j$ otherwise), and
$n_c=8\bigl(\lfloor\frac18\sum_{(\tau,l)\in\Phi}2(l+1)\rfloor+|\Phi|+1\bigr)$. For a weight $u$ and $m'\in\Z$ let
$\mathrm{wt}_{m'}(u)$ be the weight of automorphic forms with parity $2m'$ and exponents $u_{\tau,l}-m'$ at $(\tau,l)$
for $\tau\in\Psi$, $l<p_j$, $-u_{\bar\tau,N-1-l}-m'$ at $(\tau,l)$ for $\bar\tau\in\Psi$, $l<q_j$, and $0$ otherwise
(so $\mathrm{wt}_0(u)=\underline u$). Put $\sigma_u(\tau)=n_c/2+1-u_{\tau,p_j-1}$, and let the \emph{theta exponents}
$a(u)$ be the exponents of $\mathrm{wt}_{n_c/2}(u+\sigma_u)$ with negative values replaced by $0$; the \emph{theta
weight} is $G(u)=(a(u),n_c)$. A pair $(a,n)$ of exponents and an integer is \emph{dominant} if $a(\tau,\cdot)$ is
non-increasing for every $\tau$ and $n$ is even; $a$ is \emph{positive at the last entries} if $a(\tau,p_j-1)>0$ for
$\tau\in\Psi$; and $d_F(a)=\sum_{(\tau,l)\in\Phi}a(\tau,l)$. Finally, the \emph{jet numerical condition}
$\mathrm{JN}(d_F)$ means: for $k=0,1$, $4d_F+4\le g_k$, $2d_F\sqrt{2fm}<g_k$ and
$r_0(2fm)^3\bigl(r\cdot2fm/4+2d_F^2\bigr)\le g_k/2500$, where $r_0=r_0(D)$ is a constant of the global data.
\formallabel{GlobalTheory.RefSum.nCol}
\formallabel{GlobalTheory.RefSum.shift}
\formallabel{GlobalTheory.RefSum.thA}
\formallabel{GlobalTheory.RefSum.aOf}
\formallabel{GlobalTheory.RefSum.GOf}
\formallabel{GlobalTheory.RefSum.JetNum}
\formal{Formalism.L2.Pj.Dom}
\formal{Formalism.L2.Pj.LastPos}
\formal{Formalism.L2.Pj.DF}
\end{definition}

\begin{proposition}[Shape of the quotient weight]\label{pf:lift2:quoshape}
The lift condition of the following statement holds: if $j<m$, then for every quotient weight $u_1$: $p_j,q_j\ge1$,
$p_j+q_j=r$; $u_{1,\tau}$ is non-increasing for every $\tau$; $\mathrm{gap}(u_{1,\tau})_{p_j-1}\ge n_c+1$;
$|u_{1;\tau,l}|\le2b$; $(a(u_1),n_c)$ is dominant; $a(u_1)$ is positive at the last entries; and
$G(u_1)=\mathrm{wt}_{n_c/2}(u_1+\sigma_{u_1})$. The development states it under the hypothesis $\lifttwoC{Size}$, which
is not used.
\formallabel{GlobalTheory.T.quotient_weight_gap_bounds'}
\formallabel{GlobalTheory.RefSum.QuoShape.quo_shape_dom_lastPos_weight_eq}
\formallabel{GlobalTheory.RefSum.QuoShape.shape_at}
\end{proposition}

\begin{proof}
By \cref{pf:lift2:numholds}(b) it suffices to argue at a point where $\mathcal N(P)$ holds, with $j<m$; let $h_0$ be as
there, $N=r$, $d=\lifttwodF$ and $T_1=T+d+1$. By \cref{pf:lift2:wtpt}(b),(c), proved from the same hypotheses, every
quotient weight equals $u_1=u_1(P(p,i))$, which is strictly decreasing; and $p_j,q_j\ge1$, $p_j+q_j=r$. The laws of
the core datum give $n_c\le20fr^2$ (\lref{GlobalTheory.RefSum.SizeTZ.nCol_le_of_pins}: the number of factors of a
datum with the pinned signatures), and $rf^2m^2kd\ge80fr^2$ since $d=f\binom r2$, $m,k\ge2$ and $r\ge6$; hence, with
$P\ge1$, $4Tn_c\le PTrf^2m^2kd\le h_0$. Let $q$ be the gap of $u_{1,\tau}$ at $p_j-1$. As in the proof of
\cref{pf:lift2:wtpt}(c), $q\ge\frac34h_0/T_1-1\ge\frac{3h_0}{5T}-1$, so $4T(q-3N)\ge\frac{12}5h_0-4T-12NT\ge h_0$, because
$h_0\ge24NT$. Therefore $4Tn_c\le4T(q-3N)$ and $q\ge n_c+3N\ge n_c+1$. For the entries, the divided weight satisfies
$T_1|u_{1;\tau,l}|\le|u_{W;\tau,l}|+c_0+4N$ (\lref{GlobalTheory.MeanTop.divided_entries}: an estimate for the
entries of a divided weight under the hypotheses of the construction), and $|u_{W;\tau,l}|+c_0+4N\le8b$ is one of
these hypotheses; as $T_1\ge4$, $|u_{1;\tau,l}|\le2b$. Finally $n_c$ is even (it is divisible by $8$), so the
statement on theta exponents (\lref{GlobalTheory.T.thetaExponents_dominant_lastEntry_one_degree}: for $p_j\ge1$,
even $n$, and $u$ non-increasing with gap $\ge n+1$ at $p_j-1$ and entries bounded by $B$, the pair $(a(u),n)$ is
dominant, $a(u)$ is positive at the last entries, $(a(u),n)$ is the weight $\mathrm{wt}_{n/2}(u+\sigma_u)$, the entries of
$a(u)$ are at most $2B+1$, and $d_F(a(u))\le\frac{[E:\Q]}2N(2B+1)$) gives the remaining claims.
\end{proof}

\begin{proposition}[Size relations at a point]\label{pf:lift2:sizerel}
Assume $\lifttwoC{Size}$. Then the lift condition of the following statement holds: $Tz^{12}\le U$, $\ell\ge1$,
$\log\mathrm{rd}_E\le3\ell$, $n_c\le20fr^2$, $f\ge1$ and $r\ge1$.
\formallabel{GlobalTheory.T.lenT_zBase_unitE_nCol_bounds}
\formallabel{GlobalTheory.RefSum.SizeTZ.parameters_scale_rank_gap_bounds}
\end{proposition}

\begin{proof}
At every point, $Tz^{12}\le b$ and $n_c\le20fr^2$ (\lref{GlobalTheory.RefSum.SizeTZ.lenT_mul_le_b}: the parameter
data bound $\lceil P^2d^2k/m\rceil z^{12}$ by $b$, and the laws bound the factor count); $b\le U$ by
$\lifttwoP{SizeAt}$; $\ell\ge1$ by the scale axioms; $\log\mathrm{rd}_E\le3\ell$ by the parameter choice
(\lref{GlobalTheory.RefSum.SizeTZ.lrd_le}); and $f\ge2$, $r\ge6$ by the axioms of the ranks. The constant of
$\lifttwoC{Size}$ can be kept.
\end{proof}

\begin{proposition}[Numbers at a point]\label{pf:lift2:numpt}
Assume $\lifttwoC{Size}$. Then the lift condition of the following statement holds: if $j<m$, then for every
quotient weight $u_1$, $\mathrm{JN}(d_F(a(u_1)))$ holds, $2r\le n_c$, and $T(\log g_0+\log g_1)\le CU$.
\formallabel{GlobalTheory.T.jetNum_nCol_lenT_bounds_of_quo}
\formallabel{GlobalTheory.RefSum.NumPtAlt.numbers_at_point}
\end{proposition}

\begin{proof}
Combine \cref{pf:lift2:quoshape} and \cref{pf:lift2:sizerel} by \cref{pf:lift2:combine}(c), with
$C=\max(0,2+2C_Y)$, where $C_Y>0$ is a constant of the parameter choice of the family. Let $j<m$ and $u_1$ a quotient
weight. By the statement on theta exponents quoted in the proof of \cref{pf:lift2:quoshape} (with $B=2b$,
$[E:\Q]/2=f$ and $N=r$), $d_F(a(u_1))\le fr(4b+1)$. The size data of the family give $b\ge t^7$, $b\ge\mathrm{rd}_E^2$,
$r_0\le b$ and $g_k\ge b^8$ for $k=0,1$, where $t$ is the largest rank of $D$; with $d_F\le fr(4b+1)$ and
$f,m,r\le t$ the three inequalities of $\mathrm{JN}(d_F)$ follow by an elementary computation. The inequality
$2r\le n_c$ holds for every datum, since $n_c\ge8|\Phi|\ge2N$. Finally the parameter choice gives
$\log g_k\le(1+C_Y)\ell$ for $k=0,1$, so $\log g_0+\log g_1\le(2+2C_Y)\ell\le(2+2C_Y)z^{12}$, and $Tz^{12}\le U$ by
\cref{pf:lift2:sizerel}; multiply.
\end{proof}

\begin{proposition}[Size bounds on deep covers]\label{pf:lift2:sizecov}
Assume $\lifttwoC{Size}$. Then the cover condition of the following statement holds (with $u_1,u_2,\mathrm{wt}_t$ those
of $P(p,i)$ and $|v|_1=\sum_{\tau,l}|v_{\tau,l}|$): the core datum satisfies its laws, $j$ is in range, $\ell\ge1$,
$\log M'_Y\le30\ell$, $\mathrm{fs}(E)\le3\ell$, $\log(fr)\le\ell$, $e(\ell)\ge4$, $g_k>r+1$ for $k=0,1$, $f\ge8$,
$r\ge100$, the value of the flat degree at the point is at most $9r\ell$ in absolute value, $\log\kappa\le2\ell$,
$z^{12}\le P$, $TP\le\sqrt b$, $\lifttwozed(u_1)P\kappa\le b$, $\lifttwozed(u_1)\le\frac12$, $|u_1|_1,|u_2|_1\le
z^2(\kappa/T+1)$, and $|\mathrm{wt}_t|_1\le z^2\kappa$ for $t\le T$.
\formallabel{GlobalTheory.T.size_bounds_for_family_alt}
\formallabel{GlobalTheory.CovTop.CovSz.size_rows_alt}
\end{proposition}

\begin{proof}
The statement is the conjunction of six cover conditions, each proved in the development from $\lifttwoC{Size}$ and
the parameter data of the family: (1) the structural bounds from the laws of the core datum through $z^{12}\le P$
(\lref{GlobalTheory.CovTop.CovSz.sizes_of_cover_alt}, which derives them from the host numerics $\mathcal H(P,A,6)$
and the size data); (2) the bound for the flat degree
(\lref{GlobalTheory.T.abs_flatDegVal_logOm_le_nine_mul_of_size}); (3) $\log\kappa\le2\ell$ and $TP\le\sqrt b$
(\lref{GlobalTheory.T.log_kap_le_and_lenT_mul_le_sqrt}); (4) the two bounds for $\lifttwozed(u_1)$
(\lref{GlobalTheory.T.weight_parameter_size_bounds'}); (5) the bounds for $|u_1|_1$ and $|u_2|_1$
(\lref{GlobalTheory.T.weight_param_coeff_sum_bound'}); and (6) the bounds for $|\mathrm{wt}_t|_1$
(\lref{GlobalTheory.T.weight_param_coeff_sum_le_at_index'}). Cover conditions are closed under conjunction by
\cref{pf:lift2:combine}.
\end{proof}

\subsubsection*{Degrees of images and proportionality}

\begin{definition}[Linear systems and image degrees]\label{pf:lift2:def:linsys}
Let $F$ be a field of characteristic $0$. A \emph{linear system} $S$ over $F$ consists of an $F$-scheme $Y$ with an
integer $D$, a line bundle $L$ on $Y$ and global sections $b_1,\dots,b_N$ of $L$. It is a \emph{variety} if $Y$ is
smooth of relative dimension $D\ge1$, projective and geometrically connected over $F$; \emph{base point free} if the
$b_j$ have no common zero; and \emph{birational} if at the generic point $\eta$ of $Y$ the ratios $b_j/b_{j'}$
generate the function field $F(Y)$. Its \emph{image ideal} $I(S)\subseteq F[X_1,\dots,X_N]$ is the ideal generated by
the homogeneous polynomials $g$ with $g(b_1,\dots,b_N)=0$ in $H^0(Y,L^{\deg g})$, and its \emph{degree} is the top
self-intersection number $\delta(S)=(L^D)$. For a homogeneous ideal $I$ with $\dim F[X]/I=D+1$ let $e(I)$ be its
Hilbert degree ($D!$ times the leading coefficient of its Hilbert polynomial). For an $F$-linear map
$\mathrm{ev}\colon V\to\Omega$ from a finite-dimensional space to a field extension of $F$, the \emph{image degree}
$\mathrm{dg}_D(\mathrm{ev})$ is the degree of an integral Chow form, in dimension $D$, of the ideal of algebraic
relations among the images of a basis of $V$, if such a form exists, and $0$ otherwise (so that
$\mathrm{dg}_Z(u)=\mathrm{dg}_{\lifttwodF}$ of the evaluation of $V_Z(u)$ at the generic point of $Z$). A linear
system with Hermitian data $S$ \emph{presents} $(V,\Lambda,\nu,\mathrm{ev})$ if there is a basis of $V$ which is a frame
for $(\Lambda,\nu)$ with the weights of $S$ and in which $I(S)$ is the ideal of relations among the images under
$\mathrm{ev}$.
\formallabel{GlobalTheory.Adelic.LinSys}
\formallabel{GlobalTheory.Adelic.Tuple}
\formallabel{GlobalTheory.CompHt.Presents}
\formal{GlobalTheory.CompHt.imageDeg}
\end{definition}

\begin{proposition}[Degree of a birational image]\label{pf:lift2:degimage}
Assume \textup{(Q)}. Let $S$ be a linear system over a number field $F$ which is a variety, base point free and
birational. Then $\delta(S)>0$, $e(I(S))=\delta(S)$ and $\dim F[X]/I(S)=D+1$.
\formallabel{GlobalTheory.ImageDegree.degree_image_eq_topSelf}
\formallabel{GlobalTheory.T.degree_image_birational_linearSystem_eq_selfInt}
\formallabel{GlobalTheory.hilbertDegree_pos}
\formallabel{GlobalTheory.ImageDegree.finrank_quotientPiece_le_lineBundleSections}
\end{proposition}

\begin{proof}
(1) The image of $Y$ under the base-point-free birational map defined by the $b_j$ has dimension $D$, so
$\dim F[X]/I(S)=D+1$ (\lref{GlobalTheory.ImageDegree.ringKrullDim_of_birational}: the homogeneous coordinate ring of
the image of a base-point-free birational linear system on a variety of dimension $D$ has Krull dimension $D+1$).
(2) $I=I(S)$ is homogeneous, and prime because $Y$ is integral; hence its Hilbert polynomial has degree $D$
(\lref{GlobalTheory.degree_idealHilbertPoly_eq_of_isPrime}: a homogeneous prime ideal with quotient of Krull
dimension $D+1$ has a Hilbert polynomial of degree $D$), and $\dim(F[X]/I)_m=e(I)m^D/D!+O(m^{D-1})$
\cite[Ch.~2]{Fulton1998}.
(3) By definition of $I$, the map $(F[X]/I)_m\to H^0(Y,L^m)$, $X_j\mapsto b_j$, is injective, so
$\dim(F[X]/I)_m\le h^0(Y,L^m)$.
(4) $L$ is generated by global sections, hence nef \cite[1.4]{Lazarsfeld2004}. By the numerics of nef line bundles in
\textup{(Q)}, $(L^D)\ge0$ and $D!\,h^0(Y,L^m)/m^D\to(L^D)$.
(5) Conversely, the growth of $h^0(Y,L^m)$ is at most that of $\dim(F[X]/I)_m$
(\lref{GlobalTheory.ImageDegree.rankGrowth_lineBundleSections_le_image_forms}: for a base-point-free birational
linear system on a variety, $\limsup D!\,h^0(L^m)/m^D\le\limsup D!\dim(F[X]/I)_m/m^D$).
By (2)--(5), $e(I)=(L^D)=\delta(S)$. Finally, at a generic point some $b_j$ does not vanish (base point freeness), so the
linear form $X_j$ is not in the prime ideal $I$; then $X_j$ is a non-zero-divisor of degree $1$ on $F[X]/I$, the Hilbert
function is non-decreasing and eventually positive, and $e(I)>0$.
\end{proof}

\begin{corollary}[Image degree of a presented system]\label{pf:lift2:imagedeg}
Assume \textup{(Q)}. Let $S$ be a linear system with Hermitian data over a number field $F$ which is a variety, base point
free and birational, and which presents $(V,\Lambda,\nu,\mathrm{ev})$. Then $\mathrm{dg}_D(\mathrm{ev})=\delta(S)$.
\formallabel{GlobalTheory.HSCore.ImageDegree.imageDeg_eq_del}
\end{corollary}

\begin{proof}
By \cref{pf:lift2:degimage}, $I(S)$ is a homogeneous prime ideal with $\dim F[X]/I(S)=D+1$ and $e(I(S))=\delta(S)$. The
existence of integral Chow forms in \textup{(Q)} gives an integral Chow form of $I(S)$ of degree $e(I(S))$. Since $S$
presents $(V,\Lambda,\nu,\mathrm{ev})$, $I(S)$ is the ideal of relations among the images of a basis of $V$ under
$\mathrm{ev}$, so $\mathrm{dg}_D(\mathrm{ev})$ is the degree of this Chow form, namely $e(I(S))=\delta(S)$.
\end{proof}

\begin{definition}[Volume constants and proportionality]\label{pf:lift2:def:proport}
For a component $Z$ of a unitary moduli space with uniformisations $U$ and a weight $G$ of automorphic forms let
$\mathrm{gr}_Z(G)=\limsup_{\tau\to\infty}\lifttwodF!\dim V_Z(\tau G)/\tau^{\lifttwodF}$; let
$v_0(Z)=\mathrm{gr}_Z(\underline{\mathrm{std}})/\lifttwovdm(\mathrm{std})$ with $\mathrm{std}_{\tau,l}=N-1-l$ (the
\emph{volume constant} of $Z$), and $\mathrm{top}_Z(u)=v_0(Z)\lifttwovdm(u)$. The uniformisations satisfy
\emph{proportionality} if $\mathrm{gr}_Z(\underline u)=\mathrm{top}_Z(u)$ for every component $Z$ and every weight $u$
with all $u_\tau$ strictly decreasing.
\formallabel{GlobalTheory.DegCmp.Proport}
\formal{GlobalTheory.CompHt.growthOn}
\formal{GlobalTheory.CompHt.vZero}
\formal{GlobalTheory.CompHt.topDeg}
\end{definition}

\begin{proposition}[Degrees of components are bounded by volumes]\label{pf:lift2:compdeg}
Assume that $U$ satisfies proportionality (\cref{pf:lift2:def:proport}). Let $Z$ be a component and $u$ a weight with all $u_\tau$ strictly
decreasing. Then $\mathrm{dg}_Z(u)\le\mathrm{top}_Z(u)$.
\formallabel{GlobalTheory.DegTop.compDeg_le}
\formallabel{GlobalTheory.DegTop.growth_hilb_le}
\formallabel{GlobalTheory.DegTop.topDeg_nonneg}
\formallabel{GlobalTheory.DegTop.finiteDimensional_cuspOn}
\end{proposition}

\begin{proof}
As $\lifttwovdm(u)>0$ and $v_0(Z)\ge0$, $\mathrm{top}_Z(u)\ge0$, so we may assume $\delta=\mathrm{dg}_Z(u)>0$. By
proportionality the growth of $\dim V_Z(\tau\underline u)$ is finite, which forces $V_Z(u)$ to be finite-dimensional
(\lref{GlobalTheory.DegTop.finiteDimensional_of_growthOn}: a space of cusp forms of an even weight whose multiples
have finite growth is finite-dimensional). Let $\mathrm{ev}$ be the evaluation of $V_Z(u)$ at the generic point of $Z$ and
$H(\tau)$ the dimension of the span of the products of $\tau$ elements of $\mathrm{ev}(V_Z(u))$. Since an integral Chow
form of degree $\delta$ exists, $\limsup_\tau\lifttwodF!H(\tau)/\tau^{\lifttwodF}=\delta$
(\lref{GlobalTheory.T.cone_degree_growth_chowForm_integrality}: the Hilbert function of the cone over an image with
an integral Chow form of degree $\delta$ grows like $\delta\tau^D/D!$). A product of $\tau$ forms of weight $u$ supported
on $Z$ is a form of weight $\tau u$ supported on $Z$, so $H(\tau)\le\dim V_Z(\tau\underline u)$, and by proportionality
$\delta\le\mathrm{gr}_Z(\underline u)=\mathrm{top}_Z(u)$.
\end{proof}

\begin{lemma}\label{pf:lift2:delprime}
For a component $Z$ and a weight $u$ with $\lifttwovdm(u)\ne0$, $\mathrm{top}_Z(u')=(1-\rho_B(u))\,\mathrm{top}_Z(u)$.
\formallabel{GlobalTheory.CompHt.delPrime_eq}
\formallabel{GlobalTheory.CompHt.delPrime}
\end{lemma}

\begin{proof}
By definition $1-\rho_B(u)=\lifttwovdm(u')/\lifttwovdm(u)$; multiply by $\mathrm{top}_Z(u)=v_0(Z)\lifttwovdm(u)$.
\end{proof}

\begin{proposition}[Proportionality on good deep covers]\label{pf:lift2:proport}
Assume \textup{(Hg)}, \textup{(Sg)}, \textup{(KS)} and \textup{(Lie)}. Let the core datum satisfy its laws, let $j$
be in range with $p_j,q_j\ge1$, $[E:\Q]/2>2$, and let $Y$ be an $(e,o)$-good deep cover. Then the uniformisations of $Y$
satisfy proportionality, and $v_0(Z)>0$ for every component $Z$.
\formallabel{GlobalTheory.T.proportionality_principle_compactification}
\formallabel{GlobalTheory.CompProport.hirzebruch_proportionality_toroidal}
\formallabel{GlobalTheory.CompProport.proportionality_of_local_charts}
\end{proposition}

\begin{proof}
By \textup{(KS)} and \textup{(Sg)}, the moduli space of $Y$ has a model, smooth of relative dimension $\dim X^o_Y$, with
simple normal crossings boundary, a logarithmic Kodaira--Spencer isomorphism on the interior and charts at all points of
the boundary of the generic fibre; from it the development builds a toroidal flag datum with toric charts (\lref{GlobalTheory.CompProport.exists_torFlag_charts}: a toroidal flag datum whose
boundary points have toric charts exists for such a model). Fix a good embedding $\sigma$ of $F_Y$. The volume
computation (\lref{GlobalTheory.CompProport.invariant_volume_eq_covolume}: under \textup{(Hg)} and \textup{(Lie)},
for such a datum, proportionality holds, and there is a bijection $k\mapsto Z_k$ from the components of the complex
uniformisation at $\sigma$ to the components of the generic fibre with
$v_0(Z_k)=\mathrm{covol}(X_k)\,\lifttwodF!/\lifttwovdm(\mathrm{std})$, where $\mathrm{covol}(X_k)$ is the normalised volume
of a fundamental domain of the arithmetic group of the $k$-th complex component) gives proportionality. Since
$\mathrm{covol}(X_k)>0$ (\lref{GlobalTheory.TrueCx.covol_pos}: fundamental domains have positive volume) and
$\lifttwovdm(\mathrm{std})>0$ ($\mathrm{std}$ is strictly decreasing), $v_0(Z)>0$ for every component.
\end{proof}

\begin{proposition}[The Hilbert--Samuel datum presents its component]\label{pf:lift2:presents}
Let $X^o$ be a unitary moduli space over $F$ with a model, a flag cover of the model, a component $Z$,
uniformisations, a weight parameter $P$ and the auxiliary choices of the Hilbert--Samuel construction of the companion
section (a modification of the flag cover adapted to $Z$ and $P$, a canonical identification and readiness data), cusp
frames $f_1$, $f_2$ on $Z$ for $u_1(P)$ and $u_2(P)$, a section $\theta$ and constants. Then the associated
Hilbert--Samuel datum $X_Z$ presents $Z$ for the weight $u_1(P)$: its dimension is $\lifttwodF$, the frame $f_1$ is a frame
of the cusp forms of weight $u_1(P)$ on $Z$, and the image ideal of its linear system is the ideal of algebraic
relations among the values of the forms $f_1$ at the generic point of $Z$.
\formallabel{GlobalTheory.CompHS.compPresents_data}
\end{proposition}

\begin{proof}
The dimension of $X_Z$ is $\lifttwodF$ by construction, and $f_1$ is a frame by definition of a cusp frame. The moduli space
has a generic frame over the function field $\Omega_Z$ of $Z$ (\lref{GlobalTheory.Reg.genFrEx_and_genFromTest}:
generic frames exist and are obtained from test presentations); it lifts to a point $z$ of the modification over
$\Omega_Z$ (\lref{GlobalTheory.CompHS.exists_lift_genFr}), and this point is a generic point of the variety of $X_Z$
(\lref{GlobalTheory.CompHS.isGenericPoint_of_lift}), lying over the canonical point $\mathrm{Spec}\,\Omega_Z\to
\mathrm{Spec}\,F$. In a chart at $z$, after a change of frame of the linear system, the values at $z$ of the sections of
the linear system are $c$ times the values at the generic point of the forms of $f_1$, for one constant $c\ne0$
(\lref{GlobalTheory.CompHS.exists_chart_val}). If at a point over the generic point of an integral scheme the
sections of a linear system take values $c$ times given elements of a field, its image ideal is the ideal of algebraic
relations among these elements (\lref{GlobalTheory.CompHt.imageIdeal_eq_coneIdeal_of_generic}); this gives the
claim.
\end{proof}

\subsubsection*{Toroidal flag data, nefness and birationality}

\begin{definition}[Toroidal flag data]\label{pf:lift2:def:torflag}
Let $X^o$ be a unitary moduli space over a coefficient field $F$ (a CM host field $E$, a datum of signature $(p,q)$, a
level and a coefficient datum), with $N=p+q$ and $d=\lifttwodF$. A \emph{toroidal flag datum} for $X^o$ consists of a
model $\overline X$ of $X^o$ which is smooth of relative dimension $\dim X^o$ over $F$, whose boundary is a divisor with
simple normal crossings, which has a logarithmic Kodaira--Spencer isomorphism on its interior, and a flag model over
$\overline X$ (the relative flag variety of the Lie algebra of the universal semi-abelian scheme, with its projection
$\pi$). For a component $Z$ let $X_Z$ be the closure of $Z$ in the flag model, and for a weight $\lambda$ let $L_0(\lambda)$ be
the automorphic line bundle of $\lambda$ on $X_Z$. We say that the function fields are algebraically closed if the
coefficient field is algebraically closed in the function field of every component (one of the conditions of
goodness of deep covers).
\formal{GlobalTheory.CompHS.TorFlag}
\formallabel{GlobalTheory.CompHS.IsVariety₀}
\end{definition}

\begin{proposition}[Existence of toroidal flag data]\label{pf:lift2:torflag}
Assume \textup{(Hg)}. If $d\ge1$, then $X^o$ has a toroidal flag datum.
\formallabel{GlobalTheory.T.torFlag_nonempty_of_one_le_dimFlag}
\formallabel{GlobalTheory.CompHS.TorFlagEx.nonempty_torFlag_alt}
\end{proposition}

\begin{proof}
As $d\ge1$, $N\ge1$. The statement on toroidal compactifications in \textup{(Hg)} gives a model $\overline X$, smooth of
relative dimension $\dim X^o$, with simple normal crossings boundary and a logarithmic Kodaira--Spencer isomorphism on
the interior. Such a model carries theta data (\lref{GlobalTheory.ThetaNormal.hasTheta_of_smooth}: a smooth model of
this kind admits the theta normalisation of the development), hence the push-forward of the relevant line bundle is
invertible (\lref{GlobalTheory.CompHS.TorFlagEx.pushInvertible_of_hasTheta}), and a model with this property carries a
flag model (\lref{GlobalTheory.CompHS.TorFlagEx.torFlag_of_model}).
\end{proof}

\begin{proposition}[Closures of components in the flag model]\label{pf:lift2:flagvar}
For every toroidal flag datum and every component $Z$, if the function fields are algebraically closed, then $X_Z$ is
smooth of relative dimension $d$, projective and geometrically connected over $F$.
\formallabel{GlobalTheory.T.flagSpace_component_smooth_projective_connected}
\formallabel{GlobalTheory.VarFlag.isVariety₀_flag}
\end{proposition}

\begin{proof}
The generic fibre of the model is smooth, and its components are smooth projective geometrically connected varieties
(\lref{GlobalTheory.T.exists_varieties_components_smooth_genericFibre}). Over the variety $B_Z$ attached to $Z$ the
development constructs a flag bundle $C\to B_Z$ of the right dimension (\lref{GlobalTheory.FlagBundle.exists_flagCpt}).
The closure of $Z$ in the model is closed, so the inclusion of $X_Z$ into the flag model is a closed immersion and
$X_Z$ is proper and noetherian. Every point of $X_Z$ with values
in a field carries a unique point of $C$ with the same flag
(\lref{GlobalTheory.T.flagModel_point_carries_unique_flagBundle_point}), compatibly with field extensions and with
specialisation over valuation rings; by the valuative criterion this defines a proper morphism $h\colon X_Z\to C$ over
the model (\lref{GlobalTheory.T.exists_hom_of_fieldPoint_relation_proper}), which is locally quasi-finite
(\lref{GlobalTheory.T.flagModel_to_flagBundle_locallyQuasiFinite}) and an isomorphism over the open part $V$ of $C$ lying
over the interior (\lref{GlobalTheory.T.isIso_flagModel_flagBundle_over_moduli}). The preimage of $V$ is dense in
$X_Z$ and irreducible, so $X_Z$ is integral; $C$ is normal; and a proper quasi-finite birational morphism onto a normal
integral scheme is an isomorphism (\lref{GlobalTheory.isIso_of_isProper_of_quasiFinite_of_birational}, Zariski's main
theorem). Hence $X_Z\cong C$, which is smooth of relative dimension $\dim B_Z+\dim(\text{flag fibre})=d$, projective and
geometrically connected.
\end{proof}

\begin{proposition}[Nefness of automorphic line bundles]\label{pf:lift2:nef}
Assume \textup{(Hg)}, the numerics of nef line bundles in \textup{(Q)}, and \textup{(Lie)}. For every toroidal flag datum,
every component $Z$ such that the function fields are algebraically closed, and every weight $\lambda$ with all
$\lambda_\tau$ non-increasing, the line bundle $L_0(\lambda)$ on $X_Z$ is nef.
\formallabel{GlobalTheory.T.base_weight_line_bundle_is_nef}
\formallabel{GlobalTheory.NefWt.weight_line_bundle_isNef}
\formallabel{GlobalTheory.NefWt.decomp}
\formallabel{GlobalTheory.IsNef.of_cls_eq_sum}
\end{proposition}

\begin{proof}
If $N=0$ then $\lambda=0$ and $L_0(0)$ is trivial. Let $N\ge1$. For $\tau\in\Psi$ and $k<N-1$ let $e^{\le k}_\tau$ be the
weight equal to $1$ at $(\tau,l)$ for $l\le k$ and $0$ elsewhere, $-e^{>k}_\tau$ the weight equal to $-1$ at $(\tau,l)$ for
$l>k$ and $0$ elsewhere, and $\mathbf1_\tau$ the weight equal to $1$ on all entries of $\tau$. With
$g_{\tau,k}=\mathrm{gap}(\lambda_\tau)_k\ge0$ and $c_\tau=\lambda_{\tau,N-1}+\sum_{p\le k<N-1}g_{\tau,k}$ we have
\[
\lambda=\sum_{\tau}\Bigl(\sum_{k<\min(p,N-1)}g_{\tau,k}\,e^{\le k}_\tau+\sum_{p\le k<N-1}g_{\tau,k}\,(-e^{>k}_\tau)
+c_\tau\mathbf1_\tau\Bigr),
\]
since $e^{\le k}_\tau=\mathbf1_\tau-e^{>k}_\tau$ and $\lambda_{\tau,l}=\lambda_{\tau,N-1}+\sum_{l\le k<N-1}g_{\tau,k}$. By
\cref{pf:lift2:flagvar} $X_Z$ is projective, so it carries an ample line bundle, and the map $\lambda\mapsto$ class of
$L_0(\lambda)$ is additive. The line bundles $L_0(e^{\le k}_\tau)$, $k<p$, and $L_0(-e^{>k}_\tau)$, $k\ge p$, are nef
(\lref{GlobalTheory.T.hodgeBundle_headWeight_nef}, \lref{GlobalTheory.T.hodgeBundle_tailWeight_nef}: under
\textup{(Lie)}, the toroidal statements of \textup{(Hg)} and the numerics of nef line bundles, the line bundles of these
two kinds of weights on $X_Z$ are nef), and both $L_0(\mathbf1_\tau)$ and $L_0(-\mathbf1_\tau)$ are nef
(\lref{GlobalTheory.T.hodgeBundle_unitWeight_nef}). Writing $c_\tau\mathbf1_\tau$ as $|c_\tau|$ times $\pm\mathbf1_\tau$, the
class of $L_0(\lambda)$ is a combination with non-negative integer coefficients of classes of nef line bundles, hence nef.
\end{proof}

\begin{definition}[Birationality]\label{pf:lift2:def:biron}
The forms of a weight $G$ are \emph{birational} on a component $Z$ if their evaluation at the generic point of $Z$ is
nonzero and the ratios of their values generate the function field $\Omega_Z$ of $Z$ over $F$.
\formal{GlobalTheory.DegVal.BirOn}
\formal{GlobalTheory.DegVal.GenRatios}
\end{definition}

\begin{proposition}[Birationality for weights with large gaps]\label{pf:lift2:biron}
Assume \textup{(X)}, \textup{(GM)}, \textup{(Q)}, \textup{(Sg)} and \textup{(Hg)} (of which the toroidal and volume
statements, the de~Rham statement, the ray class field and semistable reduction statements, Nadel vanishing and
principalisation are used). Let the core datum satisfy its laws, let $p_j,q_j\ge1$, let $Y$ be an $(e,o)$-good deep cover
at level $j$, and let $u$ be a weight all of whose consecutive gaps are at least $(\lifttwodF+3)^2+2N$. Then for every
component $Z$ of $Y$ the forms of weight $\underline u$ are birational on $Z$.
\formallabel{GlobalTheory.T.birOn_wtOf_of_consecutive_gap_bound}
\formallabel{GlobalTheory.DegVal.birOn_of_weight_gaps_large}
\end{proposition}

\begin{proof}
The development proves, under these hypotheses:
(1) some cusp form of weight $\underline u$ supported on $Z$ has nonzero value at the generic point of $Z$
(\lref{GlobalTheory.T.exists_centre_cusp_generic_value_ne_zero}, which states exactly this);
(2) there are an embedding $\sigma\colon F_Y\to\C$ and finitely generated $F_Y$-subalgebras $B\subseteq A$ of $\Omega_Z$ such
that $A$ generates $\Omega_Z$ as a field, $B$ is contained in the field generated by the ratios of values of forms of weight
$\underline u$, and for every $b\in B\smallsetminus\{0\}$ there is an $F_Y$-algebra homomorphism $x\colon A\to\C$ over
$\sigma$ with $x(b)\ne0$ which is the only such homomorphism with its restriction to $B$
(\lref{GlobalTheory.T.cusp_values_separate_generic_points}, which states the existence of such $\sigma$, $A$ and
$B$);
(3) an algebraic lemma (\lref{GlobalTheory.DegVal.exists_mul_eq_of_generic_unique_complexPoint}): in the situation of
(2), every $a\in A$ satisfies $aq=p$ for some $p,q\in B$ with $q\ne0$.
By (3), $A$ is contained in the fraction field of $B$, hence in the field generated by the ratios of values of forms; as
$A$ generates $\Omega_Z$, this field is $\Omega_Z$. Together with (1) this is the claim.
\end{proof}

\begin{proposition}[Birationality for the weights of the parameter]\label{pf:lift2:biron12}
Assume \textup{(X)}, \textup{(GM)}, \textup{(Q)}, \textup{(Sg)}, \textup{(Hg)} and $\lifttwoC{Size}$. Then the cover condition
of the following statement holds: for every component $Z$, the forms of weights $\underline{u_1}$ and $\underline{u_2}$
(of $P(p,i)$) are birational on $Z$.
\formallabel{GlobalTheory.T.birOn_weightPair_both_components_alt}
\formallabel{GlobalTheory.HSApp.birOn_u₁_u₂_alt_alt}
\formallabel{GlobalTheory.HSApp.gaps_u₁_u₂_of_wtBig}
\end{proposition}

\begin{proof}
Combine \cref{pf:lift2:sizecov}, \cref{pf:lift2:wtpt} and \cref{pf:lift2:wtbig} by \cref{pf:lift2:combine}(c), with the
constant $0$. At a good deep cover of an admissible point with $j<m$, the core datum satisfies its laws, $j$ is in range,
$p_j,q_j\ge1$, and $P(p,i)$ is large: all gaps of $u_1$ are at least $G(d,N)=(d+3)(\binom{d+1}2+4N+d+3)\ge(d+3)^2+2N$, and
the gaps of $u_2=u_1+\varrho$ are at least those of $u_1$, as the gaps of $\varrho$ are non-negative. Apply
\cref{pf:lift2:biron} to $u_1$ and to $u_2$.
\end{proof}

\begin{proposition}[Generation at the generic point of the flag space]\label{pf:lift2:ratgen}
Let $X^o$ have a model with a flag cover, let $Z$ be a component, $U$ uniformisations and $\lambda$ a weight. Suppose that
the closure of $Z$ in the flag cover admits a modification which is an isomorphism over the open part, and that the
forms of weight $\underline\lambda$ are birational on $Z$. Then there is a generic point $\eta$ of the closure, lying in
the open part, at which the ratios of the sections of $L_0(\lambda)$ given by the cusp forms of weight $\underline\lambda$
generate the residue field.
\formallabel{GlobalTheory.T.cuspForm_ratios_generate_at_flagSpace_genericPoint}
\formallabel{GlobalTheory.HSApp.ratGen_of_birOn}
\formallabel{GlobalTheory.HSApp.RatGen}
\end{proposition}

\begin{proof}
The generic frame of $X^o$ over $\Omega_Z$ lifts to a point $z$ of the modification
(\lref{GlobalTheory.CompHS.exists_lift_genFr}); its image $\eta$ in the closure lies in the open part and is a generic
point, because the image of the modification is dense. At $\eta$ choose a chart of $L_0(\lambda)$; there are units
$t_\iota$ of $\Omega_Z$ (the diagonal of a Borel change of frame relating the chart frame to the generic frame) such that
for every cusp form $g$ of weight $\underline\lambda$ on $Z$ the value at $\eta$ of the section of $g$, read in the residue
field identified with $\Omega_Z$, is $T\cdot g(\text{generic frame})$ with $T=\prod_\iota t_\iota^{-\mathrm{ex}_\iota}\ne0$
independent of $g$ (\lref{GlobalTheory.FlagMod.FlagCover.exists_chartFr_frBor}: every framed point in a chart is related to a chart
frame by a Borel change of frame; \lref{GlobalTheory.CompHS.evR_eq_val_secOf}: the value of the section of a form in a
chart frame is the product of the diagonal entries, raised to minus the exponents, times the value of the form).
Hence the ratio of the sections of $g$ and $g'$ at $\eta$ is the ratio of the generic values of $g$ and $g'$. By
birationality some form $g_0$ has nonzero generic value, and the ratios of generic values generate $\Omega_Z$; therefore the
field generated by the ratios of sections at $\eta$ maps onto $\Omega_Z$, which is the claim.
\end{proof}

\begin{proposition}[Hilbert--Samuel data exist]\label{pf:lift2:hsready}
Assume \textup{(Hg)}, \textup{(GM)}, \textup{(Q)}, \textup{(Sg)}, the toroidal charts for fans and \textup{(Lie)} (Nadel
vanishing and principalisation being parts of \textup{(Hg)}). Let $Y$ be an $(e,o)$-good deep cover with $[E:\Q]/2>2$,
$Z$ a component and $P$ a large weight parameter such that the weights $\mathrm{wt}_t(P)$, $t\le T$, are good on $Y$. Then
there are a toroidal flag datum, a modification $R$ of the closure of $Z$ adapted to $P$ and a canonical identification
$\omega$ of the canonical bundle of the modification, such that the readiness conditions hold: the structure map of $R$ is
geometrically connected; $d\ge1$ and $d+1\le T$; the cusp forms of weights $u_1(P)$ and $u_2(P)$ on $Z$ extend to sections
of $L_0(u_1)$, $L_0(u_2)$ on the closure; these line bundles are generated by such sections over the open part; the forms
of the weights $\mathrm{wt}_t(P)$, $t\le T$, are good on $Z$; for $t\le T$ every global section of the $t$-th adjoint line
bundle of $R$ comes, via $\omega$, from a cusp form of weight $\mathrm{wt}_t(P)$; and for $t=0$ every such cusp form gives a
global section.
\formallabel{GlobalTheory.T.exists_weight_hodge_bundle_generated}
\formallabel{GlobalTheory.HSApp.exists_weight_sections_nonvanishing}
\end{proposition}

\begin{proof}
Since $P$ is large, $d\ge1$ and $d+1\le T$; \cref{pf:lift2:torflag} gives a toroidal flag datum, and \cref{pf:lift2:flagvar}
shows that the closure $X_Z$ is a smooth projective geometrically connected variety (the function fields of a good deep
cover are algebraically closed). Cusp forms of every weight on $Z$ extend to sections on the closure, vanishing on the
boundary (\lref{GlobalTheory.T.cuspForms_extend_vanish_compactifiedFlagSpace},
\lref{GlobalTheory.HSApp.vanBd_le_space}). Largeness implies that $u_1$ and $u_2$ have all gaps $\ge1$ and that
$u_i-\kappa_{\mathrm{Fl}}$ have all gaps $\ge\binom{d+1}2+1$ (\lref{GlobalTheory.HSApp.gaps_of_wtBig}); under this gap
condition the line bundles $L_0(u_i)$ are generated by cusp sections over the open part
(\lref{GlobalTheory.T.hodge_bundle_base_point_free_on_open}: if $\lambda-\kappa_{\mathrm{Fl}}$ has all gaps
$\ge\binom{d+1}2+1$ and the cusp forms of weight $\lambda$ extend, then $L_0(\lambda)$ is generated by cusp sections over the
open part). Principalisation then gives a modification $R$ adapted to $u_1,u_2$
(\lref{GlobalTheory.CompHS.nonempty_choice}); its structure map is geometrically connected since $X_Z$ is and the
open part is dense. The Kodaira--Spencer isomorphism of the flag tangent bundle
(\lref{GlobalTheory.T.flag_tangent_kodaira_spencer_iso}) and the de~Rham statement of \textup{(Hg)} give a canonical
identification $\omega$ (\lref{GlobalTheory.CompHS.canId_nonempty}), and the comparison of adjoint sections with cusp
forms is \lref{GlobalTheory.T.hodge_bundle_degree_squeeze}. Finally goodness of the weights $\mathrm{wt}_t(P)$ on $Y$
implies goodness on $Z$ (\lref{GlobalTheory.HSApp.goodOn_of_goodAt}).
\end{proof}

\subsubsection*{Cusp forms on good deep covers}

\begin{proposition}[Regularity of weights on good deep covers]\label{pf:lift2:regdom}
Assume \textup{(Hg)}. Let the core datum satisfy its laws, let $j$ be in range with $p_j,q_j\ge1$ and $[E:\Q]/2>2$, and let
$Y$ be an $(e,o)$-good deep cover. Let $\lambda$ be a weight with all $\lambda_\tau$ non-increasing and
$\mathrm{gap}(\lambda_\tau)_{p_j-1}\ge1$ for all $\tau$. Then $\mathcal G_Y(\underline\lambda)$ holds, and
$\mathcal G_Z(\lambda)$ holds for every component $Z$.
\formallabel{GlobalTheory.T.covering_domain_lattice_regularity}
\formallabel{GlobalTheory.CovTop.GoodDom.covering_regular_at_domain}
\end{proposition}

\begin{proof}
On a good deep cover, every space $V_Y(G)$ is finite-dimensional (\lref{GlobalTheory.Deep.finiteDimensional_V_of_good})
and the norms $\nu^Y_G$ have a Gram frame (\lref{GlobalTheory.HSApp.hasGram_nrm_of_good}: the Petersson-type norms of
a good deep cover come from inner products, to which \cref{pf:lift2:hasgram} applies). For weights $\lambda$ as in the
statement the lattice $\Lambda_Y(\underline\lambda)$ is finitely generated
(\lref{GlobalTheory.CovTop.GoodDom.weight_lattice_fg_on_domain}: under \textup{(Hg)}, for weights with
non-increasing entries and a positive gap at $p_j-1$ on a good deep cover, the lattice of integral cusp forms is finitely
generated) and spans $V_Y(\underline\lambda)$ (\lref{GlobalTheory.CovTop.GoodDom.weight_lattice_span_on_domain}: under
the same hypotheses it spans the space of cusp forms). This is $\mathcal G_Y(\underline\lambda)$, and on a good deep
cover with $[E:\Q]/2>2$ goodness on $Y$ implies goodness on every component
(\lref{GlobalTheory.HSApp.goodOn_of_goodAt}).
\end{proof}

For exponents $a$ and an even $n$ which are dominant (\cref{pf:lift2:def:theta}), an index $k_c$ and $s\in\Z$, let
$\lambda(a,n)$ be the weight of the datum attached to $(a,n)$ as in \cref{pf:lift2:def:wt}, $u(a,n,k_c)$ the weight
$\lambda(a,n)$ minus its value at the index $k_c-1$ (a weight constant in the second variable), and $G_s(a,n,s)$ the
weight of automorphic forms attached to $(a,n)$ shifted by $s$ times the Hodge weight; for the point $p$ these are
$\lambda$, $u_W$ and $G_s(s)$. \emph{Flat sections exist} on $Y$ if the line bundles of the weights $-\mathbf1_\tau$ have
constant sections, as in $\lifttwoP{FlatDegAt}$.

\begin{proposition}[Ranks under twisting]\label{pf:lift2:twistrank}
Assume \textup{(Hg)}. Let the core datum, $j$ and $Y$ be as in \cref{pf:lift2:regdom}, and suppose that flat sections exist
on $Y$. Then:
\begin{enumerate}[label=(\alph*)]
\item for every weight $\lambda$, every $w\colon\Psi\to\Z$ and every component $Z$,
$\mathrm{rk}_Z(\lambda+w)=\mathrm{rk}_Z(\lambda)$ ($w$ regarded as a weight constant in the second variable);
\item for dominant $(a,n)$ with $n$ even, every $k_c$ and every $s\in\Z$ such that $G_s(a,n,s)$ is good on $Y$,
\[
\lifttwork_Y(G_s(a,n,s))=\sum_{Z\in\lifttwoComp(Y)}\mathrm{rk}_Z\bigl(u(a,n,k_c)+s\,h\bigr).
\]
\end{enumerate}
\formallabel{GlobalTheory.T.twisted_flat_bundle_ranks}
\formallabel{GlobalTheory.TwistRank.flat_twist_rank_eq}
\formallabel{GlobalTheory.TwistRank.rankOn_congr}
\end{proposition}

\begin{proof}
(a) A good deep cover carries a key section (\lref{GlobalTheory.KeyCover.nonempty_keySec_of_family}), and with it and
the flat sections one obtains, for every component $Z$ and every $w$, a \emph{twisting section} for the parallel weight $w$
on $Z$ (\lref{GlobalTheory.TwistRank.exists_twSec_par}): a nowhere vanishing section by which multiplication maps the
forms of every even weight $G$ supported on $Z$ isomorphically onto those of weight $G+\underline w$, preserving the
lattices (\lref{GlobalTheory.TwistRank.rankOn_twist}). Hence the ranks agree.
(b) As $G_s(a,n,s)$ is good, the rank of $\Lambda_Y(G_s(a,n,s))$ is the sum over the finitely many components of the
ranks of the component lattices (\lref{GlobalTheory.T.supLattice_rank_eq_sum_component_ranks}). Fix $Z$ and write
$n=2m'$. Two weights of automorphic forms with the same exponents give the same component ranks, so
$\mathrm{rk}_Z(G_s(a,n,s))=\mathrm{rk}_Z\bigl(\mathrm{wt}_{m'}(L)\bigr)$ with $L=\lambda(a,n)+s\,h$. Twisting by the $m'$-th power
of a twisting section for the weight $\mathrm{wt}_1(0)$ (exponents $-1$ at all entries of the datum, parity $2$), which exists
on good deep covers with flat sections (\lref{GlobalTheory.T.exists_twisted_section_with_deep_root}), gives
$\mathrm{rk}_Z(\mathrm{wt}_{m'}(L))=\mathrm{rk}_Z(\underline L)$. Finally $L=u(a,n,k_c)+s\,h+w$ with $w$ the value of $\lambda(a,n)$ at
the index $k_c-1$, a weight constant in the second variable, and (a) gives the claim.
\end{proof}

\begin{proposition}[Short generators]\label{pf:lift2:shortgen}
Assume \textup{(Hg)}. There is a constant $C\ge0$ with the following property. Let the core datum satisfy its laws, $j$ be in
range with $p_j,q_j\ge1$ and $[E:\Q]/2>2$, and let $Y$ be an $(e,o)$-good deep cover. Let $L_b\ge1$ with $\log M'_Y\le30L_b$,
$\mathrm{fs}(E)\le3L_b$, $\log(fr)\le L_b$ and $o\le391X^2$, where $X=frL_b$. Suppose flat sections exist on $Y$ with a
constant $c_1$, $|c_1|\le9rL_b$. Let $\lambda$ be a weight with all $\lambda_\tau$ non-increasing,
$\mathrm{gap}(\lambda_\tau)_{p_j-1}\ge8fr+10$ for all $\tau$, and $\log|\lambda|_1\le10L_b$. Then for every component $Z$ there
are $n\in\N$, forms $w_1,\dots,w_n\in V_Z(\lambda)$ and $\Lambda\ge1$ such that $\Lambda_Z(\lambda)$ is the abelian group
generated by the $w_i$, $\nu^Z_{\lambda,\sigma}(w_i)\le\Lambda$ for all $i$ and $\sigma$, and
\[
\log\Lambda\le C\bigl(X^3|\lambda|_1+X^5\bigr),\qquad \log n\le CX^3 .
\]
\formallabel{GlobalTheory.T.short_generator_norm_comparison}
\formallabel{GlobalTheory.ShortGenComp.short_generator_comparison}
\end{proposition}

\begin{proof}
The development provides two constants: $C_2\ge0$ such that, under the hypotheses, $\Lambda_Z(\lambda)$ contains a family
spanning $V_Z(\lambda)$ whose norms are at most $\exp(C_2(X^3|\lambda|_1+X^5))$
(\lref{GlobalTheory.T.short_span_composition_closed}); and $C_b\ge0$ such that $\log\mathrm{rk}_Z(\lambda)\le C_bX^3$
(the rank bound in \lref{UpperTheory.Lit.Tam.tamagawa_number_formula}, which also bounds Bergman kernels and sup
norms). By \cref{pf:lift2:regdom}, $\lambda$ is good on $Z$, so $\Lambda_Z(\lambda)$ is finitely generated, the norms are
non-negative, homogeneous for the absolute values of the embeddings, and subadditive (they come from a Gram frame). A
short spanning family gives short generators (\lref{GlobalTheory.T.exists_short_generators_of_short_spanning_family}:
if a finitely generated $\OO_K$-lattice $\Lambda$ in a $K$-space with such norms contains a $K$-spanning set with norms
$\le A$, then $\Lambda$ is generated as an abelian group by $n\le2[K:\Q]\lifttwork\Lambda$ elements of norm
$\le\max(1,[K:\Q]^4\lifttwork(\Lambda)\,\mathrm{rd}_K\,A)$). Take $C=2000+C_b+C_2$. Goodness of $Y$ gives
$\log[F_Y:\Q]\le o\le391X^2$ and $\log\mathrm{rd}_{F_Y}\le o$; with $X\ge1$,
\[
\log\Lambda\le4\cdot391X^2+C_bX^3+391X^2+C_2(X^3|\lambda|_1+X^5)\le C(X^3|\lambda|_1+X^5)
\]
(if the maximum is $1$, $\log\Lambda=0$), and $\log n\le\log2+\log[F_Y:\Q]+\log\mathrm{rk}_Z(\lambda)\le1+391X^2+C_bX^3\le
CX^3$ (if $n=0$, $\log n=0$).
\end{proof}

Recall from the companion section the framed points of a unitary moduli space $X^o$ over $F$ with values in an $F$-algebra
$S$, the rules (functions on framed points), the sections of integral exponent $e\in\Z^{\mathrm{Ix}}$ (rules transforming by
the character $\chi_e$ under changes of frame), the action $y\mapsto\gamma\cdot y$ of level automorphisms on framed points and
the translates $f_\gamma(y)=f(\gamma\cdot y)$ of rules, and for a finite group $S$ of level automorphisms the level determinant
$\Delta_S(f)=\det(f_{hg})_{h,g\in S}$. For a framed point $y$ over a local $F$-algebra, $\mathrm{ev}_y(f)$ denotes the
value of $f$ at $y$.

\begin{lemma}[Borel changes of frame]\label{pf:lift2:frbor}
Let $y$ be a framed point of $X^o$ over an $F$-algebra $S$ and $t\in(S^\times)^{\mathrm{Ix}}$. Then there is a framed point
$y'$ over $S$ obtained from $y$ by a Borel change of frame with diagonal $t$.
\formallabel{GlobalTheory.SecSeparation.exists_frBor}
\end{lemma}

\begin{proof}
Let $g_\tau$ be the diagonal matrix with entries $t_{(\tau,i)}$ and $a=t_*$ the entry of $t$ at the index of the scalar part.
Since $\det g_\tau$ is a unit, there is a unique framed point $y'$ obtained from $y$ by the change of frame $(g,a)$
(\lref{GlobalTheory.Frames.existsUnique_frChg}: for every family of matrices with unit determinants and every unit $a$,
the change of frame exists and is unique), and a change of frame by triangular matrices with diagonal $t$ is a Borel
change of frame with diagonal $t$ (\lref{GlobalTheory.Frames.frBor_calculus}).
\end{proof}

\begin{lemma}[Exponent of the level determinant]\label{pf:lift2:levdet}
Assume $p+q\ge1$ and that the $A$-part of the level is prime to its modulus. If $f$ is a section of integral exponent $e$,
then $\Delta_S(f)$ is a section of integral exponent $|S|\,e$.
\formallabel{GlobalTheory.LevDesc.levDet_mem_secZ}
\end{lemma}

\begin{proof}
By the Leibniz formula $\Delta_S(f)=\sum_\pi\mathrm{sgn}(\pi)\prod_{h\in S}f_{h\pi(h)}$. Each translate $f_\gamma$ is a section of
exponent $e$ (\lref{GlobalTheory.LevDesc.transl_mem_secZ}: under the two hypotheses, translation by a level
automorphism preserves the sections of a given exponent); a product of $|S|$ sections of exponent $e$ is a section of
exponent $|S|e$ (sections of exponents $e,e'$ multiply to sections of exponent $e+e'$, and $1$ has exponent $0$); and the
sections of a given exponent form a $\Z$-module.
\end{proof}

\begin{lemma}[Evaluation of translates]\label{pf:lift2:evtransl}
Under the hypotheses of \cref{pf:lift2:levdet}, for every level automorphism $\gamma$, every section $f$ of integral exponent and
every framed point $y$ over a local $F$-algebra, $\mathrm{ev}_y(f_\gamma)=\mathrm{ev}_{\gamma\cdot y}(f)$.
\formallabel{GlobalTheory.LevelKeeps.evR_transl}
\end{lemma}

\begin{proof}
A framed point over a local ring comes from a test presentation
(\lref{GlobalTheory.FlagTan.fromTest_of_isLocalRing}), and the value of a section at it is computed from the
presentation (\lref{GlobalTheory.FlagTan.evR_eq}). The translate $f_\gamma$ is again a section of the same exponent
(\lref{GlobalTheory.LevDesc.transl_mem_secZ}), and the action of $\gamma$ on framed points is natural with respect to
test presentations (\lref{GlobalTheory.FullLevelAction.actLaws_levFull}: the level action satisfies the naturality
laws); computing both sides through the presentation of $y$ and its image under $\gamma$ gives the same value.
\end{proof}
